\documentclass[11pt,a4paper]{article}
\usepackage[T1]{fontenc}
\usepackage{lmodern}
\usepackage[margin=27mm]{geometry}
\usepackage{amsmath,amssymb,amsthm,mathtools}
\usepackage{xcolor}
\usepackage{enumitem}
\usepackage{microtype}
\usepackage[unicode,colorlinks=true,linkcolor=blue!45!black,urlcolor=blue!55!black]{hyperref}
\usepackage{bookmark}
\usepackage{fancyhdr}
\usepackage{needspace}
\makeatletter
\renewcommand*{\l@subsection}{\@dottedtocline{2}{1.5em}{3.2em}}
\makeatother
\BeforeBeginEnvironment{definition}{\Needspace{12\baselineskip}}
\BeforeBeginEnvironment{lemma}{\Needspace{8\baselineskip}}
\BeforeBeginEnvironment{proposition}{\Needspace{9\baselineskip}}
\BeforeBeginEnvironment{theorem}{\Needspace{9\baselineskip}}
\setlength{\headheight}{14pt}
\pagestyle{fancy}
\fancyhf{}
\fancyhead[L]{\small Hypernets: mathematical commentary}
\fancyhead[R]{\small Lean 4 formalization}
\fancyfoot[C]{\thepage}
\setlength{\parskip}{4pt}
\setlength{\parindent}{0pt}
\setlist[enumerate]{leftmargin=*,itemsep=3pt}
\setlist[itemize]{leftmargin=*,itemsep=3pt}
\allowdisplaybreaks[2]
\theoremstyle{definition}
\newtheorem{definition}{Definition}[section]
\theoremstyle{plain}
\newtheorem{lemma}[definition]{Lemma}
\newtheorem{proposition}[definition]{Proposition}
\newtheorem{theorem}[definition]{Theorem}
\theoremstyle{remark}
\newtheorem{remark}[definition]{Remark}
\newcommand{\steps}[1]{\mathrel{\longrightarrow^{#1}}}
\newcommand{\fin}{\mathsf{Final}}
\newcommand{\normal}{\mathsf{Normal}}
\newcommand{\stuck}{\mathsf{Stuck}}
\newcommand{\rt}{\mathsf{Root}}
\newcommand{\mode}{\mathsf{mode}}
\newcommand{\ent}{\mathsf{Enter}}
\newcommand{\remote}{\mathsf{Remote}}
\newcommand{\settled}{\mathsf{Settled}}
\newcommand{\Lft}{\mathcal L_G}
\newcommand{\upto}[3]{\mathsf{UpTo}_{#1,#2}(#3)}
\newcommand{\quasi}[2]{\mathsf{Quasi}_G(#1,#2)}
\newcommand{\safe}[2]{\mathsf{InputSafe}_G(#1,#2)}
\newcommand{\robust}[3]{\mathsf{Robust}_G(#1,#2,#3)}
\newcommand{\repobase}{https://github.com/danghica/hypernets-lean/blob/fe1b8904b1e394450111eabffaf4a7eded297740}
\newcommand{\lean}[4]{\par\nopagebreak[4]\smallskip\nopagebreak[4]\noindent{\small\href{\repobase/Hypernets/#1.lean\#L#2-L#3}{\textsf{Lean: }\texttt{\detokenize{#4}}}}\par\nopagebreak[3]\smallskip}
\newcommand{\leanfile}[2]{\href{\repobase/#1}{\texttt{\detokenize{#2}}}}
% Definition destinations and the alphabetical appendix use internal PDF links.
\newcommand{\termdef}[1]{\par\Needspace{4\baselineskip}\phantomsection\label{term:#1}\ignorespaces}
\newcommand{\glossaryentry}[3]{\par\addvspace{7pt}\Needspace{4\baselineskip}\noindent\hyperref[term:#1]{\textbf{#2}}\hfill\hyperref[term:#1]{\small p.~\pageref*{term:#1}}\par\nopagebreak[4]\noindent#3\par}
\title{Robustness and Typed Hypernet Contexts\\\large A mathematical commentary on the Lean development}
\author{Companion to \texttt{danghica/hypernets-lean}}
\date{24 October 2026\\\small Source snapshot: \texttt{fe1b890}; Lean \texttt{4.35.0}}
\begin{document}
\maketitle
\begin{abstract}
This document explains every definition and proved statement in the current Lean development, including its regression examples. It gives proofs in mathematical English and links each declaration to its complete source fragment at a fixed Git revision. The abstract counting-simulation soundness theorem is proved outright. The forward and converse robustness arguments are proved relative to an explicit specimen interface. The concrete development constructs finite typed graphs, proves well-formed shallow plugging, path and isomorphism results, operation-path uniqueness, finite boxed hierarchies, addressed deep replacement, and globally unique-hole context closure under explicit freshness conditions. Plugging at every finite box depth has an exact syntactic hole inventory and a necessary-and-sufficient freshness criterion, including exact consumption of the selected hole label. Recursive isomorphism now permits independent renaming inside boxes, preserves operation paths and context validity, and makes substitution well-defined on replacement isomorphism classes. Sites now transport across recursively isomorphic hosts, with simultaneous host/replacement congruence and invariance of exact freshness. Nested substitutions now associate at arbitrary box depth, corollas satisfy unit laws, and exact composite context closure permits repair of intermediate hole collisions. The focussed state space is now defined, and the boundary argument of Lemma~39(1) is proved for finite sequences of deep focus-free fillings. The five search rules are now defined on underlying graph coordinates, with proofs of search determinism, boundary laws, and preservation of context steps under shallow filling. Typed search coordinates now have explicit focus-edge state realizations, with proved incidence recovery and initial attachment. General erasure now applies to arbitrary closed focussed states, with a focus-preserving reinsertion isomorphism. Search on explicit states is defined through erasure, proved deterministic up to focussed isomorphism, and satisfies the initial/final search boundary laws. Both focus round trips and recursive renaming invariance are proved. Search descends to literal-label and recursive state classes; projection between them preserves and lifts executions with exact lengths and reflects search termination. Concrete rootedness and refocusing are now defined. Search preserves rootedness; rooted focus sources have operation paths from the unique input, and any incoming non-operation path rules out rootedness, giving concrete versions of Lemmas 23 and 26(1). One-way hypernets now have their term-to-term path definition, renaming invariance, and filling closure. Extraction of operation paths through a filling proves the single-hole rootedness obstruction of Lemma 26(2), including self-looping holes. Adequacy for contextual local rewriting and the full UAM remain open. An explicit dependency map distinguishes these checked results from the remaining operational obligations. The connection to a concrete universal abstract machine and the verification of the specimen interface remain outstanding; the development is therefore not yet a formalization of the paper's full Theorem~6.
\end{abstract}
\clearpage
\begingroup
\fontsize{9.5}{10.5}\selectfont
\setlength{\parskip}{0pt}
\tableofcontents
\endgroup
\clearpage
\Needspace{18\baselineskip}
\section{Scope, notation, and the status of the formalization}

The abstract part of the development follows the proof in Appendix~F of
\cite{paper}. Its completed mathematical core is the soundness of counting
simulation, corresponding to Proposition~37. The remaining argument is expressed
through a structure of abstract specimens. Given that structure and its laws,
Lean proves both directions of the simulation construction and their contextual
consequences. Constructing this structure from the paper's actual hypernets is
still a separate task. The concrete graph part now supplies finite interfaced
graphs, certified shallow plugging, renaming and path invariance, a typed boxed
hierarchy, and addressed deep replacement. These constructions and their
regression examples are explained after the abstract development. The context
calculus now also proves exact hole inventories, necessary-and-sufficient
freshness, and exact consumption of selected hole labels at arbitrary box depth.
Recursive isomorphism now identifies independently renamed box contents and
respects paths, hole observations, and deep replacement. Sites now transport across
isomorphic hosts, and substitution respects simultaneous host and replacement
renaming, including the exact freshness condition. Nested substitutions now
associate at arbitrary box depth, with corolla unit laws and a sharp final
context-validity criterion even when an intermediate collision is later removed.

The concrete development now also defines globally focus-free hypernets,
focussed hypernets, and the hole-free state space. The boundary argument of
Lemma~39(1) is proved for finite sequences of addressed focus-free fillings;
see Section~\ref{sec:concrete-focus}. The five search rules are now defined on underlying graph coordinates, with
determinism, boundary laws, and preservation under shallow filling proved
(Section~\ref{sec:search}). Explicit focus insertion now realizes typed
coordinates as actual states, preserving holes and boundary observations, and
contracting each inserted focus recovers original incidences
(Section~\ref{sec:realization}). General erasure and the focus-preserving erasure-then-insertion isomorphism are
now proved for arbitrary closed focussed states. Erasure-based search on actual
states is deterministic up to focussed isomorphism and has the expected search
boundary laws (Section~\ref{sec:state-search}). Both focus round trips and
source/target renaming invariance are proved; search and boundary observations
descend to literal-label state isomorphism classes, with equality-based search
determinism. Search also respects recursive renaming inside boxes and descends
to recursive state classes. The map from literal to recursive classes preserves
and lifts search executions with exact lengths, including successful search
termination. Canonical roots and concrete rootedness are now defined; search
preserves rootedness, and concrete rooted focus sources have operation paths
from the unique input. An incoming path containing a non-operation edge rules
out rootedness (Section~\ref{sec:rooted-path-obstruction}). Concrete refocusing
is defined and proved for search; Section~\ref{sec:uam-structure} now proves
the complete UAM structural theorem. One-way hypernets and the single-hole
obstruction of Lemma~26(2) are now proved, using extraction of actual operation
paths through shallow substitution (Section~\ref{sec:one-way-obstruction}).
Simultaneous shallow expansion now has proved graph invariants. Concrete
pre-templates, output-closure, and shallow specimens provide actual filled
states and a multi-hole proof of Lemma~26(3) for this representation
(Section~\ref{sec:shallow-specimens}). Initial/final agreement now applies
directly to these specimen states. Every finite simultaneous shallow filling
now has a finite ordinary addressed-substitution realization, and completed
schedules give isomorphic results (Section~\ref{sec:sequential-expansion}).
Independent deep substitutions now have constructed residual sites and a
recursive interchange isomorphism. A quotient of hosts with selected sites
supports substitution and sharp context validity (Section~\ref{sec:deep-interchange}).
Recursive deep specimens now supply actual filled states, with derived
focus freedom and the Lemma~26(3), boundary, and jointly rooted context-path
conclusions (Section~\ref{sec:deep-specimens}). Recursive fillings now preserve
context search with exact execution lengths. Search at host coordinates reflects
to context search away from shallow hole boundaries; invariant regions support
reflection of entire runs. These coordinate executions also realize paired
executions and rooted endpoints on actual state classes
(Section~\ref{sec:expansion-search}). General erasure now commutes with
expansion, identifying these realizations with the original deep specimen
states. Concrete entering/exiting predicates agree with the coordinate boundary
test, output closure and joint rootedness imply non-exiting, and settled
specimens have synchronized actual search transitions to rooted endpoints at
a remote coordinate (Section~\ref{sec:specimen-traversal}). Section~\ref{sec:concrete-specimen-system} now constructs the named-state interface and safety/robustness bridges.
Section~\ref{sec:tuple-safety} proves compatibility with finite original-hole
tuples and defines concrete quasi-specimens, input-safety, and robustness.
Contextual search adequacy is proved below. Concrete substitution and compute
transitions now assemble into the full UAM, with mode discipline, disjoint
transition families, and final normality proved (Section~\ref{sec:full-uam}).

\Needspace{18\baselineskip}
\subsection{Paper-to-Lean proof dependency map}
The target is the proof of Theorem~6 through Propositions~40, 37, and 38.
The following map separates checked arguments from their missing concrete
interpretations, rather than treating graph infrastructure as progress on every
part of that proof.
\begin{enumerate}
\item \textbf{Proposition~37: checked.}\par
\texttt{Machine.CountingSimulation.sound} proves counting-simulation soundness
for a deterministic machine and a reasonable cost triple. Its proof follows
the quantitative argument, including the eventual-stuck alternative.
\item \textbf{Proposition~40: concrete forward simulation checked.}\par
\texttt{SpecimenSystem.countingSimulation} and its converse construct counting
simulations from the abstract specimen laws, safety, and robustness. The abstract laws have concrete named-state and class-valued instances. Section~\ref{sec:uam-classes} proves literal full-class determinism and the concrete forward counting simulation using class safety and robustness. Section~\ref{sec:class-converse} completes the concrete converse counting simulation and two-sided refinement. Section~\ref{sec:concrete-observations} completes those observations and both contextual clauses for the fixed-depth machine. Operational transfer to the unbounded machine remains open.
\item \textbf{Lemma~39(1), Proof~F.4: boundary argument checked.}\par
\texttt{Focussed.fillings\_boundary\_agreement} proves initial/final agreement
for two finite filling sequences of one focussed hypernet. The proof preserves
the focus mode and input status of its source, as in the paper. Finite original-hole tuples and the recursive representation are now
related in both directions in Section~\ref{sec:tuple-safety}.
Both the shallow and recursive deep specimen representations now have direct
proofs of this boundary conclusion for their actual filled states.
\item \textbf{Lemma~39(2): checked for recursive specimens; point~(3): full-UAM locality checked.}\par
Concrete rootedness and output closure yield the non-operation-path obstruction
of Lemma~26(3) for actual recursive specimen states. An exiting focus would give
a singleton hole path to the focus source, so joint rootedness rules it out,
proving the concrete conclusion of point~(2).

General erasure commutes with simultaneous expansion, preserving the joined
coordinate. Consequently expansion of an already focussed context is
focus-preservingly isomorphic to filling its erasure and then inserting a focus.
Entering/exiting adjacency in the original context is equivalent to the checked
coordinate boundary test. Under the settled hypotheses of point~(3), search
enabled on either specimen side yields synchronized actual-state steps to rooted
endpoints at a remote coordinate. The proof allows recursive endpoint renaming
and uses output closure to exclude exiting backtracking focuses.

Section~\ref{sec:concrete-specimen-system} now packages that coordinate as
an allowed unique-hole specimen with recursive fillings and closed endpoints.
Its full-UAM context-step theorem constructs the abstract locality field.
Original-hole tuple compatibility is proved in Section~\ref{sec:tuple-safety};
contextual search adequacy is established in Section~\ref{sec:contextual-search}.
\item \textbf{Paper \S6.4: concrete transition definition checked.}\par
\texttt{State.Substitution} applies the explicit Fig.~9(f) fragments in a
shallow simple context. \texttt{State.Compute} specifies the required active
operation and focus modes; \texttt{Behaviour} restricts the extrinsic parameter
to this shape. \texttt{UAM.machine} assembles all three families, with proved
final normality and transition-family exclusion. This removes the transition
definition obligation.
\item \textbf{Lemma~24, Lemmas~27--28, and Lemma~3: structural theorem checked.}\par
Section~\ref{sec:uam-structure} proves chronological search revisiting,
stationary prefix replay, substitution rootedness and uniqueness, and
full-UAM determinism and refocusing exactly under the corresponding extrinsic
behaviour conditions. Full executions preserve rootedness.
\item \textbf{Paper \S\S8.1--8.3: tuple compatibility and definitions checked.}\par
\texttt{DeepSpecimen.substitutionTuple} constructs distinct, complete finite
original-hole tuples with related closed entries. \texttt{TupleFilling.specimen}
reconstructs a recursive specimen with the same filled graphs. Concrete
quasi-specimens retain unconditional witness existence and mode equality;
input-safety is unrestricted by rootedness, and robustness retains its
unrooted alternative. The rooted quasi and robustness consequences are proved.
Section~\ref{sec:concrete-specimen-system} proves the safety and robustness bridges for the constructed named-state \texttt{SpecimenSystem}.
\item \textbf{Proposition~38: logical lifting checked; initial bridge open.}\par
\texttt{Observations.contextual\_of\_lifting} proves the logical implication
once initial pairs belong to the lifting. Concrete initial focus attachment is
now constructed, including initial rooted specimens obtained by relocation.
The observation interface and its arbitrary-fragment witnesses must still
discharge the \texttt{initial\_lift} parameter.
\item \textbf{Theorem~6: conditional, not complete.}\par
The two \texttt{contextual\_sufficiency} theorems combine these abstract pieces.
The paper itself assumes determinism and refocusing of the UAM. These remain
the theorem's hypotheses. Section~\ref{sec:uam-structure} now defines their
concrete meaning and proves the full-UAM structural theorem, reducing both
properties exactly to the corresponding extrinsic behaviour conditions.
\end{enumerate}

The typed finite-depth hierarchy and explicit renaming bijections are
representation choices. Outer set-presentation normalization, recursive annotated round trips, and
compatibility across depth bounds are now checked in
Section~\ref{sec:set-presentations}. Section~\ref{sec:paper-encoding} resolves the remaining box-label correspondence
by fresh naming of actual contents, with inverse class maps and exact depth
coherence. The checked counterexample to naive erasure is repaired without
restricting the typed graph class. The existing plugging,
path, inventory, and isomorphism lemmas support that representation; they do
not replace the operational proofs in Appendix~F. The contextual-search dependency is now discharged:
Section~\ref{sec:contextual-search} identifies erasure-based state search with
finite interaction patterns plugged into explicit simple contexts, in both
directions. Opening and closing the selected slot covers arbitrary simple
contexts; output fragments and self-looping attachments are included.
The transitions of \S6.4 are now assembled into a concrete machine in
Section~\ref{sec:full-uam}. Section~\ref{sec:uam-structure} proves substitution
rootedness and uniqueness, stationary replay, and the full-UAM structural laws.
Section~\ref{sec:concrete-specimen-system} closes the successor packaging,
full-UAM locality, and concrete safety/robustness bridges. The remaining
operational connection is transfer to the union of all finite depths. Section~\ref{sec:uam-classes} constructs the deterministic full
class machine and its concrete specimen system, class safety/robustness
bridges, forward counting simulation, and remote-specimen refinement.

Section~\ref{sec:class-converse} supplies the concrete converse quasi and robustness
bridges, reverse counting simulation, two-sided refinement, and equivalence
for constructed initial specimens.

Section~\ref{sec:concrete-observations} constructs template-independent
observations, proves initial lifting, and instantiates both contextual
sufficiency clauses, including arbitrary single fragments and nested holes.
The fixed-depth concrete theorem is complete. The scope audit identifies
operational transfer to the paper's unbounded carrier as the remaining
mathematical obligation; the existing graph encoding does not prove it.
The earlier 94\% planning estimate should not be read as a measured completion
percentage. Its remaining-effort allowance did not explicitly separate this
operational transfer, so a more precise overall percentage is not justified.
The active to-do list is maintained in Appendix~\ref{app:remaining-work}.

There are three distinct kinds of assertion below. A \emph{definition} specifies
a mathematical object or predicate. A \emph{proved statement} is established by
a Lean proof term. An \emph{interface law} is a hypothesis carried by a
structure; the current code uses that hypothesis but does not prove it for
concrete hypernets. The concrete instances proved later in this document discharge those
interface laws for the fixed-depth universal abstract machine.

All code links point to revision \texttt{a3adde9}, rather than to the moving
\texttt{main} branch. Each link highlights the declaration together with its
proof or defining fields. Namespace prefixes in link labels disambiguate
declarations with the same short name. Anonymous Lean \texttt{example}
declarations are explained alongside the named theorems.

\Needspace{17\baselineskip}
\subsection{Conventions}

Fix a set of states $X$ and a machine $M$ on $X$. Write $s\to t$ for its
one-step transition relation, $s\steps{n}t$ for exactly $n$ transitions, and
$\fin(s)$ for successful finality. Natural numbers include zero. A cost
preorder $Q$ is written infix: $n\,Q\,m$. The source count is always the
left argument and the target count the right argument. Write $s\preceq_Q t$
for state refinement, and $s\simeq_Q t$ for mutual state refinement.

The paper's three cost parameters $(Q,Q',Q'')$ are named $(Q,P,S)$ in the
code and in this commentary. Thus $P$ is the auxiliary relation on the
\emph{left} of a simulation candidate, and $S$ is the one on its
\emph{right}. We spell out witnesses rather than relying on a convention for
composition of relations.

No result assumes that all executions terminate. An arbitrary state can be
final, stuck, or capable of further computation. Successful termination is the
only qualitative observation in the present state-refinement definition.
Neither equality of returned values nor equality of final states across two
different source states is required.

\Needspace{17\baselineskip}
\subsection{Dependency order}

The files are organized as follows. Finite-execution facts in
\leanfile{Hypernets/Transition.lean}{Transition.lean} support the quantitative
relations in \leanfile{Hypernets/Quantitative.lean}{Quantitative.lean} and
the soundness argument in
\leanfile{Hypernets/Simulation.lean}{Simulation.lean}. The observation interface
in \leanfile{Hypernets/Contextual.lean}{Contextual.lean} converts state
results into contextual ones. The specimen argument in
\leanfile{Hypernets/LocalReasoning.lean}{LocalReasoning.lean} constructs the
simulation to which soundness applies. Finally,
\leanfile{Hypernets/Examples.lean}{Examples.lean} checks substantive edge
cases, and \leanfile{Hypernets/Audit.lean}{Audit.lean} reports theorem
dependencies. The root file \leanfile{Hypernets.lean}{Hypernets.lean} imports
the development and its examples; it adds no mathematical declaration.

The concrete branch begins with
\leanfile{Hypernets/Graph/Finite.lean}{Graph/Finite.lean} and
\leanfile{Hypernets/Graph/Basic.lean}{Graph/Basic.lean}. Plugging, paths, and
isomorphism support the hierarchy and deep replacement modules. Operation-path
uniqueness imports the counted-execution results, providing a concrete use of the
abstract machine lemmas. Occurrence fibres and provenance then support shallow and deep context closure
in the new context modules. Both branches are imported by the root file.

The exposition follows these dependency orders. Section numbers and theorem
numbers here are local to this document; references to the paper are explicitly
identified as such.

\Needspace{18\baselineskip}
\section{Machines and finite executions}

\begin{definition}[Machine]\label{term:machine}
A machine on $X$ consists of a relation ${\to}\subseteq X\times X$, a
predicate $\fin\subseteq X$, and the law
\[
  \fin(s)\ \Longrightarrow\ \neg(s\to t)\qquad(s,t\in X).
\]
Thus every successful final state has no outgoing transition. No converse is
assumed: an unsuccessful normal form is permitted.
\end{definition}
\lean{Transition}{12}{15}{Machine}
The fields \texttt{step} and \texttt{final} supply the relation and predicate.
The field \texttt{final\_normal} supplies the displayed law. A machine may
have any state type; finiteness or a graphical representation is not required.

\begin{definition}[Determinism, normality, and stuckness]\label{term:determinism}
The machine is \emph{deterministic} if $s\to t_1$ and $s\to t_2$ imply
$t_1=t_2$. A state $s$ is \emph{normal} if $s\to t$ is false for every
$t$. It is \emph{stuck} if it is normal and is not final:
\[
 \normal(s)\iff\forall t\,\neg(s\to t),\qquad
 \stuck(s)\iff\normal(s)\land\neg\fin(s).
\]
\end{definition}
\lean{Transition}{21}{22}{Machine.Deterministic}
\lean{Transition}{24}{24}{Machine.Normal}
\lean{Transition}{26}{27}{Machine.Stuck}
Determinism is a separate hypothesis, not a field of every machine. Stuckness
distinguishes failure from successful finality even though both are normal
forms.

\begin{definition}[Counted execution]\label{term:execution}
The relation $s\steps{n}t$ is generated inductively by
\[
 \frac{}{s\steps{0}s},\qquad
 \frac{s\to u\quad u\steps{n}t}{s\steps{n+1}t}.
\]
\end{definition}
\lean{Transition}{30}{32}{Machine.Steps}
These are precisely the constructors \texttt{zero} and \texttt{cons}. The
count includes every transition of the machine. There is no separate cost
weight for searching or rewriting. In particular, zero transitions cannot
change a state.

\begin{lemma}[Zero-step equality]
If $s\steps{0}t$, then $s=t$.
\end{lemma}
\lean{Transition}{36}{38}{Machine.Steps.zero_eq}
\begin{proof}
Only the zero constructor can produce an execution with count zero. Its two
endpoints are the same state. The step constructor has a positive count and
cannot occur.
\end{proof}

\begin{lemma}[Concatenation]
If $s\steps{m}t$ and $t\steps{n}v$, then $s\steps{m+n}v$.
\end{lemma}
\lean{Transition}{40}{45}{Machine.Steps.append}
\begin{proof}
Induct on the first execution. For a zero-step execution, $s=t$, so the second
execution is the required one. Otherwise its first transition is $s\to u$
and its remaining execution has some length $k$. By induction, that remainder
followed by the second execution has length $k+n$. Prepending $s\to u$
gives length $(k+n)+1=(k+1)+n$. No determinism is used.
\end{proof}

\begin{lemma}[An execution from a normal state]
If $s\steps{n}t$ and $\normal(s)$, then $n=0$ and $s=t$.
\end{lemma}
\lean{Transition}{47}{51}{Machine.Steps.normal_eq}
\begin{proof}
The zero-step case has both conclusions by definition. A positive execution
would begin with a transition out of $s$, contradicting normality.
\end{proof}

For the following three lemmas, assume determinism.

\begin{lemma}[Every execution is a prefix of a final execution]
Suppose $s\steps{m}t$, $s\steps{n}z$, and $\fin(z)$. Then there is a
natural number $r$ such that
\[
 n=m+r\quad\text{and}\quad t\steps{r}z.
\]
In particular $m\leq n$.
\end{lemma}
\lean{Transition}{54}{67}{Machine.Steps.prefix_of_final}
\begin{proof}
Induct on the execution of length $m$, allowing $n$ to vary. If $m=0$, then
$t=s$; take $r=n$ and use the given final execution.

Otherwise the first execution starts with $s\to u$. The final execution
cannot have length zero: that would make $s=z$ final while an outgoing
transition exists. It therefore starts with $s\to u'$. Determinism gives
$u=u'$. Apply the induction hypothesis to the two remaining executions. If
their lengths are $m-1$ and $n-1$, the result gives
$n-1=(m-1)+r$ and the required execution $t\steps{r}z$. Hence $n=m+r$.
The Lean proof uses constructor indices rather than natural-number
subtraction, so no subtraction side conditions are needed.
\end{proof}

\begin{lemma}[Uniqueness of successful terminal executions]
If $s\steps{m}t$, $s\steps{n}z$, $\fin(t)$, and $\fin(z)$, then
$m=n$ and $t=z$.
\end{lemma}
\lean{Transition}{69}{74}{Machine.Steps.final_unique}
\begin{proof}
The prefix lemma gives $n=m+r$ and $t\steps{r}z$. Finality of $t$ makes
it normal. The normal-state lemma yields $r=0$ and $t=z$, and therefore
$m=n$.
\end{proof}

\begin{lemma}[A state on a successful execution is not stuck]
If $s\steps{m}t$, $s\steps{n}z$, and $\fin(z)$, then $t$ is not stuck.
\end{lemma}
\lean{Transition}{76}{82}{Machine.Steps.not_stuck_of_final}
\begin{proof}
Suppose $t$ were stuck. The prefix lemma gives a remaining execution
$t\steps{r}z$. Normality of $t$ forces $t=z$. Since $z$ is final, $t$
is final, contradicting the unsuccessful part of stuckness.
\end{proof}

\begin{lemma}[Removing a prescribed first transition]
If $s\to s'$ and $s\steps{n+1}t$, then $s'\steps{n}t$.
\end{lemma}
\lean{Transition}{84}{90}{Machine.Steps.drop_first}
\begin{proof}
The positive execution has a first transition $s\to u$ and a remainder
$u\steps{n}t$. Determinism identifies $u$ with $s'$, giving the stated
remainder. The conclusion concerns the \emph{prescribed} successor $s'$;
without determinism it would only be valid for some successor.
\end{proof}

\begin{lemma}[A stuck execution continues after an available step]
If $s\to s'$, $s\steps{n}t$, and $\stuck(t)$, then for some $k$,
$s'\steps{k}t$.
\end{lemma}
\lean{Transition}{93}{101}{Machine.Steps.stuck_after_step}
\begin{proof}
The execution cannot have length zero, since then $s=t$ would be normal
despite $s\to s'$. Its first transition is therefore $s\to u$, with a
remaining execution to $t$. Determinism gives $u=s'$. Take the length of
that remainder as $k$. The code does not retain the additional equality
$n=k+1$, because later uses need only the existence of a remainder.
\end{proof}

\Needspace{18\baselineskip}
\section{Cost relations and state refinement}

\begin{definition}[Cost preorder]\label{term:cost-preorder}
A cost order is a reflexive, transitive relation $Q$ on $\mathbb N$.
\end{definition}
\lean{Quantitative}{6}{9}{CostOrder}
The relation and its two preorder proofs are stored together. Antisymmetry,
decidability, and compatibility with addition are not part of this definition.

\begin{definition}[Reasonable triple]\label{term:reasonable}
A triple $(Q,P,S)$ of cost preorders is \emph{reasonable} if all four of the
following implications hold:
\begin{align*}
 a\,Q\,b\ \land\ c\,Q\,d &\ \Longrightarrow\ (a+c)\,Q\,(b+d),\tag{addition}\\
 a\,P\,b &\ \Longrightarrow\ b\leq a,\tag{decrease}\\
 a\,P\,b &\ \Longrightarrow\ a\,S\,b,\tag{inclusion}\\
 a\,P\,b\ \land\ b\,Q\,c\ \land\ c\,S\,d
   &\ \Longrightarrow\ a\,Q\,d.\tag{sandwich}
\end{align*}
\end{definition}
\lean{Quantitative}{12}{16}{Reasonable}
These four fields express the paper's three conditions: the decrease and
inclusion fields are the two parts of its condition~(B). The decrease field
will justify a strict descent in the soundness proof. Inclusion is used when
turning safety's auxiliary equivalences into an asymmetric up-to relation.

\Needspace{23\baselineskip}
\begin{definition}[Converse and the basic cost orders]\label{term:cost-orders}
The converse preorder is $a\,Q^{-1}\,b$ if $b\,Q\,a$. The four named
orders are
\[
 \begin{aligned}
 a\,E\,b &\iff a=b,&
 a\,I\,b &\iff b\leq a,\\
 a\,H\,b &\iff a\leq b,&
 a\,U\,b &\iff \text{true}.
 \end{aligned}
\]
Here $E$, $I$, $H$, and $U$ denote \texttt{exact}, \texttt{improvement},
\texttt{overhead}, and \texttt{any}, respectively.
\end{definition}
\lean{Quantitative}{20}{23}{CostOrder.converse}
\lean{Quantitative}{25}{25}{CostOrder.exact}
\lean{Quantitative}{27}{28}{CostOrder.improvement}
\lean{Quantitative}{30}{30}{CostOrder.overhead}
\lean{Quantitative}{32}{32}{CostOrder.any}
\begin{proof}[Verification of the preorder fields]
Converse preserves reflexivity. Given $b\,Q\,a$ and $c\,Q\,b$,
transitivity of $Q$ gives $c\,Q\,a$, establishing transitivity of
$Q^{-1}$. Equality is reflexive and transitive. The relation $I$ is
reflexive by $a\leq a$; from $b\leq a$ and $c\leq b$ one obtains
$c\leq a$. The corresponding verification for $H$ uses the usual order
direction. All requirements for the universal relation $U$ are immediate.
\end{proof}

\begin{lemma}[Exact costs are reasonable]
The triple $(E,E,E)$ is reasonable.
\end{lemma}
\lean{Quantitative}{34}{38}{CostOrder.reasonable_exact}
\begin{proof}
Substitute equals for equals in the addition condition. The decrease
condition reduces to $a\leq a$, inclusion is equality included in itself,
and the sandwich condition is transitivity of equality.
\end{proof}

\begin{lemma}[Improvement costs are reasonable]
The triple $(I,I,I)$ is reasonable.
\end{lemma}
\lean{Quantitative}{40}{44}{CostOrder.reasonable_improvement}
\begin{proof}
For addition, $b\leq a$ and $d\leq c$ give $b+d\leq a+c$. Decrease is
the definition of $I$, and inclusion is reflexive inclusion. The sandwich
premises give $d\leq c\leq b\leq a$, hence $a\,I\,d$.
\end{proof}

\begin{lemma}[Unconstrained termination costs are reasonable]
The triple $(U,I,U)$ is reasonable.
\end{lemma}
\lean{Quantitative}{46}{50}{CostOrder.reasonable_termination}
\begin{proof}
Addition, inclusion, and the sandwich condition have universally true
conclusions. Decrease is again the definition of $I$.
\end{proof}

\begin{lemma}[Overhead with exact auxiliary costs is reasonable]
The triple $(H,E,E)$ is reasonable.
\end{lemma}
\lean{Quantitative}{52}{56}{CostOrder.reasonable_overhead}
\begin{proof}
The addition condition is monotonicity of addition for $\leq$. Equality
implies the required weak decrease and includes into itself. In the sandwich
condition, substitute $a=b$ and $c=d$ into $b\leq c$ to obtain $a\leq d$.
\end{proof}

\begin{definition}[State refinement]\label{term:refinement}
For a cost preorder $Q$, define
\[
 s\preceq_Q t\quad\Longleftrightarrow\quad
 \begin{gathered}
 \text{for every $n,z$ with $s\steps{n}z$ and $\fin(z)$,}\\
 \text{there exist $m,w$ with $n\,Q\,m$, $t\steps{m}w$, and $\fin(w)$.}
 \end{gathered}
\]
State equivalence is mutual refinement:
$s\simeq_Q t$ if $s\preceq_Q t$ and $t\preceq_Q s$.
\end{definition}
\lean{Quantitative}{65}{67}{Machine.Refines}
\lean{Quantitative}{69}{70}{Machine.Equiv}
If $s$ has no successful execution, its refinement obligation is vacuous.
The definition therefore does not distinguish stuckness from divergence
unless a surrounding observation supplies additional distinctions.
For $Q=I$, the target must terminate in no more steps than the source.
For $Q=H$, it must terminate in at least as many steps. For $Q=E$, the
counts agree exactly; for $Q=U$, counts are ignored.

\begin{lemma}[Reflexivity]
For every $s$, $s\preceq_Q s$.
\end{lemma}
\lean{Quantitative}{74}{76}{Machine.Refines.refl}
\begin{proof}
Given a successful execution of length $n$ from $s$, choose that same
execution and final state as witnesses. Reflexivity supplies $n\,Q\,n$.
\end{proof}

\begin{lemma}[Transitivity]
If $s\preceq_Q t$ and $t\preceq_Q v$, then $s\preceq_Q v$.
\end{lemma}
\lean{Quantitative}{78}{83}{Machine.Refines.trans}
\begin{proof}
Take a successful execution of length $n$ from $s$. The first refinement
supplies a successful execution of length $m$ from $t$, with $n\,Q\,m$.
Apply the second refinement to this execution to obtain a successful
execution of length $k$ from $v$, with $m\,Q\,k$. Transitivity gives
$n\,Q\,k$, and these are the required witnesses. Determinism is unnecessary.
\end{proof}

\begin{lemma}[Monotonicity in the cost relation]
If $Q\subseteq P$ and $s\preceq_Q t$, then $s\preceq_P t$.
\end{lemma}
\lean{Quantitative}{85}{90}{Machine.Refines.mono}
\begin{proof}
Reuse each execution and final-state witness supplied by the original
refinement. Replace its cost proof $n\,Q\,m$ by $n\,P\,m$ using the inclusion.
\end{proof}

\Needspace{18\baselineskip}
\section{Counting simulation and its soundness}

\begin{definition}[Up-to closure]
Given a relation $R$ on states and cost orders $P,S$, define
\[
 \upto{P}{S}{R}(s,t)\quad\Longleftrightarrow\quad
 \exists u,v\;\bigl(s\preceq_P u\ \land\ R(u,v)\ \land\ v\preceq_S t\bigr).
\]
\end{definition}
\lean{Simulation}{10}{11}{Machine.UpTo}
The witnesses occur in path order. In particular the right-hand refinement
is from $v$ to $t$, not from $t$ to $v$. Auxiliary refinement can change the
states at which the candidate relation is resumed.

\begin{definition}[Counting simulation]
A relation $R$ is a $(Q,P,S)$-counting simulation if whenever $R(s,t)$:
\begin{enumerate}[label=\textup{(\Alph*)}]
\item If $s$ is final, then $t$ is final.
\item For every transition $s\to s'$, one of the following holds:
\begin{enumerate}[label=\textup{(\roman*)}]
\item There exist $k,z$ such that $s'\steps{k}z$ and $z$ is stuck.
\item There exist $k,l,h_1,h_2$ such that
\[
 s'\steps{k}h_1,\qquad t\steps{l}h_2,\qquad
 (1+k)\,Q\,l,\qquad \upto{P}{S}{R}(h_1,h_2).
\]
\end{enumerate}
\end{enumerate}
\end{definition}
\lean{Simulation}{14}{19}{Machine.CountingSimulation}
The leading $1$ charges for the already prescribed transition $s\to s'$.
Neither $k$ nor $l$ is required to be positive. The stuck alternative is a
deliberate part of the definition, consistent with observing only successful
termination. The finality condition is stronger than simply allowing the
target to reach a final state after some extra steps: it requires the target
to be final already when the source is final.

\begin{theorem}[Soundness of counting simulation]\label{thm:soundness}
Suppose $M$ is deterministic, $(Q,P,S)$ is reasonable, and $R$ is a
$(Q,P,S)$-counting simulation. Then
\[
 R(s,t)\ \Longrightarrow\ s\preceq_Q t.
\]
\end{theorem}
\lean{Simulation}{24}{51}{Machine.CountingSimulation.sound}
This is the abstract transition-system form of the paper's Proposition~37.
Its proof is a strong induction on the length of a successful source
execution. Although simulation provides a step-wise, coinductive style of
reasoning, the soundness theorem itself is proved by induction on a finite
observation; no coinductive datatype is introduced in this file.

\begin{proof}
We prove, simultaneously for all states $s,t,z$, that if $R(s,t)$,
$s\steps{N}z$, and $\fin(z)$, then there are $m,w$ such that
$N\,Q\,m$, $t\steps{m}w$, and $\fin(w)$. Use strong induction on $N$;
the induction hypothesis applies to every smaller natural number and every
pair of related states.

\emph{Base case.} If $N=0$, then $s=z$ is final. Condition~(A) makes
$t$ final. Choose the zero-step execution at $t$ and use $0\,Q\,0$.

\emph{Positive case.} Write $N=n+1$ and decompose the successful execution:
\[
 s\to s'\steps{n}z.
\]
Apply condition~(B) of simulation. Its stuck alternative would give a
finite execution from $s'$ to a stuck state. This is impossible by the
lemma that a state on a successful execution cannot be stuck, applied to
the successful execution $s'\steps{n}z$.

Consequently there are $k,l,h_1,h_2,u,v$ satisfying
\begin{gather*}
 s'\steps{k}h_1,\qquad t\steps{l}h_2,\qquad (1+k)\,Q\,l,\\
 h_1\preceq_P u,\qquad R(u,v),\qquad v\preceq_S h_2.
\end{gather*}
The prefix lemma, using determinism and finality of $z$, gives a remainder
length $r$ such that
\[
 n=k+r,\qquad h_1\steps{r}z.
\]
Apply $h_1\preceq_P u$ to this remaining successful execution. It supplies
$a,z_a$ with
\[
 r\,P\,a,\qquad u\steps{a}z_a,\qquad \fin(z_a).
\]
The decrease condition of reasonableness gives $a\leq r$. Since
$N=1+k+r$, we have $a<N$. This is the strict inequality needed to invoke
the strong induction hypothesis. The proof would not be justified for
arbitrary auxiliary cost preorders.

Applying that hypothesis to $R(u,v)$ gives $b,z_b$ such that
\[
 a\,Q\,b,\qquad v\steps{b}z_b,\qquad \fin(z_b).
\]
Now apply $v\preceq_S h_2$. We obtain $c,z_c$ with
\[
 b\,S\,c,\qquad h_2\steps{c}z_c,\qquad \fin(z_c).
\]
The sandwich condition combines $r\,P\,a$, $a\,Q\,b$, and $b\,S\,c$
to yield $r\,Q\,c$. Additivity of $Q$ then gives
\[
 (1+k+r)\,Q\,(l+c),\quad\text{that is,}\quad N\,Q\,(l+c).
\]
Concatenate $t\steps{l}h_2$ with $h_2\steps{c}z_c$. The result is a
successful execution of length $l+c$ from $t$, which supplies the required
witnesses.
\end{proof}

The code follows these witnesses exactly: its \texttt{count} is $n=k+r$,
\texttt{smaller} is $a<n+1$, \texttt{rc} is $r\,Q\,c$, and
\texttt{total} is $(n+1)\,Q\,(l+c)$. The inclusion $P\subseteq S$ is
not needed in this particular proof; it is part of the shared reasonable
triple and will be used by the specimen argument.

\Needspace{18\baselineskip}
\section{Typed observations and contextual lifting}

\begin{definition}[Observation interface]
Let $J$ be a type of interfaces, and for each $i\in J$ let $F_i$ be a
type of fragments and $C_i$ a type of contexts. An observation interface
consists of a predicate $\mathsf{Allowed}(c)$ for contexts and a relation
\[
 \mathsf{Obs}(c,a,s),\qquad c\in C_i,\ a\in F_i,\ s\in X.
\]
The intended reading is that placing $a$ in $c$ and initializing the
result can produce state $s$.
\end{definition}
\lean{Contextual}{15}{18}{Observations}
The common index $i$ prevents a fragment and a context of incompatible
interfaces from being compared. Observation is a relation rather than a
total function. It can be undefined, and the structure does not require
uniqueness of the produced state. The current code does not yet define it
by graph plugging.

\begin{definition}[Contextual refinement and equivalence]
For $a,b\in F_i$, write $a\preceq_Q^{\mathcal O} b$ if for every
$c\in C_i$ and states $s,t$,
\[
 \mathsf{Allowed}(c)\ \land\ \mathsf{Obs}(c,a,s)\ \land\ \mathsf{Obs}(c,b,t)
 \ \Longrightarrow\ s\preceq_Q t.
\]
Contextual equivalence $a\simeq_Q^{\mathcal O}b$ is the conjunction of
this refinement and the refinement in the opposite direction.
\end{definition}
\lean{Contextual}{26}{28}{Observations.Refines}
\lean{Contextual}{30}{31}{Observations.Equiv}
The quantifiers compare every pair of observations of the two fragments
in the same allowed context. The definition imposes no existence obligation
when either fragment has no observation there. Also, state-refinement
reflexivity alone does not establish contextual reflexivity for this fully
general interface: two observations of the same fragment might give
different states. A concrete, functional initialization relation can supply
the missing property, but it is not assumed here.

\begin{proposition}[Contextual refinement from a lifted relation]
Let $T_i$ be a relation on $F_i$ for each interface $i$, and let $L$ be
a relation on states. Suppose:
\begin{enumerate}
\item Whenever $T_i(a,b)$ and $c$ is allowed, observations
  $\mathsf{Obs}(c,a,s)$ and $\mathsf{Obs}(c,b,t)$ imply $L(s,t)$.
\item $L(s,t)$ implies $s\preceq_Q t$.
\end{enumerate}
Then $T_i(a,b)$ implies $a\preceq_Q^{\mathcal O}b$.
\end{proposition}
\lean{Contextual}{35}{44}{Observations.contextual_of_lifting}
\begin{proof}
Fix $T_i(a,b)$, an allowed $c$, and the two observations. The first
hypothesis, named \texttt{initial\_lift} in Lean, puts their initial states
in $L$. The second hypothesis, named \texttt{sound}, then gives the state
refinement required by the definition of contextual refinement.
\end{proof}
This isolates the logical step of the paper's Proposition~38. Its concrete
proof still has to establish the initial-specimen hypothesis for actual
plugging and initialization. The proposition above neither asserts nor
requires closure under plugging at this abstract level.

\begin{proposition}[Contextual equivalence from a lifted relation]
Under the same initial-lifting hypothesis, if $L(s,t)$ implies
$s\simeq_Q t$, then $T_i(a,b)$ implies $a\simeq_Q^{\mathcal O}b$.
\end{proposition}
\lean{Contextual}{46}{57}{Observations.contextual_equiv_of_lifting}
\begin{proof}
For the forward refinement, apply the preceding proposition using the
forward component of state equivalence. For the reverse refinement, fix
observations of $b$ and $a$ in an allowed context. Put them into
\texttt{initial\_lift} in the order $a,b$, obtaining $L$ from the state
of $a$ to the state of $b$. Use the reverse component of its state
equivalence. This does not require that $T_i$ or $L$ be symmetric.
\end{proof}

\Needspace{18\baselineskip}
\section{The abstract specimen interface}

\begin{definition}[Focus modes]\label{term:focus}
The type of modes has exactly three constructors:
$\mathsf{search}$, $\mathsf{backtrack}$, and $\mathsf{rewrite}$.
\end{definition}
\lean{LocalReasoning}{15}{17}{Focus}
They represent the paper's searching, backtracking, and triggering focuses.
The derived decidable equality allows their distinctness to be checked by
finite case analysis. This datatype records only a mode, not a graph
position or a focus edge.

\begin{definition}[Specimen system]\label{term:specimen}\label{def:specimen}
A specimen system $G$ over $M$ consists of a type $\Sigma$ of specimens,
maps $\ell,r:\Sigma\to X$, a mode map
$\mode:\Sigma\to\mathsf{Focus}$, a predicate $\ent$ on specimens,
and a predicate $\rt$ on states. Its laws are:
\begin{enumerate}
\item \emph{Rootedness preservation:} $\rt(s)$ and $s\to t$ imply $\rt(t)$.
\item \emph{Final-state agreement:} for every $c\in\Sigma$,
  $\fin(\ell(c))\iff\fin(r(c))$.
\item \emph{Context-only traversal:} suppose both endpoints of $c$ are
  rooted and $c$ is either backtracking or searching without entering.
  If either endpoint has an outgoing transition, there is a specimen $d$
  whose focus is remote and such that
  \[
    \ell(c)\to\ell(d),\qquad r(c)\to r(d).
  \]
\end{enumerate}
The terms ``rooted,'' ``entering,'' and ``remote'' receive their precise
abstract meanings from these data and the definitions immediately below.
\end{definition}
\lean{LocalReasoning}{24}{38}{SpecimenSystem}
The corresponding fields are \texttt{Specimen}, \texttt{left},
\texttt{right}, \texttt{mode}, \texttt{entering}, \texttt{rooted},
\texttt{rooted\_step}, \texttt{same\_final}, and \texttt{context\_step}.
The last three are hypotheses needed to construct an instance, not results
proved here from a graph model. The intended specimen consists of a
multi-hole focussed context and two pointwise related vectors of fragments.
That intended representation is not a field of this abstract structure.

In particular, the interface does not identify modes with a focus field on
states, define rootedness as reachability from an initial state, or define
entering by incidence with a hole. The context-traversal law captures the
consequence of graph locality and output-closure that the simulation proof
uses. Its availability premise is symmetric between the two endpoints,
which will be essential in the converse proof.

\begin{definition}[Remote, settled, and lifted specimens]\label{term:lifting}
A specimen $c$ is remote if
\[
 \remote(c)\iff \mode(c)\ne\mathsf{rewrite}\ \lor\ \neg\ent(c).
\]
It is settled if
\[
 \settled(c)\iff \mode(c)=\mathsf{backtrack}\ \lor\
   (\mode(c)=\mathsf{search}\land\neg\ent(c)).
\]
Its lifted state relation is
\[
 \Lft(s,t)\iff\exists c\in\Sigma\;
  \remote(c)\land\rt(\ell(c))\land\rt(r(c))
  \land s=\ell(c)\land t=r(c).
\]
\end{definition}
\lean{LocalReasoning}{45}{45}{SpecimenSystem.Remote}
\lean{LocalReasoning}{47}{48}{SpecimenSystem.Settled}
\lean{LocalReasoning}{50}{52}{SpecimenSystem.Lifting}
Remote excludes only an entering rewrite focus. It allows an entering
search focus. Settled is stronger: it excludes every rewrite focus and
also excludes an entering search focus. The two equalities in the lifting
definition expose a particular specimen as a witness for the endpoints.

\begin{lemma}[Settled implies remote]
If $\settled(c)$, then $\remote(c)$.
\end{lemma}
\lean{LocalReasoning}{54}{56}{SpecimenSystem.settled_remote}
\begin{proof}
A backtracking or searching mode is distinct from the rewrite mode.
Therefore the first disjunct in the remote predicate holds in either case.
\end{proof}

\begin{lemma}[Rootedness along a finite execution]
If $\rt(s)$ and $s\steps{n}t$, then $\rt(t)$.
\end{lemma}
\lean{LocalReasoning}{58}{62}{SpecimenSystem.rooted_steps}
\begin{proof}
Induct on the execution. At length zero nothing changes. For a positive
execution, the one-step rootedness law makes its first successor rooted;
the induction hypothesis propagates this property through the remainder.
\end{proof}

\begin{definition}[Abstract non-rewrite quasi-specimen]\label{term:quasi}
For relations $A,B$ on states, define $\quasi{A}{B}(s,t)$ to mean that
\emph{if} $s$ and $t$ are rooted, then there is a specimen $c$ with
\[
 \begin{gathered}
 \mode(c)\ne\mathsf{rewrite},\qquad
 \rt(\ell(c)),\quad\rt(r(c)),\\
 A(s,\ell(c)),\qquad B(r(c),t).
 \end{gathered}
\]
\end{definition}
\lean{LocalReasoning}{67}{70}{SpecimenSystem.Quasi}
The outer implications matter: no witness is required if either endpoint
is unrooted. Subsequent proofs establish both rootedness hypotheses before
using a quasi-specimen. Unlike the paper's full definition, this abstract
consequence does not compare focus labels of $s,t$ with that of $c$;
there is no state-focus map in the interface. It specifically packages the
non-rewrite situation used by safety and robustness.

\begin{lemma}[A rooted quasi-specimen supplies an up-to witness]
If $\quasi{\preceq_P}{\preceq_S}(s,t)$ and both $s,t$ are rooted, then
$\upto{P}{S}{\Lft}(s,t)$.
\end{lemma}
\lean{LocalReasoning}{72}{77}{SpecimenSystem.quasi_upTo}
\begin{proof}
Obtain its witness $c$. Its non-rewrite mode makes it remote. Together
with rootedness of its endpoints, this gives $\Lft(\ell(c),r(c))$.
Use $\ell(c)$ and $r(c)$ as the intermediate states in the up-to
definition, with the two refinements furnished by the quasi-specimen.
\end{proof}

\Needspace{24\baselineskip}
\begin{definition}[Input-safety]\label{term:safety}
The predicate $\safe{Q}{P}$ requires the following for every specimen $c$
whose endpoints are rooted, whose mode is search, and which is entering.
At least one of these three cases must hold:
\begin{enumerate}[label=\textup{(\Roman*)}]
\item Both endpoints have finite executions to stuck states. Their counts
  may be unrelated, and either count may be zero.
\item For some settled specimen $d$ and $k,l\in\mathbb N$,
  \[
    \ell(c)\steps{1+k}\ell(d),\quad r(c)\steps{l}r(d),\quad
    (1+k)\,Q\,l.
  \]
\item For some states $a,b$ and $k,l\in\mathbb N$,
  \[
    \ell(c)\steps{1+k}a,\quad r(c)\steps{1+l}b,\quad
    (1+k)\,Q\,(1+l),
  \]
  and $\quasi{\simeq_P}{\simeq_P}(a,b)$ holds.
\end{enumerate}
\end{definition}
\lean{LocalReasoning}{81}{91}{SpecimenSystem.InputSafe}
Case~II requires a positive left execution but permits zero right steps.
Case~III requires positive executions on both sides and auxiliary
\emph{equivalences}. Retaining both directions of those equivalences is
necessary for reusing safety in the converse theorem.

\begin{definition}[Robustness]\label{term:robustness}
The predicate $\robust{Q}{P}{S}$ requires that for every specimen $c$
with rooted endpoints, rewrite mode, and $\neg\ent(c)$, each transition
$\ell(c)\to s'$ satisfies one of:
\begin{enumerate}[label=\textup{(\roman*)}]
\item $s'$ has a finite execution to a stuck state;
\item for some $k,l,a,b$,
  \[
   s'\steps{k}a,\quad r(c)\steps{l}b,\quad (1+k)\,Q\,l,
  \]
  and $\quasi{\preceq_P}{\preceq_S}(a,b)$ holds.
\end{enumerate}
\end{definition}
\lean{LocalReasoning}{95}{101}{SpecimenSystem.Robust}
This is a rooted, remote restriction of the consequences used from the
paper's robustness condition. The paper's alternative that an endpoint is
unrooted is contradictory under the present premises. The abstract
machine does not classify its transition relation into search and rewrite
rules; it is the specimen's mode that determines where this predicate is
invoked. Establishing the correspondence with concrete rewrite transitions
is part of the missing hypernet instance.

\Needspace{18\baselineskip}
\section{The forward robustness argument}

Throughout this section $M$ is deterministic, $(Q,P,S)$ is reasonable,
and $G$ is a specimen system satisfying $\safe{Q}{P}$ and
$\robust{Q}{P}{S}$. In particular, all the interface laws of
Definition~\ref{def:specimen} are available.

\begin{theorem}[Forward construction of a counting simulation]
The relation $\Lft$ is a $(Q,P,S)$-counting simulation.
\end{theorem}
\lean{LocalReasoning}{104}{163}{SpecimenSystem.countingSimulation}
This is an abstract form of the forward clause of the paper's
Proposition~40. The proof checks the two simulation fields directly.

\begin{proof}
Take $\Lft(s,t)$ and choose a witnessing specimen $c$. Thus
$s=\ell(c)$, $t=r(c)$, the specimen is remote, and both endpoints are
rooted.

\emph{Finality.} If $s$ is final, final-state agreement for $c$ makes
$t$ final. This verifies simulation condition~(A).

\emph{A common context-traversal case.} Suppose $c$ is settled and
there is a prescribed transition $s\to s'$. The context-traversal law
gives a remote specimen $d$ and transitions
\[
 \ell(c)\to\ell(d),\qquad r(c)\to r(d).
\]
Determinism identifies $s'$ with $\ell(d)$. Both new endpoints are
rooted, so $\Lft(\ell(d),r(d))$. The simulation matching case can use
zero further left steps and one right step. Its cost obligation is
$(1+0)\,Q\,1$, which is reflexivity of $Q$. To obtain the up-to
relation, choose $\ell(d),r(d)$ as intermediate states and insert
reflexive refinements for $P$ and $S$ on either side of their lifted
relation. These are the local proof witnesses named
\texttt{context\_case} and \texttt{settled\_case} in Lean.

It remains to cover all possible modes of $c$.

\emph{Backtracking.} Such a specimen is settled by definition, so the
common context-traversal argument applies.

\emph{Searching without entering.} This is the other settled case and
uses the same argument.

\emph{Entering search.} Apply input-safety, obtaining one of its three
cases.
\begin{enumerate}[label=\textup{(\Roman*)}]
\item There is a stuck execution from $\ell(c)$. Since
  $\ell(c)\to s'$ is available, the stuck-after-step lemma removes
  that first step and gives a stuck execution from $s'$. This is the
  stuck alternative of simulation. Safety also supplies a stuck execution
  on the right, but the forward simulation condition does not need it.
\item There are $d,k,l$ with $d$ settled,
  \[
    \ell(c)\steps{1+k}\ell(d),\qquad r(c)\steps{l}r(d),\qquad
    (1+k)\,Q\,l.
  \]
  Remove the prescribed first transition from the left execution using
  determinism. The resulting execution from $s'$ to $\ell(d)$ has
  length $k$. Rootedness propagates to both endpoints of $d$, and
  settledness implies remoteness. Hence
  $\Lft(\ell(d),r(d))$. Add reflexive auxiliary refinements as in
  the context-traversal case. The matching counts are $k,l$, with the
  cost condition already supplied by safety.
\item There are $k,l,a,b$ with the positive executions and quasi-specimen
  of safety~III. Remove the prescribed first step from the left execution,
  giving $s'\steps{k}a$; retain the right execution of length $1+l$.
  Both $a$ and $b$ are rooted. Unpack the quasi-specimen to get a
  non-rewrite specimen $d$ with rooted endpoints and
  \[
     a\simeq_P\ell(d),\qquad r(d)\simeq_P b.
  \]
  Use the forward half of the first equivalence to obtain
  $a\preceq_P\ell(d)$. The forward half of the second gives
  $r(d)\preceq_P b$, and monotonicity with $P\subseteq S$ changes
  this to $r(d)\preceq_S b$. Thus the quasi-up-to lemma supplies
  $\upto{P}{S}{\Lft}(a,b)$. The matching counts are $k,1+l$,
  and their required cost relation is exactly
  $(1+k)\,Q\,(1+l)$.
\end{enumerate}

\emph{Rewrite mode.} Remoteness now implies $\neg\ent(c)$, since
the other remote disjunct contradicts the current mode. Apply robustness
to the prescribed transition $\ell(c)\to s'$. If it provides a stuck
execution from $s'$, use the stuck alternative of simulation. Otherwise
it provides matching executions of lengths $k,l$, the cost proof
$(1+k)\,Q\,l$, and a quasi-specimen up to $\preceq_P,\preceq_S$.
Its left endpoint is rooted by the prescribed step followed by the
$k$-step execution; its right endpoint is rooted by its $l$-step
execution. The quasi-up-to lemma therefore supplies the required up-to
relation. This completes every mode and every simulation condition.
\end{proof}

\begin{theorem}[Forward state sufficiency]
Under the hypotheses of this section,
\[
 \Lft(s,t)\ \Longrightarrow\ s\preceq_Q t.
\]
\end{theorem}
\lean{LocalReasoning}{167}{171}{SpecimenSystem.sufficiency}
\begin{proof}
The preceding theorem constructs a counting simulation on $\Lft$.
Apply counting-simulation soundness, Theorem~\ref{thm:soundness}, with
the same deterministic machine and reasonable triple.
\end{proof}
All graph-specific information in this theorem is supplied through the
specimen system and the safety and robustness premises. The theorem's
short proof is therefore not evidence that the missing graph-locality
verification is automatic.

\Needspace{18\baselineskip}
\section{Opposite specimens and the converse argument}

\begin{definition}[Opposite specimen system]
The opposite system $G^{\mathrm{op}}$ has the same specimens, modes,
entering predicate, and rooted-state predicate as $G$. It swaps the two
endpoint maps:
\[
 \ell_{G^{\mathrm{op}}}(c)=r_G(c),\qquad
 r_{G^{\mathrm{op}}}(c)=\ell_G(c).
\]
\end{definition}
\lean{LocalReasoning}{174}{187}{SpecimenSystem.opposite}
\begin{proof}[Verification of its interface laws]
Rootedness preservation is unchanged. Final-state agreement is the
reverse implication pair of the original agreement. For context-only
traversal, exchange the two rootedness hypotheses and the two alternatives
of the availability premise. Apply the original law to get a remote
specimen $d$ and two transitions. Return those transitions in the
opposite order. This explains why the interface law permits availability
on either side, rather than only on the left.
\end{proof}

\begin{lemma}[Lifting commutes with taking the opposite]
For all states $s,t$,
\[
 \mathcal L_{G^{\mathrm{op}}}(s,t)\iff\Lft(t,s).
\]
\end{lemma}
\lean{LocalReasoning}{189}{194}{SpecimenSystem.lifting_opposite}
\begin{proof}
In either direction retain the witnessing specimen and its remoteness.
Exchange the two rootedness proofs and the two endpoint equalities.
The resulting tuple is precisely the witness required by the other
lifting definition.
\end{proof}

\begin{theorem}[Converse construction of a counting simulation]
Suppose $M$ is deterministic, $(Q,P,S)$ is reasonable,
$\safe{Q^{-1}}{P}$ holds, and $G^{\mathrm{op}}$ is robust for
$(Q,P,S)$. Then $\mathcal L_{G^{\mathrm{op}}}$ is a
$(Q,P,S)$-counting simulation.
\end{theorem}
\lean{LocalReasoning}{199}{278}{SpecimenSystem.countingSimulation_converse}
The safety premise is a property of the \emph{original} system at the
converse cost relation. It is not assumed that the opposite system
already satisfies input-safety. This distinction is important when the
right execution in safety~II has length zero.

\begin{proof}
Choose a witness $c$ of $\mathcal L_{G^{\mathrm{op}}}(s,t)$, so
$s=r(c)$, $t=\ell(c)$, $c$ is remote, and both endpoints are rooted.
Finality of $r(c)$ implies finality of $\ell(c)$ by the reverse
direction of final-state agreement. This establishes condition~(A).
For condition~(B), fix a transition $r(c)\to s'$.

\emph{Settled case.} If $c$ is settled, the symmetric availability
premise of context-only traversal applies to this right transition.
Obtain a remote specimen $d$ with transitions
\[
 \ell(c)\to\ell(d),\qquad r(c)\to r(d).
\]
Determinism gives $s'=r(d)$. Propagate rootedness, and obtain
$\mathcal L_{G^{\mathrm{op}}}(r(d),\ell(d))$. Choose zero further
steps from $s'$ and one step from $\ell(c)$. The cost condition is
$1\,Q\,1$, and reflexive refinements give the up-to witness. This
covers backtracking and non-entering search.

\emph{Entering search.} Input-safety for $Q^{-1},P$ supplies three
cases, which must be read in the original endpoint order.
\begin{enumerate}[label=\textup{(\Roman*)}]
\item Use the supplied stuck execution from $r(c)$. The prescribed
  step $r(c)\to s'$ and the stuck-after-step lemma give a stuck
  execution from $s'$. The supplied stuck execution from $\ell(c)$
  is not needed in this direction.
\item There are a settled specimen $d$ and counts $k,l$ with
  \[
    \ell(c)\steps{1+k}\ell(d),\qquad r(c)\steps{l}r(d),\qquad
    l\,Q\,(1+k).
  \]
  Both endpoints of $d$ are rooted. We distinguish the two possible
  forms of $l$.

  If $l=j+1$, remove the prescribed transition from the right execution
  to get $s'\steps{j}r(d)$. Keep the left execution of length $1+k$.
  The opposite lifted pair is $(r(d),\ell(d))$, and its cost condition
  is $(1+j)\,Q\,(1+k)$, exactly the supplied relation with $l=j+1$.
  Reflexive auxiliary refinements finish this subcase.

  If $l=0$, the zero-step equality gives $r(c)=r(d)$. Thus the
  prescribed transition is an outgoing transition from $r(d)$.
  Apply context-only traversal to the \emph{settled} specimen $d$.
  Obtain a remote specimen $e$ and transitions
  \[
    \ell(d)\to\ell(e),\qquad r(d)\to r(e).
  \]
  Determinism identifies $s'$ with $r(e)$. Propagating rootedness
  gives the opposite lifted pair $(r(e),\ell(e))$. There are zero
  further source steps, whereas the target execution is
  \[
    \ell(c)\steps{1+k}\ell(d)\to\ell(e),
  \]
  of length $(1+k)+1$. Its cost obligation is
  $1\,Q\,((1+k)+1)$. The safety relation in this subcase is
  $0\,Q\,(1+k)$. Add $1\,Q\,1$ using additivity of $Q$ to
  obtain the required relation. This additional traversal is the
  essential zero-step correction; the proof never attempts to remove
  a step from a zero-step execution.
\item There are $k,l,a,b$ with
  \[
    \ell(c)\steps{1+k}a,\qquad r(c)\steps{1+l}b,\qquad
    (1+l)\,Q\,(1+k),
  \]
  and a quasi-specimen for auxiliary equivalences. Remove the prescribed
  right transition to get $s'\steps{l}b$. Both $a,b$ are rooted;
  unpacking the quasi-specimen yields a non-rewrite specimen $d$ with
  rooted endpoints and
  \[
    a\simeq_P\ell(d),\qquad r(d)\simeq_P b.
  \]
  This time take the \emph{reverse} halves:
  $b\preceq_P r(d)$ and $\ell(d)\preceq_P a$. Enlarge the
  latter auxiliary order using $P\subseteq S$ to obtain
  $\ell(d)\preceq_S a$. Since
  $\mathcal L_{G^{\mathrm{op}}}(r(d),\ell(d))$, these form
  \[
    \upto{P}{S}{\mathcal L_{G^{\mathrm{op}}}}(b,a).
  \]
  The matching counts are $l,1+k$, with the already supplied cost
  relation $(1+l)\,Q\,(1+k)$.
\end{enumerate}

\emph{Rewrite mode.} Remoteness gives $\neg\ent(c)$. Apply
robustness in $G^{\mathrm{op}}$ to the prescribed transition from
$r(c)$, which is its left endpoint. The stuck alternative is already
the required one. Otherwise the two matching executions and their cost
relation are supplied directly. Rootedness propagates along the
prescribed transition and both executions. Apply the quasi-up-to lemma
\emph{in the opposite system} to obtain the up-to relation on
$\mathcal L_{G^{\mathrm{op}}}$. This completes the proof.
\end{proof}

\Needspace{18\baselineskip}
\section{State and contextual consequences}

The following results compose the two simulation constructions with the
soundness and observation lemmas. They are short in Lean because their
mathematical work has been isolated in those earlier results.

\begin{theorem}[Converse state sufficiency]
If $M$ is deterministic, $(Q,P,S)$ is reasonable,
$\safe{Q^{-1}}{P}$ holds, and $G^{\mathrm{op}}$ is robust for
$(Q,P,S)$, then
\[
 \Lft(s,t)\ \Longrightarrow\ t\preceq_Q s.
\]
\end{theorem}
\lean{LocalReasoning}{281}{286}{SpecimenSystem.sufficiency_converse}
\begin{proof}
The lifting-opposite lemma changes the premise into
$\mathcal L_{G^{\mathrm{op}}}(t,s)$. The converse construction supplies
a counting simulation on this relation. Apply counting-simulation
soundness to obtain the stated refinement.
\end{proof}

\begin{theorem}[Mutual state refinement]
Suppose $M$ is deterministic, $(Q,P,S)$ is reasonable, $G$ is
input-safe for both $(Q,P)$ and $(Q^{-1},P)$, and both $G$ and
$G^{\mathrm{op}}$ are robust for $(Q,P,S)$. Then
\[
 \Lft(s,t)\ \Longrightarrow\ s\simeq_Q t.
\]
\end{theorem}
\lean{LocalReasoning}{289}{295}{SpecimenSystem.equivalence}
\begin{proof}
Forward state sufficiency supplies $s\preceq_Q t$. Converse state
sufficiency supplies $t\preceq_Q s$. These are exactly the two
components of state equivalence.
\end{proof}
The main cost relation $Q$ itself need not be symmetric. The two
directional refinement statements use the same $Q$, with source and
target counts exchanged when the states are exchanged.

For the contextual results, fix the observation interface of the earlier
section and a family $T_i$ of relations on fragments. Assume the
following \emph{initial-lifting property}: for $T_i(a,b)$, every
allowed context $c$ and observations
$\mathsf{Obs}(c,a,s)$, $\mathsf{Obs}(c,b,t)$ give $\Lft(s,t)$.
This property is the explicit premise \texttt{initial\_lift} in the
source. It connects the abstract specimen system with the separately
supplied observation relation.

\begin{theorem}[Forward contextual sufficiency]
If $M$ is deterministic, $(Q,P,S)$ is reasonable,
$\safe{Q}{P}$ and $\robust{Q}{P}{S}$ hold, and the initial-lifting
property holds, then
\[
 T_i(a,b)\ \Longrightarrow\ a\preceq_Q^{\mathcal O}b.
\]
\end{theorem}
\lean{LocalReasoning}{311}{316}{SpecimenSystem.contextual_sufficiency}
\begin{proof}
Use contextual refinement from a lifted relation, taking $L=\Lft$.
Its initial-lifting premise is precisely the one assumed here. Its
state-soundness premise is forward state sufficiency. Applying that
proposition to $T_i(a,b)$ gives the desired contextual refinement.
\end{proof}

\begin{theorem}[Converse contextual sufficiency]
If $M$ is deterministic, $(Q,P,S)$ is reasonable,
$\safe{Q^{-1}}{P}$ holds, $G^{\mathrm{op}}$ is robust for
$(Q,P,S)$, and the same initial-lifting property holds, then
\[
 T_i(a,b)\ \Longrightarrow\ b\preceq_Q^{\mathcal O}a.
\]
\end{theorem}
\lean{LocalReasoning}{318}{323}{SpecimenSystem.contextual_sufficiency_converse}
\begin{proof}
Fix $T_i(a,b)$, an allowed context, an observation of $b$ producing
$s$, and an observation of $a$ producing $t$. Supply the observations
to initial lifting in the original $a,b$ order to get $\Lft(t,s)$.
Converse state sufficiency then gives $s\preceq_Q t$, the required
contextual comparison.
\end{proof}

These are conditional abstract counterparts of the two clauses of the
paper's Theorem~6. Refocusing, output-closure, type-correct plugging, and
the paper's concrete specimen definitions have not disappeared from the
mathematics: their required consequences are currently represented by
the specimen-system laws and initial-lifting premise. Instantiating
those hypotheses is the remaining concrete formalization problem.

\Needspace{18\baselineskip}
\section{Semantic examples and regression proofs}

The examples are small abstract machines used to check the meaning and
edge cases of the definitions. They do not formalize the paper's
lambda-calculus translation, copying UAM, or program-equivalence examples.

\Needspace{17\baselineskip}
\subsection{A countdown machine}

\begin{definition}[Countdown]
The countdown machine has state set $\mathbb N$, transition
$s\to t$ if $s=t+1$, and final predicate $s=0$.
\end{definition}
\lean{Examples}{11}{14}{Examples.countdown}
\begin{proof}[Verification of the machine law]
If $s=0$, an outgoing transition would imply $0=t+1$, impossible for
a natural number $t$. Thus final states are normal as required by
the machine structure.
\end{proof}

\begin{lemma}[Countdown is deterministic]
If $s\to a$ and $s\to b$, then $a=b$.
\end{lemma}
\lean{Examples}{16}{19}{Examples.countdown_deterministic}
\begin{proof}
The premises give $s=a+1$ and $s=b+1$. Cancellation of successor
gives $a=b$.
\end{proof}

\begin{lemma}[Canonical countdown execution]
For every $n\in\mathbb N$, the countdown machine has
$n\steps{n}0$.
\end{lemma}
\lean{Examples}{21}{24}{Examples.countdown_run}
\begin{proof}
Induct on $n$. For $n=0$, use the zero-step constructor. For $n+1$,
take the transition $n+1\to n$ and prepend it to the induction
hypothesis $n\steps{n}0$.
\end{proof}

\begin{proposition}[Countdown characterizes the cost relation]
For every cost preorder $Q$ and states $a,b\in\mathbb N$,
\[
 a\preceq_Q b\quad\Longleftrightarrow\quad a\,Q\,b.
\]
\end{proposition}
\lean{Examples}{27}{36}{Examples.countdown_refines_iff}
\begin{proof}
Suppose $a\preceq_Q b$. Apply it to the canonical successful execution
$a\steps{a}0$. It gives a successful execution of some length $m$
from $b$ and $a\,Q\,m$. Compare this execution with
$b\steps{b}0$. Determinism and final-execution uniqueness give
$m=b$, so $a\,Q\,b$.

Conversely assume $a\,Q\,b$. For any successful execution of length
$n$ from $a$, compare it with $a\steps{a}0$. Uniqueness gives
$n=a$. Choose $b\steps{b}0$ as the target execution; its cost
condition is $n\,Q\,b$, obtained by replacing $n$ with $a$.
\end{proof}

The file then proves three anonymous examples, each by the proposition
and the indicated numerical fact:
\begin{enumerate}
\item $5\preceq_I3$, because $3\leq5$.
\item $3\not\preceq_I5$, because $5\nleq3$.
\item $5\not\preceq_E3$, because $5\ne3$.
\end{enumerate}
\lean{Examples}{38}{39}{Examples: improvement from 5 to 3}
\lean{Examples}{41}{44}{Examples: no improvement from 3 to 5}
\lean{Examples}{46}{49}{Examples: unequal exact counts}
These verify that improvement uses the target count as the smaller
quantity. Reversing the intended cost orientation would invalidate the
first two examples.

\Needspace{17\baselineskip}
\subsection{A crash and an infinite loop}

\begin{definition}[Fault states and the fault machine]
There are four distinct states:
\[
 \mathsf{done},\quad\mathsf{crash},\quad\mathsf{stuck},\quad\mathsf{loop}.
\]
The only transitions are
\[
 \mathsf{crash}\to\mathsf{stuck},\qquad
 \mathsf{loop}\to\mathsf{loop}.
\]
Only $\mathsf{done}$ is final.
\end{definition}
\lean{Examples}{51}{53}{Examples.FaultState}
\lean{Examples}{55}{60}{Examples.fault}
The four names are constructors of a finite datatype with derived
decidable equality. Here the constructor $\mathsf{stuck}$ names a
particular state; $\stuck(s)$ remains the general predicate defined
earlier.
\begin{proof}[Verification of the machine law]
No transition has source $\mathsf{done}$, because both possible source
constructors are distinct from it. Consequently a final state has no
outgoing transition.
\end{proof}

\begin{lemma}[The fault machine is deterministic]
Any two transitions with the same source have the same target.
\end{lemma}
\lean{Examples}{62}{65}{Examples.fault_deterministic}
\begin{proof}
Case-analyze the two transition proofs. If both use the crash rule,
both targets are $\mathsf{stuck}$. If both use the loop rule, both
targets are $\mathsf{loop}$. Either mixed case would equate
$\mathsf{crash}$ with $\mathsf{loop}$ and is impossible.
\end{proof}

\begin{definition}[The crash--loop candidate]
Let $R(s,t)$ mean $s=\mathsf{crash}$ and $t=\mathsf{loop}$.
\end{definition}
\lean{Examples}{67}{67}{Examples.crashLoop}

\begin{lemma}[A simulation using the stuck alternative]
The crash--loop candidate is an $(E,E,E)$-counting simulation.
\end{lemma}
\lean{Examples}{71}{88}{Examples.crash_loop_simulation}
\begin{proof}
For its only related pair, the source is $\mathsf{crash}$, which is
not final; hence the finality condition is vacuous. Its only successor
is $\mathsf{stuck}$. Use the zero-step execution from that successor
to itself in the stuck alternative of simulation. It is normal because
neither transition rule starts there, and it is not final because it
differs from $\mathsf{done}$.
\end{proof}

The next anonymous example concludes
$\mathsf{crash}\preceq_E\mathsf{loop}$.
\lean{Examples}{90}{91}{Examples: crash refines loop}
\begin{proof}
Apply counting-simulation soundness to the preceding simulation, fault
determinism, the reasonable triple $(E,E,E)$, and the only pair of $R$.
\end{proof}
There is no requirement to match a successful source execution, because
none exists. Thus the result is consistent even with exact costs. The
example exercises the stuck branch of the implemented proof; it does
not prove that every conceivable up-to simulation proof of this same
refinement must use that branch.

\Needspace{23\baselineskip}
\Needspace{17\baselineskip}
\subsection{The zero-step converse case}

\begin{definition}[Visit specimens and endpoint maps]
The type of specimens has three distinct constructors
$\mathsf{entering},\mathsf{settled},\mathsf{finished}$, with derived
decidable equality. The endpoint maps into the countdown machine are
\[
 \begin{aligned}
 \ell(\mathsf{entering})&=2,& r(\mathsf{entering})&=1,\\
 \ell(\mathsf{settled})&=1,& r(\mathsf{settled})&=1,\\
 \ell(\mathsf{finished})&=0,& r(\mathsf{finished})&=0.
 \end{aligned}
\]
\end{definition}
\lean{Examples}{93}{95}{Examples.Visit}
\lean{Examples}{97}{100}{Examples.visitLeft}
\lean{Examples}{102}{105}{Examples.visitRight}

\begin{definition}[Visit specimen system]
Every visit specimen has search mode. The entering predicate holds
exactly of the constructor $\mathsf{entering}$. Every state is declared
rooted. Together with the endpoint maps above, these data define a
specimen system over countdown.
\end{definition}
\lean{Examples}{107}{126}{Examples.visitSystem}
\begin{proof}[Verification of the interface laws]
Rootedness preservation is immediate since rootedness is the true
predicate. Final-state agreement holds because the endpoint pairs are
$(2,1)$, $(1,1)$, and $(0,0)$: only in the third pair is either endpoint
zero, and then both are.

For context-only traversal, consider the specimen. The constructor
$\mathsf{entering}$ cannot satisfy the settled premise: its mode is
search and its entering predicate holds. At the constructor
$\mathsf{settled}$, both endpoints are $1$. Choose
$\mathsf{finished}$ and the countdown transitions $1\to0$ on both
sides; its search mode is remote. At $\mathsf{finished}$, both
endpoints are zero and neither has an outgoing transition, contradicting
the availability premise. These cases exhaust the datatype.
\end{proof}
This is a deliberately abstract instance. For example, all its modes are
search, including that of the pair of final states. The specimen interface
does not enforce a concrete UAM convention for final focus modes.

\begin{lemma}[Visit input-safety]
The visit system is input-safe for $(H^{-1},E)$.
\end{lemma}
\lean{Examples}{129}{137}{Examples.visit_safe}
\begin{proof}
The entering premise forces the input specimen to be
$\mathsf{entering}$. Use safety~II with output specimen
$\mathsf{settled}$ and $k=l=0$. The left execution is the single
countdown transition $2\to1$; the right execution has zero steps at
$1$. The output specimen is searching without entering and is therefore
settled. The cost requirement is $1\,H^{-1}\,0$, or $0\leq1$.
\end{proof}

\begin{lemma}[Visit converse robustness]
The opposite visit system is robust for $(H,E,E)$.
\end{lemma}
\lean{Examples}{139}{142}{Examples.visit_converse_robust}
\begin{proof}
Its mode is always search. The rewrite-mode premise of robustness is
therefore impossible, so every required implication holds vacuously.
\end{proof}

\begin{proposition}[The converse zero-step regression]
In countdown, $1\preceq_H2$.
\end{proposition}
\lean{Examples}{146}{149}{Examples.converse_zero_step_case}
\begin{proof}
Use converse state sufficiency for the visit system. Countdown is
deterministic, $(H,E,E)$ is reasonable, and the preceding two lemmas
supply its safety and converse robustness premises. The entering
specimen is remote, both its endpoints are rooted, and it witnesses
$\Lft(2,1)$. The theorem consequently gives $1\preceq_H2$.

In the underlying converse simulation construction, safety~II provides
the pair $(1,1)$ after one original left step and zero original right
steps. The prescribed right step is $1\to0$. Context traversal at
$(1,1)$ supplies $(0,0)$ and adds one transition to the original left
execution, now $2\to1\to0$. Thus the source terminates in one step
and the target in two, satisfying $1\,H\,2$. This is precisely the
zero-step subcase proved in the converse argument.
\end{proof}

\Needspace{18\baselineskip}
\section{Finite carriers and concrete interfaced graphs}

We now turn to actual graph objects. The constructions in the following sections
prove their own incidence and typing obligations. They do not assume the specimen
laws used in the preceding abstract argument. We work with small finite vertex
and edge carriers; operation labels may live in a larger universe.

\Needspace{24\baselineskip}
\subsection{Enumerations and bijections}
\begin{definition}[Finite enumeration]\label{term:enumeration}
An enumeration of a type $A$ consists of a finite list $l$ and a proof that every
$a\in A$ belongs to $l$. Repetitions are allowed. This witnesses finiteness without
requiring decidable equality or a canonical choice of names.
\end{definition}
\lean{Graph/Finite}{9}{11}{Enumeration}
The empty type is enumerated by the empty list: completeness follows because
there is no element to check. The finite ordinal $\{0,\ldots,n-1\}$ is enumerated
in increasing order by Lean's finite range; its library membership theorem proves
completeness, including when $n=0$.
\lean{Graph/Finite}{15}{15}{Enumeration.empty}
\lean{Graph/Finite}{17}{17}{Enumeration.fin}
For a disjoint union $A+B$, concatenate the tagged enumerations of $A$ and $B$.
An element has exactly one tag, and completeness of the corresponding list places
it in the appropriate half of the concatenation.
\lean{Graph/Finite}{19}{25}{Enumeration.sum}
For a subset $\{a\in A\mid p(a)\}$, attach membership certificates to the enumeration
of $A$ and retain just entries satisfying $p$. Classical decidability of $p$
allows this filtering. Every subset element is retained because its certificate
proves the test true. The resulting definition is noncomputable for an arbitrary
proposition $p$; no algorithm deciding $p$ is claimed.
\lean{Graph/Finite}{27}{37}{Enumeration.subtype}
For a dependent sum $\sum_{a\in A}B_a$, enumerate each fibre, tag it with $a$, and
concatenate over the enumeration of $A$. A pair $(a,b)$ appears because $a$ occurs
in the outer list and $b$ in the enumeration of its fibre.
\lean{Graph/Finite}{39}{43}{Enumeration.sigma}

\begin{definition}[Bijection]\label{term:bijection}
A bijection $f:A\cong B$ consists of functions $f:A\to B$ and $f^{-1}:B\to A$
with both inverse equations, pointwise.
\end{definition}
\lean{Graph/Finite}{47}{51}{Bijection}
Identity uses the identity function twice. Inversion interchanges the two functions
and their inverse laws. Composition uses $g\circ f$ forwards and
$f^{-1}\circ g^{-1}$ backwards; cancelling the inner inverse pair and then the
outer pair proves each law.
\lean{Graph/Finite}{57}{57}{Bijection.refl}
\lean{Graph/Finite}{59}{60}{Bijection.symm}
\lean{Graph/Finite}{62}{66}{Bijection.trans}
Bijections extend to disjoint unions by applying the appropriate map inside each
tag. Case analysis on the tag reduces both inverse laws to the component laws.
\lean{Graph/Finite}{68}{73}{Bijection.sum}
Every bijection is injective: apply its inverse to an equality of images and
cancel. An enumeration transports along a bijection by mapping its forward
function over the list. For $b\in B$, the listed preimage $f^{-1}(b)$ maps to $b$,
proving completeness.
\lean{Graph/Finite}{75}{77}{Bijection.injective}
\lean{Graph/Finite}{79}{81}{Bijection.enumeration}
\begin{lemma}[Distinctness under restricted injectivity]
If a list $l$ has no repeated elements and $f$ is injective on elements of $l$,
then $\operatorname{map}(f,l)$ has no repeated elements.
\end{lemma}
\begin{proof}
Distinct positions in $l$ contain distinct elements. If their images were equal,
restricted injectivity would identify those elements, a contradiction. Lean
expresses this as preservation of the pairwise inequality predicate under mapping.
\end{proof}
\lean{Graph/Finite}{86}{92}{nodup_map_on}
This restricted version will be essential: replacement can identify an input
with an output, while remaining injective on each incidence list separately.

\Needspace{24\baselineskip}
\subsection{Ports, signatures, and graph invariants}
\begin{definition}[Signature and ports]\label{term:ports}
A signature $\sigma$ assigns an ordered input-label list and output-label list
to every operation. With internal carrier $V$ and boundary lengths $n,m$, vertices
are the disjoint union
\[
 \operatorname{Port}(V,n,m)=V+\operatorname{Fin}(n)+\operatorname{Fin}(m).
\]
The three constructors are internal, input, and output. In particular the two
boundary sets are disjoint even if their wire labels agree.
\end{definition}
\lean{Graph/Basic}{13}{15}{Signature}
\lean{Graph/Basic}{17}{21}{Port}
The predicates $\mathsf{IsInput}$ and $\mathsf{IsOutput}$ test the corresponding
constructor. A carrier map $f:V\to W$ acts only on internal vertices and fixes
each boundary position. It is injective whenever $f$ is: unlike constructors
cannot have equal images, and equality within the internal constructor reduces
to injectivity of $f$. The input and output predicates are preserved and reflected
by this map, by the same three-case analysis.
\lean{Graph/Basic}{27}{29}{Port.IsInput}
\lean{Graph/Basic}{31}{33}{Port.IsOutput}
\lean{Graph/Basic}{35}{38}{Port.map}
\lean{Graph/Basic}{40}{43}{Port.map_injective}
\lean{Graph/Basic}{45}{46}{Port.map_isInput}
\lean{Graph/Basic}{48}{49}{Port.map_isOutput}
\termdef{vertex-label}
Given boundary-label lists $I,O$ and internal labelling $\lambda:V\to L$, the label
of an internal vertex is $\lambda(v)$, of input $i$ is $I_i$, and of output $j$
is $O_j$. The finite-index types ensure that these lookups are defined.
\lean{Graph/Basic}{51}{54}{Port.label}

\begin{definition}[Concrete typed net]\label{term:net}
A net $G:\!I\to O$ has finite internal and edge carriers, labels for both, and
ordered source and target vertex lists $s(e),t(e)$ for each edge. It satisfies:
\begin{enumerate}
\item Each individual source list and target list has no repetitions.
\item A vertex belongs to the sources of at most one edge, and to the targets
of at most one edge.
\item A vertex belongs to some source list exactly when it is not an output;
it belongs to some target list exactly when it is not an input.
\item Mapping vertex labels over $s(e)$ and $t(e)$ gives precisely the input
and output lists of the signature of the edge label.
\end{enumerate}
\end{definition}
\lean{Graph/Basic}{61}{77}{Net}
These fields are part of a net's certificate of well-formedness. Coverage excludes
isolated vertices. It does not exclude cycles or a vertex belonging to both the
source and target of the same edge. The full vertex type and its label function
are abbreviations for the port type and the labelling just defined.
\lean{Graph/Basic}{84}{84}{Net.Vertex}
\lean{Graph/Basic}{86}{86}{Net.vertexLabel}
From forward source coverage, membership in $s(e)$ implies that a vertex is not
an output. From target coverage, membership in $t(e)$ implies that it is not an
input. Apply the appropriate equivalence to the witness edge $e$.
\lean{Graph/Basic}{88}{89}{Net.source_not_output}
\lean{Graph/Basic}{91}{92}{Net.target_not_input}
Taking lengths in the two typing equations shows that incidence-list lengths are
exactly the corresponding signature arities: mapping labels does not change length.
\lean{Graph/Basic}{94}{96}{Net.source_length}
\lean{Graph/Basic}{98}{100}{Net.target_length}
\termdef{indexed-port}
The $i$th source port of $e$ is its source list indexed by a position in the
signature's input list; the preceding length equality justifies the lookup.
Target ports are defined in the same way with the output list.
\lean{Graph/Basic}{102}{103}{Net.sourcePort}
\lean{Graph/Basic}{105}{106}{Net.targetPort}
Each indexed port belongs to its incidence list, by the list lookup membership
lemma. Within either list, equality of indexed ports implies equality of indices,
because that list has no duplicates; equality of finite ordinals follows from
equality of their underlying natural numbers.
\lean{Graph/Basic}{108}{108}{Net.sourcePort_mem}
\lean{Graph/Basic}{110}{110}{Net.targetPort_mem}
\lean{Graph/Basic}{112}{114}{Net.sourcePort_injective}
\lean{Graph/Basic}{116}{118}{Net.targetPort_injective}
The label of source port $i$ is the $i$th signature input label, and the label of
target port $i$ is the $i$th signature output label. Apply indexed lookup to the
appropriate list typing equation, then commute lookup with mapping.
\lean{Graph/Basic}{120}{123}{Net.sourcePort_label}
\lean{Graph/Basic}{125}{128}{Net.targetPort_label}
Finally, membership in a source list is equivalent to being its indexed port for
some signature input position. Given membership, choose its list index and
transport the bound along the arity equality. Conversely, indexed ports belong
to their lists. The target statement uses the output arity and the identical
membership argument.
\lean{Graph/Basic}{130}{137}{Net.mem_source_iff_port}
\lean{Graph/Basic}{139}{146}{Net.mem_target_iff_port}

\Needspace{24\baselineskip}
\subsection{Elementary graph constructors}
The ordered finite range has no duplicates: equality of its entries at two
positions gives equality of their underlying natural numbers and hence of the
positions.
\lean{Graph/Constructors}{5}{9}{finRange_nodup}
\begin{definition}[Corolla and empty graph]\label{term:corolla}
The corolla of operation $o$ has one edge labelled $o$, no internal vertices,
and all of its prescribed inputs and outputs in signature order. Its sources
are precisely the input constructors and its targets the output constructors.
\end{definition}
\begin{proof}[Well-formedness]
Finite ranges and their constructor maps give distinct incidence lists. There
is only one edge, so edge uniqueness is immediate. Every input occurs in its
source range and every output in its target range, and the opposite boundary
constructors cannot occur there. There are no internal vertices to check.
Mapping labels over either ordered range recovers the corresponding signature
list, proving typing.
\end{proof}
\lean{Graph/Constructors}{18}{52}{Net.corolla}
\termdef{empty-graph}
The empty graph has empty internal and edge carriers and empty boundaries. Each
incidence or label function has empty domain; all obligations follow by eliminating
the purported element of an empty type. A bare identity wire is not provided:
a single vertex serving as both boundary input and boundary output would violate
the disjoint-boundary representation used here.
\lean{Graph/Constructors}{56}{78}{Net.empty}

\Needspace{18\baselineskip}
\section{Shallow plugging, including self-loop holes}
Fix a host $C:I\to O$, an edge $h$ to remove, and a replacement $H$ whose boundary
lists are exactly the input and output signature lists of $C$'s label at $h$.
All constructions in this section are in the namespace $\mathsf{Net.Replacement}$
unless explicitly stated otherwise.

\Needspace{24\baselineskip}
\subsection{Carriers and boundary identification}
\begin{definition}[Replacement carriers]
The new internal vertices are $C_V+H_V$, with disjoint tags providing freshness.
The new edges are $(C_E\setminus\{h\})+H_E$. The outer boundary remains $I,O$.
The surviving-edge subtype carries a proof of inequality to $h$, so the removed
edge cannot accidentally survive as an old edge.
\end{definition}
\lean{Graph/Plugging}{19}{19}{Net.Replacement.Vertex}
\lean{Graph/Plugging}{20}{20}{Net.Replacement.Edge}
The host embedding $j$ tags each old internal vertex on the left and fixes the
outer boundary. The insertion map $k$ tags a replacement internal vertex on the
right, sends replacement input $i$ to $j(s_h(i))$, and sends replacement output
$i$ to $j(t_h(i))$. Thus gluing is expressed by a function rather than a quotient
of all vertex names.
\lean{Graph/Plugging}{22}{22}{Net.Replacement.host}
\lean{Graph/Plugging}{24}{27}{Net.Replacement.insert}
The host embedding is injective by injectivity of the left sum tag and the port
map lemma. Its image is disjoint from freshly tagged replacement internals:
case analysis on a host vertex gives either a boundary constructor or an internal
left tag, neither equal to an internal right tag.
\lean{Graph/Plugging}{29}{30}{Net.Replacement.host_injective}
\lean{Graph/Plugging}{32}{33}{Net.Replacement.host_ne_internal}

\begin{lemma}[The two restricted insertion injections]
The insertion map is injective on vertices that are not outputs. It is also
injective on vertices that are not inputs.
\end{lemma}
\begin{proof}
For the first claim, exclude output cases using the hypotheses. Two internal
images agree only when their right-tagged names agree. An internal image cannot
agree with an input image, by freshness. Two input images agree only when their
host source ports agree, and source-port injectivity identifies their indices.
For the second claim, exclude input cases instead, use target-port injectivity
on two outputs, and use the same freshness argument on mixed cases.
\end{proof}
\lean{Graph/Plugging}{35}{50}{Net.Replacement.insert_source_injective}
\lean{Graph/Plugging}{52}{67}{Net.Replacement.insert_target_injective}
If a non-output replacement vertex has insertion image $j(v)$, it must be some
input $i$, with $s_h(i)=v$. Outputs are excluded by hypothesis, internals by
freshness, and in the input case host injectivity gives the equation. Dually, a
non-input replacement vertex mapping to $j(v)$ must be output $i$ with
$t_h(i)=v$. These inversion lemmas control mixed old/new incidence cases.
\lean{Graph/Plugging}{69}{75}{Net.Replacement.insert_source_host}
\lean{Graph/Plugging}{77}{83}{Net.Replacement.insert_target_host}

\Needspace{24\baselineskip}
\subsection{Preservation of incidence invariants}
Sources and targets of surviving host edges are their old lists mapped by $j$.
Those of inserted edges are the replacement lists mapped by $k$.
\lean{Graph/Plugging}{85}{87}{Net.Replacement.source}
\lean{Graph/Plugging}{89}{91}{Net.Replacement.target}
\begin{lemma}[Distinct incidence entries]
Every new source list and every new target list has no duplicates.
\end{lemma}
\begin{proof}
A surviving host list uses the injective map $j$. A replacement source list
contains no outputs, so the first restricted injection of $k$ applies to its
members. A replacement target list contains no inputs, so the second injection
applies. In all cases use preservation of distinctness under restricted injectivity.
\end{proof}
\lean{Graph/Plugging}{93}{98}{Net.Replacement.source_nodup}
\lean{Graph/Plugging}{100}{105}{Net.Replacement.target_nodup}

\begin{lemma}[Unique source and target incidence]
A vertex belongs to the sources of at most one new edge, and to the targets
of at most one new edge.
\end{lemma}
\begin{proof}
For source incidence, split on the two edge tags. Two host edges reduce, by
injectivity of $j$, to old source uniqueness. Two inserted edges reduce, by
restricted source injectivity of $k$, to source uniqueness in $H$. In the mixed
case the shared vertex is an inserted input, by the source inversion lemma;
its corresponding host vertex is therefore a source port of $h$. Old source
uniqueness identifies the surviving host edge with $h$, contradicting its subtype
certificate. The other mixed ordering is the same contradiction with the equality
reversed. For target incidence, repeat the four cases with target uniqueness,
restricted target injectivity, and the target inversion lemma: the mixed vertex
is now an inserted output and a target port of $h$.
\end{proof}
\lean{Graph/Plugging}{107}{136}{Net.Replacement.source_unique}
\lean{Graph/Plugging}{138}{167}{Net.Replacement.target_unique}
A vertex in a new source list is not an outer output. In the host case use the
old source property and preservation of boundary constructors. In the inserted
case an internal vertex remains internal, an input maps to a source port of $h$,
and an output is impossible in a replacement source list. For a new target list,
exchange source/input with target/output: inserted inputs are impossible and
inserted outputs map to host target ports, which cannot be outer inputs.
\lean{Graph/Plugging}{169}{180}{Net.Replacement.source_not_output}
\lean{Graph/Plugging}{182}{193}{Net.Replacement.target_not_input}

\begin{lemma}[Coverage of surviving host vertices]
For every host vertex $v$ that is not an output, $j(v)$ belongs to a new source
list. If $v$ is not an input, $j(v)$ belongs to a new target list.
\end{lemma}
\begin{proof}
Choose an old source edge by coverage in $C$. If it differs from $h$, retain it.
If it equals $h$, express $v$ as source port $i$; replacement input $i$ belongs
to a source edge by coverage in $H$, and its image is $j(v)$. For target coverage,
choose an old target edge. Retain it unless it is $h$; in that case express $v$
as target port $i$ and use coverage at replacement output $i$.
\end{proof}
\lean{Graph/Plugging}{195}{204}{Net.Replacement.source_host_covered}
\lean{Graph/Plugging}{206}{215}{Net.Replacement.target_host_covered}
Full source coverage now follows by splitting a new vertex into a host internal,
a fresh internal, an outer input, or an outer output. Host internals and inputs
use the preceding lemma. Fresh internals use replacement coverage and map the
witness edge into the right edge tag. Outer outputs contradict the required
non-output hypothesis. The reverse implication is the source exclusion already
proved. Full target coverage uses host internals and outer outputs, replacement
coverage for fresh internals, and excludes outer inputs; its reverse implication
is target exclusion.
\lean{Graph/Plugging}{217}{231}{Net.Replacement.source_coverage}
\lean{Graph/Plugging}{233}{247}{Net.Replacement.target_coverage}

\Needspace{24\baselineskip}
\subsection{Labels, typing, and the certified result}
New internal labels use the host label on left tags and replacement label on
right tags. Surviving edges keep their labels; inserted edges receive their
replacement labels.
\lean{Graph/Plugging}{249}{249}{Net.Replacement.internalLabel}
\lean{Graph/Plugging}{251}{253}{Net.Replacement.edgeLabel}
Labelling commutes with $j$, by the three port cases. It also commutes with $k$:
internal labels agree by definition, while input and output labels agree by the
host indexed-port typing lemmas and the matching replacement interface.
\lean{Graph/Plugging}{255}{257}{Net.Replacement.host_label}
\lean{Graph/Plugging}{259}{264}{Net.Replacement.insert_label}
For source typing, split the edge tag and compose the incidence map with the
vertex-label map. The commuting equation just proved reduces the result to
source typing in $C$ or $H$, respectively. For target typing the same composition
argument reduces to their target typing equations.
\lean{Graph/Plugging}{266}{276}{Net.Replacement.source_typed}
\lean{Graph/Plugging}{278}{288}{Net.Replacement.target_typed}

\begin{theorem}[Well-formed shallow plugging]
There is a net $C[h:=H]:I\to O$ with the carriers, labels, and incidence lists
above. It preserves the entire ordered external interface and all graph invariants,
even when $h$ has a vertex in both its source and target lists.
\end{theorem}
\begin{proof}
Enumerate the new internal carrier by the sum enumeration. Enumerate surviving
host edges by subset filtering, and concatenate the tagged replacement-edge
enumeration. Supply each of the distinctness, uniqueness, coverage, and typing
proofs established above as the corresponding net field. This assembles the
certified graph. Classical filtering of the surviving-edge subtype accounts for
the noncomputable qualifier.
\end{proof}
\lean{Graph/Plugging}{294}{312}{Net.plug}
There is deliberately no claim that $k$ is globally injective in this theorem.
If a host source port equals a host target port, its replacement input and output
are identified. The separate source and target injection arguments are sufficient
for every net invariant.

\addtocontents{toc}{\protect\newpage}
\Needspace{18\baselineskip}
\section{Paths, renaming, and isomorphism}
\Needspace{24\baselineskip}
\subsection{Paths and their behaviour under plugging}
\begin{definition}[Adjacency and path]\label{term:path}
Edges $e,f$ are adjacent when some vertex belongs to both $t(e)$ and $s(f)$.
A path is a nonempty edge list: a singleton is a path; prepending $e$ to a path
starting at $f$ gives a path when $e,f$ are adjacent. Edges may repeat, and this
notion does not specify a chosen boundary vertex at either end.
\end{definition}
\lean{Graph/Paths}{11}{12}{Net.Adjacent}
\lean{Graph/Paths}{15}{17}{Net.Path}
For a list beginning $e,f$, the path predicate is equivalent to adjacency of
$e,f$ together with the path predicate for the tail starting at $f$. Invert the
path constructor for the forward implication and apply it for the reverse one.
\lean{Graph/Paths}{19}{23}{Net.path_pair_iff}
Any map on edges preserving adjacency maps paths to paths. Induct on the path:
singletons remain singletons, and the inductive step transports its first
adjacency and then its tail path. If the edge map preserves and reflects
adjacency, it preserves and reflects paths. Induct on the list: the empty list
is a path on neither side, singletons are paths on both sides, and longer lists
reduce to the first adjacency equivalence and the tail induction hypothesis.
\lean{Graph/Paths}{25}{31}{Net.Path.map}
\lean{Graph/Paths}{33}{42}{Net.path_map_iff}

\begin{lemma}[Host and inserted adjacencies]
Two surviving host edges are adjacent after plugging exactly when they were
adjacent before plugging. Two inserted edges that were adjacent remain adjacent.
\end{lemma}
\begin{proof}
For host reflection, pull both incidence witnesses back through their mapped
lists. Equality of their host images implies equality of the original vertices
by injectivity of $j$, giving an old shared witness. For host preservation map
an old shared witness through $j$. For inserted preservation map the shared
replacement witness through $k$; injectivity is unnecessary in this direction.
\end{proof}
\lean{Graph/Paths}{51}{61}{Net.Replacement.adjacent_host_iff}
\lean{Graph/Paths}{63}{66}{Net.Replacement.adjacent_insert}
\begin{definition}[Non-self-loop hole]\label{term:non-self-loop}
The hole $h$ is non-self-looping if no vertex belongs to both $s(h)$ and $t(h)$.
This requires the two lists to be disjoint; it is stronger than distinctness
within each individual list, already required of every net.
\end{definition}
\lean{Graph/Paths}{70}{70}{Net.Replacement.NonSelfLoop}
For such a hole, $k$ is globally injective. Equal internal images reduce to equal
names; mixed internal/boundary cases contradict freshness. Two input images or
two output images reduce to the appropriate port injection. Equality of an
input image and an output image would exhibit a host vertex in both $s(h)$ and
$t(h)$, contrary to non-self-looping.
\lean{Graph/Paths}{72}{93}{Net.Replacement.insert_injective}
Consequently inserted adjacency is reflected as well as preserved: pull the two
mapped-list witnesses back and use global injectivity of $k$ to identify them.
The other implication is the unconditional preservation result above.
\lean{Graph/Paths}{95}{104}{Net.Replacement.adjacent_insert_iff}

\begin{theorem}[Path preservation and reflection]
A list of surviving host edges is a path after plugging if and only if its old
edge list is a path. Any replacement path gives an inserted path. For a
non-self-loop hole the converse holds for inserted paths too.
\end{theorem}
\begin{proof}
For host paths, induct on the surviving-edge list. Empty and singleton lists
are immediate; a longer list reduces to host adjacency equivalence and the tail
hypothesis. For inserted preservation, apply path mapping to inserted adjacency
preservation. For inserted reflection under the stated restriction, apply the
general path equivalence lemma to inserted adjacency equivalence.
\end{proof}
\lean{Graph/Paths}{109}{120}{Net.Replacement.path_host_iff}
\lean{Graph/Paths}{122}{125}{Net.Replacement.path_insert}
\lean{Graph/Paths}{129}{132}{Net.Replacement.path_insert_iff}

\Needspace{24\baselineskip}
\subsection{Transport along carrier bijections}
A bijection on internal names extends to ports by mapping only internal vertices.
Its inverse laws follow by cases on a port: internal cases use the original
inverse laws and boundary cases are reflexive. Mapping the identity is pointwise
the identity, also by port cases; function extensionality gives the corresponding
identity of functions.
\lean{Graph/Isomorphism}{9}{14}{Port.bijection}
\lean{Graph/Isomorphism}{16}{17}{Port.map_id}
\lean{Graph/Isomorphism}{19}{20}{Port.map_id_fun}
Fix bijections $a:G_V\cong V'$ and $b:G_E\cong E'$. Reindexed sources and targets
at $e'\in E'$ are the old incidence lists at $b^{-1}(e')$, mapped by the port
extension of $a$.
\lean{Graph/Isomorphism}{35}{36}{Net.Reindex.source}
\lean{Graph/Isomorphism}{38}{39}{Net.Reindex.target}
Source uniqueness survives this transport: two shared mapped witnesses are equal
by port injectivity, so old source uniqueness identifies their old edge indices;
apply $b$ and cancel $b\circ b^{-1}$. The target uniqueness proof uses old target
uniqueness and the same cancellation.
\lean{Graph/Isomorphism}{41}{48}{Net.Reindex.source_unique}
\lean{Graph/Isomorphism}{50}{57}{Net.Reindex.target_unique}
For source coverage, a new source witness pulls back to an old source vertex,
which is not an output; port maps preserve this property. Conversely, pull a
non-output vertex backwards through $a^{-1}$, choose an old source edge by
coverage, and map both the edge and vertex forwards. The two inverse laws prove
that this is a witness for the original new vertex. Target coverage repeats this
argument with non-input vertices and old target coverage.
\lean{Graph/Isomorphism}{59}{70}{Net.Reindex.source_coverage}
\lean{Graph/Isomorphism}{72}{83}{Net.Reindex.target_coverage}
Give an internal name $w$ the label of $a^{-1}(w)$. Labelling an old vertex after
mapping by $a$ then agrees with its old label: the internal case cancels
$a^{-1}\circ a$, and boundary labels are fixed.
\lean{Graph/Isomorphism}{85}{88}{Net.Reindex.label_map}
\begin{definition}[Reindexed graph]\label{term:reindexing}
The reindexed net uses the new carriers, inverse-transported labels, and the
incidence lists just defined.
\end{definition}
\begin{proof}[Well-formedness]
Transport finite enumerations along the bijections. Distinctness follows from
port injectivity, uniqueness and coverage from the preceding lemmas. For typing,
compose incidence mapping with the vertex-label map, use label commutation, and
apply the original typing equation. The boundary lists are unchanged.
\end{proof}
\lean{Graph/Isomorphism}{94}{121}{Net.reindex}

\Needspace{24\baselineskip}
\subsection{Graph isomorphism and shallow context congruence}
\begin{definition}[Isomorphism]\label{term:isomorphism}
An isomorphism $G\cong H$ between nets with the same signature and boundary lists
is a bijection of internal vertices and a bijection of edges preserving internal
labels, edge labels, and both ordered incidence lists. Its port action fixes
boundary positions.
\end{definition}
\lean{Graph/Isomorphism}{123}{129}{Net.Iso}
Identity bijections and identity list maps yield a reflexive isomorphism. For the
inverse isomorphism, invert both bijections. Inverse label preservation follows
by applying the original equation at an inverse image and cancelling. To invert
an incidence equation, map its two sides by the inverse port map and cancel the
forward/inverse maps pointwise. For composition, compose carrier bijections;
compose label equalities, and use successive list mapping followed by the two
incidence-preservation equations.
\lean{Graph/Isomorphism}{135}{141}{Net.Iso.refl}
\lean{Graph/Isomorphism}{143}{157}{Net.Iso.symm}
\lean{Graph/Isomorphism}{159}{175}{Net.Iso.trans}
Isomorphisms preserve and reflect adjacency: transport a shared witness forwards
for preservation; for reflection pull both witnesses back from the incidence
maps and identify their preimages by port injectivity. Applying the general path
equivalence lemma then preserves and reflects paths, with their edge lists
mapped by the edge bijection.
\lean{Graph/Isomorphism}{177}{191}{Net.Iso.adjacent_iff}
\lean{Graph/Isomorphism}{193}{195}{Net.Iso.path_iff}
The original graph is isomorphic to its reindexing. Use the chosen carrier
bijections; inverse cancellation proves label preservation and each incidence
identity directly from the definition of reindexing.
\lean{Graph/Isomorphism}{199}{207}{Net.reindexIso}

Suppose $f:H\cong K$ and both graphs have the interface of host hole $h$.
The intended isomorphism between $C[h:=H]$ and $C[h:=K]$ is the identity on
host internal names and $f$ on replacement internal names, and likewise the
identity on surviving host edges and $f$ on inserted edges. Sum bijections
supply their inverse laws.
\lean{Graph/PlugIsomorphism}{16}{17}{Net.PlugIso.vertices}
\lean{Graph/PlugIsomorphism}{19}{20}{Net.PlugIso.edges}
The resulting port map commutes with host embedding: check the three port
constructors, all of which are reflexive. It commutes with insertion too:
internal vertices use the replacement part of the sum map, and boundary vertices
reduce to the host-embedding equation at the corresponding hole port.
\lean{Graph/PlugIsomorphism}{22}{24}{Net.PlugIso.map_host}
\lean{Graph/PlugIsomorphism}{26}{32}{Net.PlugIso.map_insert}
\begin{theorem}[Congruence of shallow plugging]
For every typed host and hole, $H\cong K$ implies $C[h:=H]\cong C[h:=K]$.
No non-self-loop hypothesis is required.
\end{theorem}
\begin{proof}
Use the sum bijections above. For each label equation split the carrier tag:
host labels are unchanged and inserted labels use $f$'s corresponding equation.
For a host source list, commute the new port map past host embedding pointwise.
For an inserted source list, commute it past insertion and then use source
preservation by $f$. Repeat with target lists and target preservation. These
arguments prove all isomorphism fields without reflecting through insertion.
\end{proof}
\lean{Graph/PlugIsomorphism}{38}{85}{Net.plugIso}
Here edge labels are preserved by literal equality. For boxed labels containing
bodies, this is not yet the recursive notion that would identify independently
renamed but isomorphic bodies. No such stronger congruence is claimed.

\Needspace{18\baselineskip}
\section{Concrete uniqueness of operation paths}
Choose a predicate $P$ on operation labels in a net $G$. This section first proves
a generic graph theorem; the next section shows that the single-source condition
is automatic for the typed operation labels of the linear hierarchy.

\begin{definition}[Unary operations and traversal]\label{term:unary}
Selected operations are unary when each edge whose label satisfies $P$ has a
source list of length one. An operation step $s\rightsquigarrow t$ chooses such
an edge $e$ with $s\in s(e)$ and $t\in t(e)$.
\end{definition}
\lean{Graph/OperationPaths}{17}{18}{Net.UnaryOperations}
\lean{Graph/OperationPaths}{20}{21}{Net.OperationStep}
\termdef{reverse-traversal}
Reverse traversal is the machine whose transition from $t$ to $s$ is
$s\rightsquigarrow t$ and whose successful final states are boundary inputs.
Inputs are normal: any reverse step from an input would put that input in a
target list, contradicting target exclusion.
\lean{Graph/OperationPaths}{24}{29}{Net.reverseOperations}
\begin{lemma}[Determinism derived from incidence]
If selected operations are unary, reverse operation traversal is deterministic.
\end{lemma}
\begin{proof}
Two reverse steps from the same vertex select edges both targeting that vertex.
Target uniqueness identifies the edges. Their source list has length one, so
both predecessor vertices equal its only entry and hence equal one another.
There is no additional determinism premise.
\end{proof}
\lean{Graph/OperationPaths}{31}{44}{Net.reverseOperations_deterministic}

\begin{definition}[Edge-recording trace]\label{term:operation-trace}
A reverse operation trace from $v$ to $w$ records its actual edge list. The empty
trace starts and ends at the same vertex. A nonempty trace first chooses an
operation edge targeting the current vertex and a vertex in its source, then
continues recursively from that source vertex.
\end{definition}
\lean{Graph/OperationPaths}{47}{52}{Net.OperationTrace}
A trace with edge list $l$ yields exactly $|l|$ steps in the reverse machine.
Induct on its construction: the empty trace is a zero-step execution and each
trace constructor supplies the first step of the counted execution.
\lean{Graph/OperationPaths}{54}{58}{Net.OperationTrace.steps}

\begin{theorem}[Exact uniqueness to inputs]
Under the unary condition, two reverse traces from the same vertex to inputs
have equal edge lists and equal ending input vertices.
\end{theorem}
\begin{proof}
Induct on the first trace. If it is empty, the starting vertex is an input;
the second trace cannot be nonempty, since no input belongs to a target list.
Thus it too is empty and both conclusions follow. If the first trace is
nonempty, the second cannot be empty, for the same target-exclusion reason.
Their first edges are equal by target uniqueness. Their next vertices agree by
the determinism lemma. Apply the induction hypothesis to their tails and prefix
the common edge to the resulting list equality.
\end{proof}
\lean{Graph/OperationPaths}{62}{83}{Net.OperationTrace.unique_to_inputs}
\termdef{operation-path}
A forward operation path from $s$ to $t$ with edge list $l$ means a reverse trace
from $t$ to $s$ with list $\operatorname{reverse}(l)$. Empty paths are allowed in
this endpoint-sensitive notion, unlike the nonempty edge-only path predicate of
the preceding section.
\lean{Graph/OperationPaths}{86}{87}{Net.OperationPath}
Two forward operation paths from inputs to the same vertex therefore have equal
forward edge lists and equal starting input indices. Apply exact trace uniqueness,
use injectivity of list reversal, and then injectivity of the input constructor.
\lean{Graph/OperationPaths}{89}{94}{Net.operationPath_unique_from_inputs}

\begin{theorem}[No positive cycle on an input-rooted operation path]
If a vertex $v$ is reached by a forward operation path from an input, every
reverse operation trace from $v$ back to $v$ has an empty edge list.
\end{theorem}
\begin{proof}
The input-rooted path gives a reverse execution from $v$ to a successful input.
Prefix that execution by the proposed cycle to obtain another successful
execution from $v$ to the same input. Deterministic terminal uniqueness identifies
their lengths. Thus $|\mathit{cycle}|+|\mathit{path}|=|\mathit{path}|$, so the
cycle has length zero and is empty.
\end{proof}
\lean{Graph/OperationPaths}{99}{111}{Net.no_operation_cycle_on_input_path}
This is a statement about cycles at input-reachable vertices, not a global
acyclicity restriction on nets. The self-loop example below is a valid net with
a cycle and no boundary inputs.

\Needspace{18\baselineskip}
\section{Typed boxes, deep occurrences, and contexts}
\Needspace{24\baselineskip}
\subsection{Relabelling without changing incidence}
A signature-preserving map $f$ of operation alphabets relabels a net by replacing
each edge label $o$ with $f(o)$. All carriers, enumerations, vertex labels, and
incidence data are retained. Hence their invariant proofs are retained too.
For typing, compose the old typing equality with the inverse of the stipulated
signature equality for $f(o)$. This yields the new signature list at that edge.
\lean{Graph/Hierarchy}{17}{36}{Net.relabel}

\Needspace{24\baselineskip}
\subsection{The linear wire and label discipline}
\begin{definition}[Wire types, operation arities, and holes]
A wire type is either a term wire $*$ or a thunk wire $T^n(*)$, carrying a natural
number $n$ of bound variables. An operation specification assigns each operation
an argument count $a(o)$ and an ordered list of binder counts $b(o)$. A hole label
contains a natural-number name together with arbitrary input and output wire lists.
Equality of hole labels includes these lists as well as the name.
\end{definition}
\lean{Graph/Hierarchy}{44}{47}{Linear.Wire}
\lean{Graph/Hierarchy}{51}{53}{Linear.Operations}
\lean{Graph/Hierarchy}{56}{60}{Linear.Hole}
Primitive edge labels are operations, variables, focus labels from the abstract
three-mode type, or typed holes. This introduces concrete focus \emph{labels};
it does not yet define the UAM transition relation.
\lean{Graph/Hierarchy}{62}{66}{Linear.Primitive}
Their signature is as follows. Every non-hole primitive has one term input.
Operation $o$ has $a(o)$ term outputs followed by the thunk outputs specified in
$b(o)$. A variable or focus has one term output. A hole uses precisely the input
and output lists stored in its label. Ordering is part of each of these definitions.
\lean{Graph/Hierarchy}{68}{76}{Linear.primitiveSignature}

\begin{definition}[Alphabets and typed boxes]
An alphabet is a label type and a signature on it. A box over alphabet $A$ stores
numbers $n,m$ and a body net over $A$ of type
\[
 *\longrightarrow *^{n+m}.
\]
The corresponding outer box edge will have type
$T^n(*)\longrightarrow *^m$. Thus body boundary and outer boundary differ:
$n$ body outputs account for binding, and $m$ for exposed free variables.
\end{definition}
\lean{Graph/Hierarchy}{80}{82}{Linear.Alphabet}
\lean{Graph/Hierarchy}{84}{87}{Linear.Box}
All alphabet label types live in one fixed higher universe, while graph carriers
remain small types. The base alphabet contains just primitives, raised into that
universe by a universe lift; its signature reads the underlying primitive. The
successor alphabet is the disjoint union of primitives and boxes over the previous
alphabet. Its signature uses the primitive signature on the left and the stated
external box interface on the right.
\lean{Graph/Hierarchy}{89}{94}{Linear.baseAlphabet}
\lean{Graph/Hierarchy}{96}{105}{Linear.nextAlphabet}
Recursively, depth zero uses the base alphabet and depth $d+1$ uses the successor
of the depth-$d$ alphabet. A hypernet at depth bound $d$ is a net over that
alphabet. The bound limits possible nesting; it need not be the actual minimum
depth of the particular graph.
\lean{Graph/Hierarchy}{107}{109}{Linear.alphabet}
\lean{Graph/Hierarchy}{111}{112}{Linear.Hypernet}

The operation predicate selects primitive operation labels only. At depth zero
it inspects the universe-lifted primitive; at a successor depth it accepts only
the primitive-operation summand, excluding boxes and all other primitive kinds.
\lean{Graph/Hierarchy}{114}{118}{Linear.IsOperation}
Selected operation edges have exactly one source. Rewrite the incidence length
using signature typing, split on the depth and label constructors, discard labels
excluded by the predicate, and compute the operation input list $[*]$.
\lean{Graph/Hierarchy}{120}{131}{Linear.unary_operations}
It follows that two operation paths from inputs to the same vertex have identical
edge lists and input indices in every typed hypernet. Instantiate the generic
operation-path uniqueness theorem with the preceding derived unary proof. This
is the graph-theoretic uniqueness conclusion associated with the paper's Lemma~22
for this representation; neither determinism nor a locality field is assumed.
\lean{Graph/Hierarchy}{135}{140}{Linear.operation_path_unique}

\Needspace{24\baselineskip}
\subsection{Box construction and promotion of the depth bound}
The box constructor takes a body with the specified $n+m$ output interface and
forms the corolla of its box label in the successor alphabet. Corolla
well-formedness proves all outer graph conditions. The result has exactly one
outer box edge and the external type $T^n(*)\to *^m$.
\lean{Graph/Hierarchy}{143}{147}{Linear.box}
Promotion raises the available depth bound by one without changing the outer
graph. At depth zero, send a lifted primitive to the primitive summand at depth
one. At a successor depth, keep primitives and recursively promote each box
body, retaining its bound/free metadata. These maps preserve signatures because
neither primitive signatures nor external box arities change. Apply the generic
relabel construction at each recursive step.
\lean{Graph/Hierarchy}{151}{161}{Linear.raise}
Outer sources and targets are unchanged by promotion. The code records these
as heterogeneous equalities: the dependent graph types are indexed by different
depths, although after unfolding either defining case the incidence functions
are definitionally the same. Splitting on depth proves both equalities by
reflexivity.
\lean{Graph/Hierarchy}{163}{165}{Linear.raise_source}
\lean{Graph/Hierarchy}{167}{169}{Linear.raise_target}

\Needspace{24\baselineskip}
\subsection{Finite occurrence addresses at arbitrary depth}
\begin{definition}[All edge occurrences]
At depth zero, an occurrence is a shallow edge. At depth $d+1$, an occurrence is
either a shallow edge or a pair $(e,a)$ where $e$ is a box edge and $a$ an
occurrence in its body. A primitive edge has the empty type of child occurrences.
\end{definition}
\lean{Graph/Hierarchy}{173}{178}{Linear.AllEdge}
The parent edge is part of every deep address. In particular identical bodies
placed at different parent edges do not identify their child occurrences.

\begin{theorem}[Finiteness of all occurrences]
Every finite-depth hypernet has an explicit finite enumeration of its shallow
and deep edge occurrences.
\end{theorem}
\begin{proof}
Recurse on the depth bound. At zero use the graph's edge enumeration. At a
successor bound, enumerate shallow edges in the left summand. For the right
summand enumerate the dependent sum over outer edges: a primitive contributes
an empty enumeration, and a box contributes the recursively obtained enumeration
of its body. The sum and dependent-sum enumeration constructions already prove
completeness of this list. Recursion decreases the depth at every box descent.
\end{proof}
\lean{Graph/Hierarchy}{182}{190}{Linear.allEdgesFinite}
Every given occurrence belongs to this enumeration, by its completeness field.
This is the finiteness content of Appendix~A, Lemma~19, for the typed hierarchy.
The list need not be duplicate-free, since completeness suffices here.
\lean{Graph/Hierarchy}{192}{195}{Linear.every_edge_enumerated}

\Needspace{24\baselineskip}
\subsection{Hole lookup and the context condition}
A primitive hole lookup returns its stored label for a hole constructor and no
label otherwise. Full occurrence lookup uses that function on a shallow primitive,
returns no hole at a shallow box edge, and recursively looks inside a deep box
address. A purported child occurrence under a primitive has empty type and is
eliminated. Dependent transports in Lean express that the selected label really
is the box whose body is being searched.
\lean{Graph/Hierarchy}{197}{199}{Linear.primitiveHole}
\lean{Graph/Hierarchy}{202}{214}{Linear.holeAt}
\begin{definition}[Unique holes and contexts]
A hypernet has unique holes if, for each complete typed hole label $h$, any two
occurrences whose lookup returns $h$ are equal. A context is a hypernet together
with this property. This permits multiple different hole labels and also permits
no holes; it is an at-most-once condition for each label over all depths.
\end{definition}
\lean{Graph/Hierarchy}{218}{221}{Linear.UniqueHoles}
\lean{Graph/Hierarchy}{223}{224}{Linear.Context}
\begin{lemma}[Boxing preserves unique holes]
If a body has unique holes, its one-edge box has unique holes.
\end{lemma}
\begin{proof}
A shallow occurrence in the outer box is not a hole, so either proposed hole
occurrence must be deep. The outer graph has only one edge; therefore their
parent indices agree. Their lookups are the body lookups, whose occurrences
agree by the body's uniqueness assumption. Tagging these equal child addresses
with their common parent gives equality of the original occurrences.
\end{proof}
\lean{Graph/Hierarchy}{228}{249}{Linear.box_unique_holes}
Boxing a context now means boxing its underlying graph and supplying this lemma
as the context certificate. Shallow plugging at any fixed hierarchy depth is
just the generic typed net plugging construction with that depth's alphabet;
its interface constraints are checked by the type of its replacement argument.
\lean{Graph/Hierarchy}{251}{254}{Linear.boxContext}
\lean{Graph/Hierarchy}{257}{261}{Linear.plug}
This last operation returns a hypernet. The later context-calculus development
now proves that its result is a context under explicit freshness conditions
(Section~\ref{sec:context-closure}), including a sharp necessary-and-sufficient
condition for shallow plugging. The typed hierarchy has not yet been proved
adequate for the paper's original set-based presentation of graph-labelled edges.

\Needspace{18\baselineskip}
\section{Addressed deep replacement and concrete locality}
\Needspace{24\baselineskip}
\subsection{Changing one box body}
A per-edge signature-preserving relabelling generalizes alphabet relabelling:
it assigns a new label directly to each edge, rather than requiring the assignment
to factor through its old label. All carriers, incidence, and vertex labels stay
unchanged. The two new typing equations follow by composing each old typing
proof with the inverse of the specified per-edge signature equality.
\lean{Graph/DeepPlugging}{15}{35}{Net.relabelEdges}

\begin{definition}[Replacement of a box label]
Suppose outer edge $e$ is labelled by a box $(n,m,B)$. Given a new body $B'$ of
the same body type $*\to *^{n+m}$ and depth bound, change just the label at $e$
to $(n,m,B')$, retaining every other edge label.
\end{definition}
\begin{proof}[Well-formedness]
Use classical equality testing on outer edges to select $e$. At $e$ the old and
new labels have the same external input $[T^n(*)]$ and output $*^m$; this follows
from the given equality identifying the old label as the box. At every other
edge the label is unchanged. These two cases prove the per-edge signature
conditions for relabelling, which supplies the certified net.
\end{proof}
\lean{Graph/DeepPlugging}{45}{63}{Linear.replaceBox}
Selection uses an edge occurrence, not equality of box contents. Equal bodies
at other edges therefore do not cause simultaneous replacement.

\Needspace{24\baselineskip}
\subsection{Typed sites and replacement at a site}
\begin{definition}[Site]
A site in an outer hypernet $G$ carries a local depth bound and two wire-label
lists $J,K$. It is either a shallow edge of $G$, in which case $J,K$ are exactly
that edge's signature lists, or a box edge with a proof identifying its box
label together with a recursively chosen site in its body. In the latter case
the site's local depth and interface are inherited from the child site.
\end{definition}
\lean{Graph/DeepPlugging}{67}{75}{Linear.Site}
The indices prevent supplying a replacement of the wrong interface or local
depth. A site may select an arbitrary edge, not only an edge labelled as a hole;
the separate unique-hole context condition is not assumed here.

\begin{definition}[Deep plugging]
To replace at a shallow site use shallow net plugging. To replace at a site
inside a box, recursively replace in its body and use box-label replacement to
put that modified body back at the chosen outer edge. The recursion is on the
site's finite constructor structure, and the result retains the outer depth
bound and boundary interface.
\end{definition}
\lean{Graph/DeepPlugging}{79}{85}{Linear.plugAt}
At a shallow site this definition agrees with shallow plugging by reflexivity:
both sides unfold to the same term.
\lean{Graph/DeepPlugging}{87}{91}{Linear.plugAt_shallow}

\begin{theorem}[Locality at every enclosing box]
Replacement strictly inside a selected box leaves all outer source lists and
all outer target lists unchanged. It also leaves the label of every other outer
edge unchanged.
\end{theorem}
\begin{proof}
The inside branch changes only an edge label using per-edge relabelling. The
source and target fields of that construction are definitionally those of the
old graph, so their equalities are reflexive, independent of what the recursive
replacement did inside. For another outer edge $f\ne e$, the conditional defining
box-label replacement chooses the old-label branch, proving the label equation.
Repeated descent applies this fact at each enclosing layer.
\end{proof}
\lean{Graph/DeepPlugging}{94}{98}{Linear.plugAt_inside_source}
\lean{Graph/DeepPlugging}{100}{104}{Linear.plugAt_inside_target}
\lean{Graph/DeepPlugging}{108}{113}{Linear.plugAt_inside_other_label}
These are unconditional, concrete graph locality results about the defined
replacement operation. They do not yet prove UAM traversal locality, refocusing,
or the specimen interface laws: no concrete machine transition relation has
been constructed. The subsequent provenance development proves unique-hole context closure for
this concrete replacement at any depth. Associativity of context plugging and
a recursive isomorphism theorem for deep replacement remain separate tasks.

\Needspace{18\baselineskip}
\section{Concrete graph examples and regression proofs}
\Needspace{24\baselineskip}
\subsection{Self-loop gluing and the necessary reflection restriction}
Use a singleton wire-label set and a singleton operation alphabet, whose only
operation has one input and one output of that label. The graph named
$\mathsf{loop}$ has one internal vertex, one edge, no outer boundary, and that
vertex as both the edge's source and target.
\lean{Graph/Examples}{7}{9}{Examples.unary}
\lean{Graph/Examples}{12}{44}{Examples.loop}
Its finite carriers are both $\operatorname{Fin}(1)$. Singleton incidence lists
are distinct; all edges are equal since their indices are bounded by one. The
only internal vertex is incident on both sides, providing coverage, while
boundary cases have empty finite-index types. Label typing is reflexive. Thus
this is a well-formed net despite its cycle. Let $\mathsf{wireEdge}$ be the
corolla of the same unary operation, with its two distinct boundary vertices.
\lean{Graph/Examples}{46}{46}{Examples.wireEdge}

The corolla is not adjacent to itself: a shared target/source witness would have
to be both its output constructor and its input constructor, an impossible
constructor equality.
\lean{Graph/Examples}{49}{54}{Examples.corolla_not_self_adjacent}
When the corolla replaces the loop's sole edge, its input and output images are
equal. Both are, by definition, the same old internal vertex, because the host
hole's source and target port coincide. Hence the inserted edge is adjacent to
itself: use that old vertex as the witness and compute the mapped singleton
source and target lists.
\lean{Graph/Examples}{58}{60}{Examples.inserted_boundary_identified}
\lean{Graph/Examples}{62}{67}{Examples.inserted_self_adjacent}
Consequently the inserted two-entry list $[e,e]$ is a path, while $[e,e]$ is not
a path in the original corolla. The positive part uses its new self-adjacency and
a singleton tail path. The negative part extracts its first adjacency and
contradicts the corolla result. This proves that unconditional reflection of
inserted paths would be false; the non-self-loop hypothesis cannot simply be
omitted from that theorem.
\lean{Graph/Examples}{71}{76}{Examples.unconditional_path_reflection_fails}
The original loop itself has the two-entry path $[e,e]$, using its sole internal
vertex as the shared witness. This verifies that the graph representation has
not silently imposed global acyclicity.
\lean{Graph/Examples}{80}{81}{Examples.loop_is_cyclic}

The finite ordinal with one element is bijective to the unit type: send its
index to the sole unit element and send unit back to zero. The inverse laws
follow respectively from the ordinal bound and unit case analysis. An anonymous
example reindexes the corolla's edge carrier along this bijection and applies
shallow plugging congruence. It obtains isomorphic filled loops despite the
different replacement-edge carrier types, testing actual renaming rather than
reflexivity alone.
\lean{Graph/Examples}{83}{87}{Examples.oneToUnit}
\lean{Graph/Examples}{89}{91}{Examples.example (line 89)}

\Needspace{24\baselineskip}
\subsection{Values, boxes, and distinguished occurrence addresses}
For the linear examples, take two operation labels, value and other-value. Give
both zero term arguments and no thunk arguments; each therefore has signature
$*\to []$. The graph $\mathsf{value}$ is the corolla of the first operation at
depth zero.
\lean{Graph/Examples}{99}{100}{Linear.Examples.Op}
\lean{Graph/Examples}{102}{104}{Linear.Examples.operations}
\lean{Graph/Examples}{106}{107}{Linear.Examples.value}
Boxing it with zero bound and free variables yields a closed box of type
$T^0(*)\to []$ at depth one. A further example constructs the analogous graph
at depth bound two by putting a primitive value at depth one inside a box.
This is a larger \emph{available bound}, not a claim that the example has two
actually nested boxes.
\lean{Graph/Examples}{110}{110}{Linear.Examples.boxedValue}
\lean{Graph/Examples}{112}{115}{Linear.Examples.higherBoxedValue}
An anonymous proof verifies that any occurrence in the depth-one boxed value
belongs to the recursive occurrence enumeration, by the general completeness
theorem. The outer box edge is represented by a shallow left-tagged occurrence;
the value edge in its body is a right-tagged pair of parent edge and child edge.
\lean{Graph/Examples}{117}{118}{Linear.Examples.example (line 117)}
\lean{Graph/Examples}{122}{122}{Linear.Examples.outerOccurrence}
\lean{Graph/Examples}{124}{124}{Linear.Examples.innerOccurrence}
These two addresses are unequal. If equal, applying the sum-tag test would give
$\mathsf{true}=\mathsf{false}$, which is impossible. This checks separation of
outer and inner occurrences even though each individual edge carrier is singleton.
\lean{Graph/Examples}{126}{129}{Linear.Examples.outer_ne_inner}

The other-value corolla has exactly the same interface as the value corolla.
A site inside the boxed value descends through its only box edge and selects the
value edge in its body. Replacing at this site with other-value defines
$\mathsf{replacedInside}$.
\lean{Graph/Examples}{131}{132}{Linear.Examples.otherValue}
\lean{Graph/Examples}{135}{136}{Linear.Examples.insideValue}
\lean{Graph/Examples}{138}{139}{Linear.Examples.replacedInside}
Its outer source and target functions equal those of the original boxed value;
instantiate the two deep-locality theorems at the selected site. Its sole outer
edge label is the box with body $\mathsf{value}[e:=\mathsf{otherValue}]$.
For this label equation, unfold the site and replacement, simplify the equality
test selecting the sole outer edge, and use reflexivity for the resulting body.
The outer box is retained rather than replaced by the value operation itself.
\lean{Graph/Examples}{141}{142}{Linear.Examples.deep_replacement_source_unchanged}
\lean{Graph/Examples}{144}{145}{Linear.Examples.deep_replacement_target_unchanged}
\lean{Graph/Examples}{147}{151}{Linear.Examples.deep_replacement_body}

\Needspace{24\baselineskip}
\subsection{An input-rooted operation path}
Give the single unit operation one term argument and no binders. Its corolla
$\mathsf{argument}$ has type $*\to *$. The input, its only edge, and the output
form a forward operation path: read this as the one-step reverse trace, check
membership in the two singleton incidence lists, and end with an empty trace.
\lean{Graph/Examples}{153}{155}{Linear.Examples.argumentOperations}
\lean{Graph/Examples}{157}{158}{Linear.Examples.argument}
\lean{Graph/Examples}{160}{166}{Linear.Examples.argument_path}
Any reverse operation trace from its output back to that output has empty edge
list. Apply the no-cycle theorem to the constructed input-rooted path, deriving
the required unary condition from the typed signature. This exercises the
connection between concrete graph typing and generic machine terminal uniqueness.
\lean{Graph/Examples}{168}{172}{Linear.Examples.no_cycle_after_argument}

\Needspace{24\baselineskip}
\subsection{Filling a hole inside a certified context}
Choose the typed hole label with name zero, one term input, and no outputs.
Its corolla $\mathsf{singleHole}$ is a depth-zero hypernet. It has unique holes
because any two edge occurrences are equal: both are elements of
$\operatorname{Fin}(1)$, so their underlying indices are zero. Pairing this graph
with the proof constructs a context.
\lean{Graph/Examples}{174}{174}{Linear.Examples.holeLabel}
\lean{Graph/Examples}{176}{177}{Linear.Examples.singleHole}
\lean{Graph/Examples}{179}{184}{Linear.Examples.singleHole_unique}
\lean{Graph/Examples}{186}{186}{Linear.Examples.singleContext}
Box the graph with zero bound and free variables to obtain $\mathsf{boxedHole}$,
and box the certified context to obtain $\mathsf{boxedContext}$. The unique-hole
certificate in the latter is supplied by the general boxing theorem. Looking
up the deep address consisting of the only outer edge and the only body edge
returns the chosen hole label by definition.
\lean{Graph/Examples}{188}{189}{Linear.Examples.boxedHole}
\lean{Graph/Examples}{191}{192}{Linear.Examples.boxedContext}
\lean{Graph/Examples}{194}{195}{Linear.Examples.hole_is_deep}
The corresponding site descends into that box and chooses its shallow hole edge.
Filling this site with the value corolla produces $\mathsf{filledBox}$. Its outer
label is the same zero-bound, zero-free box, now containing
$\mathsf{singleHole}[e:=\mathsf{value}]$; unfold deep replacement and simplify the
selected-edge equality test to prove this equation.
\lean{Graph/Examples}{197}{198}{Linear.Examples.boxedHoleSite}
\lean{Graph/Examples}{200}{200}{Linear.Examples.filledBox}
\lean{Graph/Examples}{202}{206}{Linear.Examples.filled_box_body}
Finally every edge of this filled body has no primitive hole label. Split on
its edge tag. An old surviving edge would be an element of the singleton old
carrier different from its only element, contradicting the surviving-edge
certificate. An inserted value edge is an operation, whose primitive hole lookup
is definitionally empty. Thus this example proves removal of the old hole, not
merely the well-typedness of the filled graph.
\lean{Graph/Examples}{210}{216}{Linear.Examples.filled_body_no_holes}

\Needspace{18\baselineskip}
\section{Occurrence fibres and transport of hole observations}
The next milestone completes context closure under explicit hole freshness,
including replacement at arbitrary box depth. Its proofs need to distinguish
individual occurrences, rather than merely sets of labels. This section provides
coordinates for those occurrences without changing the existing hypernet type.

\Needspace{24\baselineskip}
\subsection{Coordinates local to an outer edge}
\begin{definition}[Children and occurrence fibre]\label{term:fibre}
For a label $l$ at depth zero, its child occurrence type is empty. At depth $d+1$,
a primitive label again has no children; a box label with body $B$ has child type
$\operatorname{AllEdge}(B)$. Define the fibre of $l$ to be
\[
 F_d(l)=\{\bullet\}+\operatorname{Children}_d(l).
\]
The left element is the edge itself; a right element is an occurrence below it.
\end{definition}
\lean{Graph/Occurrences}{12}{16}{Linear.Children}
\lean{Graph/Occurrences}{19}{20}{Linear.Fiber}
\termdef{fibre-lookup}
The shallow hole observation reads the primitive hole label at depth zero or in
the primitive summand at positive depth, and returns no hole on a box label.
The child observation eliminates empty child types and otherwise uses the
body's existing hole lookup. The fibre observation combines these functions:
left elements use shallow lookup and right elements use child lookup.
\lean{Graph/Occurrences}{22}{25}{Linear.shallowHole}
\lean{Graph/Occurrences}{27}{31}{Linear.childHole}
\lean{Graph/Occurrences}{33}{36}{Linear.fiberHole}
\termdef{occurrence-coordinates}
An occurrence of a graph $G$ in fibre coordinates is a dependent pair
\[
 (e,x),\qquad e\in G_E,\quad x\in F_d(\lambda_G(e)).
\]
The first component is always an outer edge, even when $x$ goes deep into a box.
This representation depends on edge labels and occurrence addresses, not on
vertex incidence or a particular enumeration of the carriers.
\lean{Graph/Occurrences}{39}{40}{Linear.Occurrences}

\begin{lemma}[Occurrence coordinates]
The existing occurrence type $\operatorname{AllEdge}(G)$ is bijective to the
fibre-coordinate occurrence type.
\end{lemma}
\begin{proof}
At depth zero, send an edge $e$ to $(e,\mathsf{left}(\bullet))$; the inverse
forgets that singleton component. A right component would belong to the empty
type and is impossible. At positive depth, send a shallow edge to the same
left-tagged pair, and a deep address $(e,a)$ to $(e,\mathsf{right}(a))$.
The inverse removes the fibre tag to reconstruct the corresponding shallow or
deep address. Both inverse laws follow by cases on the tags and uniqueness of the unit value.
\end{proof}
\lean{Graph/Occurrences}{42}{60}{Linear.occurrenceBijection}
Define a graph's occurrence observation at $(e,x)$ by applying the fibre
observation to $\lambda_G(e)$ and $x$.
\lean{Graph/Occurrences}{62}{63}{Linear.occurrenceHole}
The bijection preserves hole observations. At depth zero this is definitional.
For a positive-depth shallow occurrence, split the label into primitive or box
and compute both lookups. For a deep occurrence, both lookups read the same
body address: exposing the label cases makes the equality reflexive. The original
hierarchy lookup now uses a directly dependent reader, avoiding unnecessary
proof-dependent casts while retaining the same definition by label cases.
\lean{Graph/Occurrences}{65}{77}{Linear.occurrenceBijection_hole}

\begin{lemma}[The context condition in fibre coordinates]
A graph has unique holes exactly when any two fibre-coordinate occurrences
observing the same hole label are equal.
\end{lemma}
\begin{proof}
For the forward direction, apply unique holes to the inverse images of the two
occurrences. Observation preservation supplies their lookup hypotheses. Map
the resulting equality forwards and cancel the inverse pairs. Conversely, map
the two original addresses forwards, use fibre-coordinate uniqueness and
observation preservation, and conclude by injectivity of the bijection.
\end{proof}
\lean{Graph/Occurrences}{80}{96}{Linear.uniqueHoles_iff_occurrences}

\Needspace{24\baselineskip}
\subsection{Transport along equality and edge maps}
If labels $l,l'$ are equal, transport of $x\in F_d(l)$ along that equality gives
an element of $F_d(l')$ with the same hole observation. Such transport is injective.
For each assertion, eliminate the label equality; transport becomes the identity,
so the observation equation is reflexive and injectivity is the given equality
of transported elements.
\lean{Graph/Occurrences}{98}{101}{Linear.fiberHole_transport}
\lean{Graph/Occurrences}{103}{106}{Linear.fiber_transport_injective}
A shallow edge map $f:G_E\to H_E$ with equations
$\lambda_G(e)=\lambda_H(f(e))$ maps all occurrences by
\[
 (e,x)\longmapsto (f(e),\operatorname{transport}(x)).
\]
The descendant address is preserved through the label equality; it is not
searched for again by comparing box contents.
\lean{Graph/Occurrences}{109}{112}{Linear.mapOccurrences}
This map preserves hole observation by the equality-transport observation lemma.
If $f$ is injective, the occurrence map is injective: equality of mapped pairs
first identifies their edge indices by injectivity of $f$. After substituting
that equality, equality of their transported second components implies equality
of the original components by transport injectivity. The dependent pairs are
therefore equal.
\lean{Graph/Occurrences}{114}{119}{Linear.mapOccurrences_hole}
\lean{Graph/Occurrences}{121}{132}{Linear.mapOccurrences_injective}

\Needspace{18\baselineskip}
\section{Hole inventories and sharp shallow context closure}
\label{sec:context-closure}
All hole-label comparisons in this section compare complete typed labels: the
name and both interface lists. The occurrence quantifiers include every box depth.

\Needspace{24\baselineskip}
\subsection{Hole inventories, absence, and isomorphism}
\begin{definition}[Presence, absence, and disjointness]
Write $\mathsf{HasHole}(G,h)$ when some occurrence of $G$ looks up to $h$.
Write $\mathsf{HoleFree}(G)$ when no label has such an occurrence. Two graphs
have disjoint holes when no complete hole label occurs in both. The graphs may
have different external interfaces and depth bounds.
\end{definition}
\lean{Graph/ContextPlugging}{11}{12}{Linear.HasHole}
\lean{Graph/ContextPlugging}{14}{15}{Linear.HoleFree}
\lean{Graph/ContextPlugging}{18}{19}{Linear.HolesDisjoint}
Hole presence is equivalently witnessed by a fibre-coordinate occurrence with
the same observation. Transport a witness by the occurrence bijection for one
direction and by its inverse for the other, using observation preservation and
the inverse equation.
\lean{Graph/ContextPlugging}{21}{31}{Linear.hasHole_iff_occurrences}
A hole-free graph has unique holes vacuously: the first purported hole occurrence
already contradicts absence. Disjointness is symmetric by exchanging its two
presence premises. Any graph is disjoint from a hole-free graph because a hole
in the second graph would contradict its absence property.
\lean{Graph/ContextPlugging}{33}{36}{Linear.holeFree_unique}
\lean{Graph/ContextPlugging}{38}{40}{Linear.HolesDisjoint.symm}
\lean{Graph/ContextPlugging}{42}{44}{Linear.holeFree_disjoint}

Boxing neither creates nor removes a hole label from the inventory of the body.
A shallow box occurrence cannot be a hole. A deep occurrence is exactly an
address in its body with the same lookup, so it supplies a body witness. Conversely,
tag a body witness with the box's only outer edge. Consequently a hole-free body
gives a hole-free box: transfer any proposed hole witness to the body and contradict
its absence property.
\lean{Graph/ContextPlugging}{47}{58}{Linear.box_hasHole_iff}
\lean{Graph/ContextPlugging}{60}{63}{Linear.box_holeFree}

A label-preserving edge map preserves hole presence: map the fibre-coordinate
witness using the occurrence map and its observation equation. An injective such
edge map reflects unique holes: map two proposed equal-label occurrences into
the target, apply target uniqueness, and use occurrence-map injectivity.
\lean{Graph/ContextPlugging}{66}{71}{Linear.hasHole_map}
\lean{Graph/ContextPlugging}{74}{82}{Linear.uniqueHoles_of_embedding}
It follows that graph isomorphisms preserve and reflect hole inventories, and
preserve and reflect the context condition. Apply presence preservation to the
forward and inverse edge bijections for the first equivalence. Apply reflection
of uniqueness to the inverse bijection for forward preservation and to the
forward bijection for reflection. Label preservation is supplied by the
isomorphism equations with the appropriate equality direction.
\lean{Graph/ContextPlugging}{84}{87}{Linear.iso_hasHole_iff}
\lean{Graph/ContextPlugging}{89}{94}{Linear.iso_uniqueHoles_iff}
These results use the current isomorphism notion that fixes edge labels. They do
not yet extend it to isomorphic, independently renamed box contents.

\Needspace{24\baselineskip}
\subsection{Provenance and the abstract collision argument}
\begin{definition}[Provenance certificate]
For result $R$, host $C$, and replacement $H$, a provenance certificate is an
injective function
\[
 p:\operatorname{Occurrences}(R)\longrightarrow
 \operatorname{Occurrences}(C)+\operatorname{Occurrences}(H)
\]
that preserves hole observation in the chosen summand. Its left values identify
retained host occurrences and its right values identify inserted occurrences.
Different result occurrences must have different tagged origins.
\end{definition}
\lean{Graph/ContextPlugging}{98}{106}{Linear.Provenance}
This is a reusable intermediate interface. Unlike the earlier specimen laws,
the provenance certificates required for actual plugging are constructed below
from the concrete operations, rather than left as unproved assumptions.

Every result hole label belongs to the host or replacement inventory. Choose its
fibre-coordinate witness, inspect the origin tag, and use preservation to obtain
a witness in the indicated graph.
\lean{Graph/ContextPlugging}{113}{119}{Linear.Provenance.holes_subset}
A provenance certificate is \emph{compatible} when no inserted hole label is
observed at any retained host origin. Formally, if $p(a)=\mathsf{left}(c)$ and
$c$ observes $h$, then $H$ does not contain $h$. Host occurrences outside the
image of the left branch are irrelevant to this condition.
\lean{Graph/ContextPlugging}{122}{124}{Linear.Provenance.Compatible}

\begin{theorem}[Unique holes from compatible provenance]
If $C,H$ have unique holes and a provenance certificate for $R,C,H$ is compatible,
then $R$ has unique holes.
\end{theorem}
\begin{proof}
Take two result occurrences observing the same $h$ and inspect their origins.
If both origins are retained, host uniqueness identifies their host addresses.
If both are inserted, replacement uniqueness identifies their replacement
addresses. In either case the tagged origins agree and provenance injectivity
identifies the result occurrences. If one origin is retained and one inserted,
observation preservation gives a retained host observation of $h$ and a hole
witness for $h$ in $H$, contradicting compatibility. This covers all four tag
combinations and proves fibre-coordinate uniqueness, hence the original condition.
\end{proof}
\lean{Graph/ContextPlugging}{126}{151}{Linear.Provenance.unique_of_compatible}
Full disjointness of the host and replacement inventories implies compatibility:
any retained host occurrence is in particular a host occurrence. Substituting
this fact into the preceding theorem gives unique holes under full disjointness.
\lean{Graph/ContextPlugging}{154}{157}{Linear.Provenance.unique}

\Needspace{24\baselineskip}
\subsection{Concrete shallow provenance and a first context operation}
The provenance map for $C[e:=H]$ reads its outer edge tag. A surviving host edge
with fibre address $x$ maps to the corresponding old host edge with $x$;
an inserted edge maps to its replacement edge with $x$. In either case the
edge label, and hence the fibre, is definitionally unchanged.
\begin{proof}[Injectivity and observation preservation]
Mixed edge tags give different origin tags and cannot agree. Two retained
origins first identify the old edge values. Their surviving-edge subtype values
then agree because their inequality certificates are propositions; equality of
the dependent second components identifies their fibre addresses. Two inserted
origins are already the same replacement pair. In every case fibre observation
is unchanged by definition. Thus these data construct a provenance certificate
for any shallow replacement, including when the removed edge is a box.
\end{proof}
\lean{Graph/ContextPlugging}{163}{190}{Linear.shallowProvenance}
The general provenance theorem now proves context closure for shallow plugging
under full hole disjointness. A context-valued operation packages the filled
graph with this proof, taking certified host and replacement contexts and a
proof of disjointness. All graph incidence invariants still come from the earlier
certified net plugging construction.
\lean{Graph/ContextPlugging}{193}{198}{Linear.plug_uniqueHoles}
\lean{Graph/ContextPlugging}{200}{204}{Linear.Context.plug}

\Needspace{24\baselineskip}
\subsection{Exact inventory and the sharp freshness condition}
\begin{definition}[Retained hole label]
A retained hole label for removal of edge $e$ in $C$ is a hole observation in the
fibre of an outer edge $f\ne e$. This includes deep occurrences below $f$.
It excludes both the removed edge and every occurrence inside its label.
\end{definition}
\lean{Graph/ContextPlugging}{207}{209}{Linear.RetainedHole}
\begin{theorem}[Exact shallow hole inventory]
For each label $h$,
\[
 \mathsf{HasHole}(C[e:=H],h)
 \quad\Longleftrightarrow\quad
 \mathsf{RetainedHole}(C,e,h)\ \lor\ \mathsf{HasHole}(H,h).
\]
\end{theorem}
\begin{proof}
Express presence using fibre coordinates. A result occurrence has an outer edge
tag that is either a surviving host edge or an inserted edge. The first case
supplies exactly the retained witness; the second supplies a replacement witness.
Conversely, a retained witness gives a result occurrence at the left edge tag,
using its inequality proof to construct the surviving-edge subtype. A replacement
witness gives a result occurrence at the right tag. Their hole observations
are unchanged by definition.
\end{proof}
\lean{Graph/ContextPlugging}{212}{229}{Linear.plug_hasHole_iff}

Define sharp plugging freshness to mean that no retained host hole label occurs
in the replacement. This is weaker than requiring the replacement to be disjoint
from the entire original host: labels confined to the removed edge's region may
be reused.
\lean{Graph/ContextPlugging}{232}{235}{Linear.PlugFresh}
Sharp freshness implies uniqueness of the filled graph when both inputs have
unique holes. Apply the compatible-provenance theorem to shallow provenance.
Any result occurrence with a retained origin must have a surviving host edge
tag; its fibre is the retained witness required by sharp freshness. An inserted
edge cannot have a retained origin, so the other case is impossible.
\lean{Graph/ContextPlugging}{237}{249}{Linear.plug_uniqueHoles_of_fresh}

\begin{theorem}[Sharp freshness is necessary]
If $C[e:=H]$ has unique holes, sharp plugging freshness holds.
\end{theorem}
\begin{proof}
Suppose a retained host fibre observes $h$ and $H$ contains $h$. Form the two
result occurrences from their witnesses as in the inventory theorem. Result
uniqueness identifies them. Equality of their first components would identify
a surviving host edge tag with an inserted edge tag, which is impossible.
\end{proof}
\lean{Graph/ContextPlugging}{252}{261}{Linear.plug_fresh_necessary}
Consequently, for two contexts the filled graph has unique holes if and only if
sharp plugging freshness holds. Necessity is the preceding theorem, which did
not require uniqueness of either input; sufficiency uses the two input certificates.
\lean{Graph/ContextPlugging}{264}{269}{Linear.plug_uniqueHoles_iff_fresh}

\Needspace{24\baselineskip}
\subsection{Consumption of a selected hole and certified composition}
In a context, if the removed edge's shallow lookup is $h$, that label has no
retained occurrence elsewhere. Otherwise the retained occurrence and the selected
edge's shallow occurrence would be two occurrences of $h$. Context uniqueness
identifies their pairs, so their outer edges agree, contradicting the retained
edge's inequality certificate.
\lean{Graph/ContextPlugging}{272}{277}{Linear.removed_hole_not_retained}
It follows that the consumed label $h$ occurs in $C[e:=H]$ exactly when it occurs
in $H$. Use the exact inventory theorem and exclude its retained alternative
by the preceding result. Thus consumption removes the old occurrence; it does
not forbid the replacement from deliberately reintroducing that label.
\lean{Graph/ContextPlugging}{280}{287}{Linear.plug_consumed_hole_iff}
The operation $\mathsf{Context.plugFresh}$ packages shallow composition with this
sharp freshness premise. Its result graph is ordinary net plugging, and its
context certificate is the sharp sufficiency proof. It supports arbitrary finite
depth and arbitrary selected outer edges, rather than requiring that the selected
edge's label itself be a hole.
\lean{Graph/ContextPlugging}{289}{293}{Linear.Context.plugFresh}

\Needspace{18\baselineskip}
\section{Provenance and context closure at arbitrary box depth}
The shallow origin map alone is insufficient for deep replacement: an enclosing
box edge stays in place while occurrences in its body may change. We construct
provenance through each selected enclosing box and then induct on the typed site.

\Needspace{24\baselineskip}
\subsection{Local provenance for a changed label}
\begin{definition}[Fibre provenance]\label{term:fibre-provenance}
For old and new labels $l,l'$ at the same outer depth and a replacement $H$, a
fibre provenance certificate is an injective map
\[
 F_d(l')\longrightarrow F_d(l)+\operatorname{Occurrences}(H)
\]
preserving the hole observation in the indicated summand. It accounts for
changing one label while retaining or inserting the occurrences within that label.
\end{definition}
\lean{Graph/DeepContext}{12}{19}{Linear.FiberProvenance}
An unchanged label has the left inclusion as its origin map, which is injective
and preserves observation by reflexivity. If the new label merely has a proof
of equality to the old label, first transport its fibre address along that
equality and then use the left inclusion. Injectivity follows from transport
injectivity; the hole equation is the inverse direction of the transport
observation equation.
\lean{Graph/DeepContext}{22}{26}{Linear.FiberProvenance.unchanged}
\lean{Graph/DeepContext}{29}{33}{Linear.FiberProvenance.retained}

\begin{lemma}[Assembly of local provenance]
Suppose a signature-preserving per-edge relabelling of $G$ has a fibre provenance
certificate at every edge. If only a single designated edge can produce inserted
origins, these certificates assemble into provenance for the relabelled graph
relative to $G,H$.
\end{lemma}
\begin{proof}
At a result pair $(e,x)$ inspect its fibre origin. A retained origin $y$ gives
the old graph occurrence $(e,y)$; an inserted origin is already an occurrence
of $H$. Observation preservation is precisely the local certificate's equation.
For injectivity, consider the two output tags. Two retained origins identify
their outer edges and then their local retained addresses; local injectivity
identifies the original fibre addresses. Mixed output tags cannot agree.
If both origins are inserted, each outer edge must be the designated one:
otherwise the hypothesis that all its origins are retained contradicts its
inserted tag. With the outer edges identified, equality of inserted addresses
and local injectivity again identify the fibre addresses. Thus the global
origin map is injective in every case.
\end{proof}
\lean{Graph/DeepContext}{36}{91}{Linear.relabelProvenance}
The single-designated-edge restriction prevents two different outer edges from
claiming the same inserted occurrence. The lemma proves this separation; it
does not assume global injectivity as an extra graph law.

\Needspace{24\baselineskip}
\subsection{Lifting a body's provenance through one box}
Suppose the old and new bodies $B,B'$ have the same bound/free interface and a
provenance certificate $p$ for $B',B,H$. Their outer box labels differ only in
body content. Construct fibre provenance as follows. The shallow occurrence of
the box is retained as the old box's shallow occurrence. A new child address is
converted to the fibre-coordinate address of $B'$ and sent through $p$. Retained
body origins are converted back to body addresses and tagged as old box children;
inserted origins remain occurrences of $H$.
\begin{proof}[Injectivity and observation preservation]
A retained shallow box occurrence is distinct from every retained child and
every inserted origin by its tags. For two children, inspect the two body-origin
tags. Mixed tags cannot give equal lifted origins. If both are retained,
equality of the lifted children, followed by the old body's occurrence bijection,
identifies the retained body origins. If both are inserted, their equality is
already equality of body origins. In both cases body-provenance injectivity and
the new body's occurrence-bijection injectivity identify the original children.
For observations, the shallow box lookup is empty on both sides. A child's lookup
agrees with its body-coordinate observation by the bijection lemma; body
provenance gives the intermediate equality. In the retained case the inverse
bijection of the old body supplies the remaining observation equation. In the
inserted case the body equation already gives the replacement observation.
\end{proof}
\lean{Graph/DeepContext}{95}{147}{Linear.boxFiberProvenance}

To lift through a specified edge of a larger graph, use this box certificate at
that edge and an equality-transported retained certificate at every other edge.
Classical testing of edge equality selects the two cases. At the selected edge,
the given label equality identifies its old box, and the new label contains the
modified body with unchanged metadata. Elsewhere the label is unchanged, so all
local origins are retained. The assembly lemma therefore constructs provenance
for the existing concrete box-replacement operation.
\lean{Graph/DeepContext}{150}{170}{Linear.replaceBoxProvenance}

\Needspace{24\baselineskip}
\subsection{The deep context closure theorem}
\begin{theorem}[Provenance at every typed site]
For every typed site $s$ in $G$ and matching replacement $H$, the defined result
$\mathsf{plugAt}(s,H)$ has injective, hole-observation-preserving provenance
relative to $G,H$.
\end{theorem}
\begin{proof}
Induct on the site. A shallow site uses the concrete shallow provenance map.
For an inside site, the induction hypothesis gives provenance for the replacement
in the chosen body's inner site. Lift this certificate through the selected
outer box using the preceding construction. That is exactly the inside branch
of the definition of deep plugging. Induction follows the finite site address,
so the proof works at every finite nesting depth.
\end{proof}
\lean{Graph/DeepContext}{173}{180}{Linear.siteProvenance}

\begin{theorem}[Deep replacement preserves contexts under compatibility]
If $G,H$ have unique holes and the concrete site provenance is compatible, then
$\mathsf{plugAt}(s,H)$ has unique holes.
\end{theorem}
\begin{proof}
Use the constructed site provenance in the compatible-provenance uniqueness
theorem. No locality or occurrence-injectivity hypothesis remains to be supplied:
those properties have been proved for the actual replacement operation.
The remaining compatibility premise is explicitly a restriction on colliding
hole labels at retained origins.
\end{proof}
\lean{Graph/DeepContext}{183}{188}{Linear.plugAt_uniqueHoles_of_compatible}
Full disjointness of $G,H$ gives a simpler sufficient condition: instantiate the
full-disjointness provenance theorem with the same site certificate. This
condition is stronger than compatibility because it also inspects discarded host
occurrences; it is not asserted to be necessary.
\lean{Graph/DeepContext}{190}{194}{Linear.plugAt_uniqueHoles}
Every hole label of the deep result belongs to the inventory of the host or
replacement, by the provenance inventory-subset theorem. This is an upper bound
on the deep inventory, not an exact description of which host labels survive.
\lean{Graph/DeepContext}{197}{200}{Linear.plugAt_hasHole_subset}
If $H$ is hole-free, it has unique holes and is disjoint from $G$. Substituting
these two facts gives context closure for replacing any addressed edge by
hole-free material, with no separate freshness witness required.
\lean{Graph/DeepContext}{203}{207}{Linear.plugAt_holeFree_unique}
Finally, $\mathsf{Context.plugAt}$ accepts a certified host context, a typed site,
a certified replacement context, and full disjointness. It returns the concrete
deeply filled graph paired with its proved context certificate. The original
ordered outer interface and enclosing incidence laws are those already supplied
by deep replacement.
\lean{Graph/DeepContext}{209}{212}{Linear.Context.plugAt}

This provenance construction supplies globally unique-hole context closure and
concrete origin tracking through boxes. Section~\ref{sec:deep-inventory}
strengthens the inventory bound to an exact syntactic characterization and gives
a necessary-and-sufficient freshness condition at every finite box depth.
Associativity of plugging and recursive isomorphism of box contents remain
separate obligations. None of these results yet defines the concrete UAM.

\Needspace{18\baselineskip}
\section{Context composition: absence, reuse, and a collision}
\Needspace{24\baselineskip}
\subsection{Hole-free filling inside a box}
The value corolla is hole-free: any hole witness would claim that the lookup
of a primitive operation is a hole label, whereas that lookup is empty.
Applying deep context closure for a hole-free replacement to the old boxed-hole
site proves that its filled result has unique holes. The context-valued operation
also builds a certified filled context from the boxed context and the value
corolla, using hole-free uniqueness and disjointness for the replacement.
\lean{Graph/ContextExamples}{10}{12}{Linear.Examples.value_holeFree}
\lean{Graph/ContextExamples}{14}{15}{Linear.Examples.filledBox_unique}
\lean{Graph/ContextExamples}{17}{19}{Linear.Examples.filledContext}

The single-hole corolla has no retained occurrence when its only outer edge is
removed. A proposed retained witness includes an edge different from zero, but
every element of its edge carrier $\operatorname{Fin}(1)$ equals zero. This
contradicts the witness's inequality certificate.
\lean{Graph/ContextExamples}{21}{25}{Linear.Examples.singleHole_retained_absent}
Consequently replacing its only hole by the very same single-hole corolla still
gives a context. Use sharp context closure: both input uniqueness proofs are
available, and the retained part of the freshness premise is impossible by the
preceding result. However, these two identical input graphs do not have fully
disjoint inventories: their sole edge is a witness for the same label in each.
Applying a proposed full-disjointness proof to those two witnesses gives a
contradiction. Together these examples show that sharp freshness admits valid
label reuse which the stronger sufficient condition would exclude.
\lean{Graph/ContextExamples}{28}{30}{Linear.Examples.reuse_removed_label}
\lean{Graph/ContextExamples}{33}{34}{Linear.Examples.reuse_not_fully_disjoint}

Filling the single-hole corolla with the value is hole-free, not just a context.
The exact inventory theorem leaves a retained alternative or a value alternative;
the first is impossible by absence of surviving old edges and the second by
hole-freeness of the value.
\lean{Graph/ContextExamples}{36}{40}{Linear.Examples.filled_single_holeFree}
The boxed filled result is hole-free at every depth as well. Express any proposed
hole in fibre coordinates. Its outer edge has singleton carrier, so it is zero.
The already proved label equation identifies that edge as a box whose body is
the filled single-hole graph. Transport the fibre address along that equation,
which preserves its hole observation. Its shallow alternative cannot observe a
hole because it is a box; its child alternative contradicts the filled body's
hole-freeness. Thus no shallow or deep witness exists in the result.
\lean{Graph/ContextExamples}{43}{55}{Linear.Examples.filledBox_holeFree}

\Needspace{24\baselineskip}
\subsection{Two distinct shallow holes in series}
For a natural-number name $n$, define a hole label $\chi_n$ with one term input
and one term output. The two-hole graph has external type $*\to *$, one internal
term vertex $v$, and edges $e_0,e_1$ with labels $\chi_0,\chi_1$ and incidence
\[
 \mathsf{input}_0\xrightarrow{e_0}v\xrightarrow{e_1}\mathsf{output}_0.
\]
This notation records the two source and target singleton lists, not an extra
path assumption.
\lean{Graph/ContextExamples}{58}{58}{Linear.Examples.chainHole}
\lean{Graph/ContextExamples}{61}{128}{Linear.Examples.twoHoles}
\begin{proof}[Graph invariants]
The carriers are $\operatorname{Fin}(1)$ and $\operatorname{Fin}(2)$, with finite
range enumerations. Every incidence list is singleton and hence distinct.
For source uniqueness, two edges of different indices have different source
constructors, namely input and internal, so a shared source is impossible;
equal-index cases give equal edges. The target argument uses internal and output
constructors instead. Coverage chooses $e_0$ for the input source and internal
target, and $e_1$ for the internal source and output target. Outer outputs have
no source occurrence and outer inputs no target occurrence by constructor
disjointness. There is only one internal index and one index at each boundary,
so these cases exhaust all vertices. Every displayed vertex label is term,
matching the singleton input and output lists of both typed hole labels.
\end{proof}

This graph has unique holes. Two equal observed labels give equal numeric hole
names. If one edge index is zero and the other is not, their names are respectively
zero and one, a contradiction. If both are zero they agree directly. If neither
is zero, both are one because their ordinal carrier has size two, so they agree
again. At depth zero these are all occurrence addresses.
\lean{Graph/ContextExamples}{130}{149}{Linear.Examples.twoHoles_unique}

\Needspace{24\baselineskip}
\subsection{A checked failure without freshness}
Let the replacement be the corolla of $\chi_1$, with the same external interface
as $e_0$. It is itself a context: any two occurrences belong to its singleton
edge carrier and hence are equal, independently of their lookup hypotheses.
\lean{Graph/ContextExamples}{151}{152}{Linear.Examples.repeatedHole}
\lean{Graph/ContextExamples}{154}{159}{Linear.Examples.repeatedHole_unique}
Nevertheless, replacing $e_0$ in the two-hole graph by this corolla violates
unique holes. In the result, the retained host edge $e_1$ and the inserted
replacement edge both observe $\chi_1$. Context uniqueness would identify their
occurrences, but their edge tags are different constructors of the disjoint sum.
This is a contradiction. Thus well-formed graph plugging and uniqueness of each
input graph alone do not guarantee the context condition in the result.
\lean{Graph/ContextExamples}{162}{167}{Linear.Examples.collision_breaks_unique}
The same witnesses explicitly refute sharp freshness: $e_1$ and its shallow fibre
address witness retained presence of $\chi_1$, and the replacement's only edge
witnesses inserted presence. Feeding both to a proposed freshness proof gives
false. This regression exercises the actual collision excluded by the
necessary-and-sufficient theorem, alongside the valid label-reuse example above.
\lean{Graph/ContextExamples}{169}{172}{Linear.Examples.collision_not_fresh}

\Needspace{18\baselineskip}
\section{Exact inventories and sharp context closure at every depth}
\label{sec:deep-inventory}

The provenance argument establishes a useful upper bound on hole inventories.
We now strengthen it to an equality described entirely by the original host and
the selected address. This also gives a necessary-and-sufficient freshness
condition. Throughout this section $s$ is a site in $G$, $H$ has the interface
and local depth required by $s$, and $G[s:=H]$ denotes $\mathsf{plugAt}(s,H)$.
Write $\mathcal H(G)$ for the set of full typed hole labels present in $G$.
These statements concern the concrete typed hierarchy; they do not presuppose
any specimen laws or a concrete UAM semantics.

\subsection{Separating one outer edge from the rest of a graph}

\begin{definition}[Inventory of an edge]\label{term:edge-hole-inventory}
For an outer edge $e$ of $G$, let $E_G(e,h)$ mean that some position in the fibre
of $e$ observes $h$. This includes the edge itself and all descendants within
its label. Thus $E_G(e,h)$ is the predicate $\mathsf{EdgeHasHole}(G,e,h)$.
\end{definition}
\lean{Graph/DeepInventory}{9}{10}{Linear.EdgeHasHole}

\begin{definition}[Uniqueness away from an edge]\label{term:away-unique}
The predicate $\mathsf{AwayUnique}(G,e)$ says that any two occurrences with outer
edges different from $e$ and with the same hole observation are equal. It places
no condition on occurrences within the fibre of $e$.
\end{definition}
\lean{Graph/DeepInventory}{13}{15}{Linear.AwayUnique}

\begin{lemma}[Partition of uniqueness]
The graph $G$ has unique holes if and only if all three conditions hold:
occurrences away from $e$ are unique; hole observations within the fibre of $e$
are unique; and no label is observed both away from $e$ and within that fibre.
\end{lemma}
\begin{proof}
Global uniqueness restricts to each part. A collision between the two parts
would force equality of their outer edges, contradicting the requirement that
the first outer edge differ from $e$. Conversely, take two occurrences of the
same hole label. If both lie away from $e$, use away-uniqueness. If both have
outer edge $e$, use fibre uniqueness. The two mixed cases contradict the
cross-part disjointness condition. These cases exhaust all occurrences.
\end{proof}
\lean{Graph/DeepInventory}{19}{43}{Linear.uniqueHoles_partition}

\begin{lemma}[Inventories of a box and of a whole graph]
If the label of $e$ is a box with body $B$, then $E_G(e,h)$ holds exactly when
$h\in\mathcal H(B)$. For an arbitrary edge $e$, a label $h$ occurs in $G$ exactly
when it occurs away from $e$ or satisfies $E_G(e,h)$.
\end{lemma}
\begin{proof}
The shallow position of a box is not a hole, while each of its child positions
is precisely an occurrence of its body, with the same observation. For the
whole-graph statement, write an occurrence as an outer edge and a fibre
position, and split according to whether that edge equals $e$. In the reverse
direction either witness is already an occurrence of $G$.
\end{proof}
\lean{Graph/DeepInventory}{46}{56}{Linear.edgeHasHole_box}
\lean{Graph/DeepInventory}{59}{69}{Linear.hasHole_partition}

\begin{lemma}[Exact inventory after replacing a box body]
Changing the body of box edge $e$ to $B'$ preserves the inventory away from $e$
exactly. The entire new inventory is the union of the old inventory away from
$e$ and $\mathcal H(B')$.
\end{lemma}
\begin{proof}
Every other edge has exactly the same label before and after the change.
Transport its fibre positions along this label equality; fibre transport
preserves observations in both directions. This proves the first claim.
Partition the new graph at $e$ and use the preceding box-inventory lemma to
identify its selected fibre with the new body's inventory.
\end{proof}
\lean{Graph/DeepInventory}{72}{84}{Linear.replaceBox_retained_iff}
\lean{Graph/DeepInventory}{87}{92}{Linear.replaceBox_hasHole_iff}

\begin{lemma}[Box uniqueness]
The hole observations in a box fibre are unique exactly when its body has
unique holes. Consequently, the body of every box in a context has unique holes.
\end{lemma}
\begin{proof}
A box's shallow position cannot witness a hole. Two child positions observing
the same hole are equal precisely when their body occurrences are equal.
Injectivity of the child constructor proves the reverse implication. For the
consequence, restrict the graph's uniqueness to the selected fibre using the
partition lemma, then apply this equivalence.
\end{proof}
\lean{Graph/DeepInventory}{95}{109}{Linear.boxFiber_unique_iff}
\lean{Graph/DeepInventory}{112}{117}{Linear.uniqueHoles_box_body}

\begin{lemma}[Uniqueness after changing one box]
Suppose $G$ is a context. Replacing the body of box edge $e$ by $B'$ yields a
context if and only if $B'$ has unique holes and its hole inventory is disjoint
from the inventory of $G$ away from $e$.
\end{lemma}
\begin{proof}
First the unchanged part remains unique: transport two of its occurrences to
$G$ along their unchanged label equalities, apply uniqueness in $G$, and use
injectivity of fibre transport to recover equality in the new graph.
The uniqueness partition for the changed graph now has a proved first
condition. Its second condition is equivalent to uniqueness of $B'$, by box
uniqueness. Its third condition is exactly the stated disjointness, by the
preservation of the outside inventory and the exact inventory of a box.
\end{proof}
\lean{Graph/DeepInventory}{120}{133}{Linear.replaceBox_awayUnique}
\lean{Graph/DeepInventory}{137}{157}{Linear.replaceBox_uniqueHoles_iff}

\subsection{A syntactic inventory of retained holes}

\begin{definition}[Retention along a site]\label{term:site-retained-hole}
Let $R_s(h)$ be defined recursively on the site address. At a shallow site
selecting $e$ in $G$, it is $\mathsf{RetainedHole}(G,e,h)$: there is an occurrence
of $h$ under another outer edge. If the site descends through box edge $e$ to an
inner site $t$, set
\[
 R_s(h)\quad\Longleftrightarrow\quad
 \mathsf{RetainedHole}(G,e,h)\ \lor\ R_t(h).
\]
This is $\mathsf{SiteRetainedHole}$. It depends only on the host and address,
not on the replacement graph or a provenance certificate. It excludes all
occurrences below the edge finally removed, while retaining the siblings of
every enclosing box. Enclosing box edges are not themselves holes.
\end{definition}
\lean{Graph/DeepInventory}{161}{165}{Linear.SiteRetainedHole}

\begin{theorem}[Exact deep hole inventory]
For every site and every type-correct replacement, without uniqueness or
freshness assumptions,
\[
 h\in\mathcal H(G[s:=H])\quad\Longleftrightarrow\quad
 R_s(h)\ \lor\ h\in\mathcal H(H).
\]
\end{theorem}
\lean{Graph/DeepInventory}{168}{176}{Linear.plugAt_hasHole_iff}
\begin{proof}
Induct on the site address. The shallow case is the exact shallow inventory
theorem. For a site inside a box, the exact box-replacement inventory is the
union of the unaffected outer inventory and the changed body's inventory.
The induction hypothesis identifies the latter as the inner retained inventory
union the replacement inventory. Reassociate these three alternatives; the
first two are exactly $R_s$. Both directions are preserved at each step, so no
retained or inserted label is lost from the characterization.
\end{proof}

\begin{definition}[Exact deep freshness]\label{term:site-fresh}
The predicate $\mathsf{SiteFresh}(s,H)$ says
\[
 \forall h,\quad R_s(h)\ \Longrightarrow\ h\notin\mathcal H(H).
\]
Only retained host labels constrain the replacement; a label that is wholly
consumed is available for reuse.
\end{definition}
\lean{Graph/DeepInventory}{179}{181}{Linear.SiteFresh}

\begin{lemma}[Retention implies original presence]
If $R_s(h)$, then $h\in\mathcal H(G)$.
\end{lemma}
\begin{proof}
Induct on the site. An outside witness is already a host occurrence. For an
inner retained witness, the induction hypothesis supplies an occurrence in the
original box body; the box-inventory lemma and whole-graph partition lift it
back into the host.
\end{proof}
\lean{Graph/DeepInventory}{184}{193}{Linear.siteRetainedHole_hasHole}

\begin{lemma}[Replacement uniqueness is necessary for shallow plugging]
If a shallow plugged graph has unique holes, its replacement has unique holes.
\end{lemma}
\begin{proof}
The right-tagged injection embeds every replacement edge into the result,
preserving labels and all fibre observations. Uniqueness reflects along this
injective edge map by the established embedding theorem.
\end{proof}
\lean{Graph/DeepInventory}{196}{200}{Linear.plug_uniqueHoles_replacement}

\subsection{Necessary and sufficient context closure}

\begin{theorem}[Exact context criterion at arbitrary depth]
Suppose $G$ has unique holes. Then
\[
 \mathsf{UniqueHoles}(G[s:=H])
 \quad\Longleftrightarrow\quad
 \mathsf{UniqueHoles}(H)\ \land\ \mathsf{SiteFresh}(s,H).
\]
Thus uniqueness of the replacement is a necessary conclusion in the forward
direction, not an additional assumption of the theorem.
\end{theorem}
\begin{proof}
Induct on the address. At a shallow site, uniqueness of the result forces both
replacement uniqueness and sharp shallow freshness. The reverse implication is
the proved shallow closure theorem.

For an inside site, let $B$ be the old body and $t$ the inner address. Since $G$
is a context, $B$ is a context. The box-replacement criterion says that the
result is a context exactly when $B[t:=H]$ is a context and its inventory is
disjoint from the outer inventory away from the selected box. The induction
hypothesis replaces its first condition by uniqueness of $H$ and inner
freshness.

For necessity, inner freshness excludes collisions with inner retained labels.
A label inserted from $H$ is present in the changed body by the exact inventory
theorem; outer disjointness therefore excludes collisions with every retained
outer label. These are precisely the two alternatives defining $R_s$.

For sufficiency, site freshness supplies inner freshness, so induction proves
uniqueness of the changed body. Consider a label shared by that body and an
unaffected outer edge. The exact inventory theorem gives two cases. If the
label comes from $H$, site freshness is contradicted. If it is retained inside
the body, retention implies its presence in the old body. It would then have
occurred both inside and outside that box already in $G$, contradicting the
uniqueness partition of $G$. This proves the required outer disjointness and
completes the induction.
\end{proof}
\lean{Graph/DeepInventory}{204}{234}{Linear.plugAt_uniqueHoles_iff}

If $H$ is already a context, its uniqueness conjunct is known, leaving the
necessary-and-sufficient condition $\mathsf{SiteFresh}(s,H)$ alone.
\lean{Graph/DeepInventory}{237}{242}{Linear.plugAt_uniqueHoles_iff_fresh}
\termdef{deep-certified-fresh}
The constructor $\mathsf{Context.plugAtFresh}$ pairs the deep replacement graph
with the proof supplied by this equivalence. Unlike full-label disjointness,
its required certificate permits reuse of consumed labels at any box depth.
\lean{Graph/DeepInventory}{245}{248}{Linear.Context.plugAtFresh}

\begin{theorem}[Exact hole-freeness criterion]
The result $G[s:=H]$ is hole-free if and only if no label satisfies $R_s$ and $H$
is hole-free. No uniqueness assumption is needed.
\end{theorem}
\begin{proof}
The exact inventory is empty precisely when both of its constituent inventories
are empty. In either direction, an alleged witness in one part gives, or comes
from, a witness in the union.
\end{proof}
\lean{Graph/DeepInventory}{252}{260}{Linear.plugAt_holeFree_iff}

\subsection{Consumption of the selected label}

\begin{definition}[Selected hole label]\label{term:selected-hole}
The partial observation $\mathsf{Site.selectedHole}(s)$ follows the address to
its selected edge and reads that edge's own shallow hole label. It does not
inspect descendants of the selected edge. Thus selecting a box yields no hole
label, even if its body contains holes.
\end{definition}
\lean{Graph/DeepInventory}{263}{267}{Linear.Site.selectedHole}

\begin{lemma}[Presence and non-retention of a selected hole]
If the selected label is $h$, then $h$ occurs in the original host. If the host
is a context, $R_s(h)$ is false.
\end{lemma}
\begin{proof}
For presence, the shallow case uses the selected edge's shallow fibre position.
The inside case lifts the induction hypothesis from the body through the box.
For non-retention, the shallow statement is the existing consumed-label lemma.
Inside a box, a retained witness outside the box contradicts the old graph's
uniqueness partition, because the selected label occurs inside its body. An
inner retained witness contradicts induction applied to that body's uniqueness.
\end{proof}
\lean{Graph/DeepInventory}{270}{276}{Linear.Site.selectedHole_present}
\lean{Graph/DeepInventory}{280}{291}{Linear.site_selected_not_retained}

\begin{theorem}[Exact consumption at arbitrary depth]
If $G$ is a context and $s$ selects a hole labelled $h$, then
\[
 h\in\mathcal H(G[s:=H])\quad\Longleftrightarrow\quad h\in\mathcal H(H).
\]
\end{theorem}
\begin{proof}
Use the exact deep inventory theorem and eliminate its retained-host alternative
by non-retention of the selected label. This permits both genuine removal and
intentional reintroduction, depending on the replacement.
\end{proof}
\lean{Graph/DeepInventory}{294}{299}{Linear.plugAt_consumed_hole_iff}

\subsection{Concrete regressions inside boxes}

In the existing one-hole boxed graph, no label is retained when the inner hole
is selected. At the outer level there is no other edge (the carrier has one
element), and inside the box there is no other edge either. This proves
$\neg R_s(h)$ for every $h$.
\lean{Graph/DeepInventoryExamples}{9}{14}{Linear.Examples.boxedHole_retained_absent}
Replacing that hole by the same one-hole context therefore preserves uniqueness,
by sharp deep freshness. Full disjointness fails because both original graphs
contain the selected label. The consumed-label theorem also proves that this
label is genuinely present after replacement, witnessed by its occurrence in
the inserted graph.
\lean{Graph/DeepInventoryExamples}{17}{19}{Linear.Examples.deep_reuse_unique}
\lean{Graph/DeepInventoryExamples}{21}{23}{Linear.Examples.deep_reuse_not_disjoint}
\lean{Graph/DeepInventoryExamples}{26}{27}{Linear.Examples.deep_reuse_present}
Replacing instead with the hole-free value graph yields a hole-free result:
combine absence of retained labels and hole-freeness of the value using the
exact hole-freeness criterion. This is a shorter independent derivation of the
previous global filling result.
\lean{Graph/DeepInventoryExamples}{30}{31}{Linear.Examples.deep_fill_holeFree}

For a negative example, box the earlier two-hole chain with zero bound variables
and one free output. Its external interface is one zero-binding thunk input and
one term output. Define a site descending into this box and selecting the first
hole. The boxed graph is a context because the two-hole body is a context and
boxing preserves uniqueness.
\lean{Graph/DeepInventoryExamples}{34}{35}{Linear.Examples.boxedTwoHoles}
\lean{Graph/DeepInventoryExamples}{37}{38}{Linear.Examples.boxedFirstSite}
\lean{Graph/DeepInventoryExamples}{40}{41}{Linear.Examples.boxedTwoHoles_unique}
The one-hole replacement carrying label~$1$ is also a context, but it collides
with the retained second hole. These two witnesses refute deep freshness. The
necessary direction of the deep context theorem then proves that the result
fails uniqueness.
\lean{Graph/DeepInventoryExamples}{44}{46}{Linear.Examples.deep_collision_not_fresh}
\lean{Graph/DeepInventoryExamples}{48}{53}{Linear.Examples.deep_collision_breaks_unique}
Finally, the exact inventory distinguishes the labels: label~$0$ is absent and
label~$1$ is present. The first claim uses consumption and the fact that the
replacement's label has name~$1$, not~$0$; the second uses its inserted witness.
An inventory records presence rather than multiplicity, so this is consistent
with the independently proved duplicate occurrence of label~$1$.
\lean{Graph/DeepInventoryExamples}{56}{65}{Linear.Examples.deep_collision_inventory}

\addtocontents{toc}{\protect\newpage}
\Needspace{18\baselineskip}
\section{Recursive isomorphism and substitution modulo renaming}
\label{sec:recursive-isomorphism}

Literal equality of edge labels is too restrictive when a label contains a
whole graph. Renaming the body of a box changes its graph record, although its
mathematical incidence structure is unchanged. This section defines recursive
isomorphism, proves its invariants, and proves congruence of replacement at any
addressed depth. All comparisons are within a fixed depth bound and fixed
ordered outer interface. The quotient construction below does not identify
different depth bounds or prove the representation theorem for the paper's
set-based presentation.

\subsection{Isomorphisms with related labels}

\begin{definition}[Isomorphism relative to a label relation]\label{term:relational-isomorphism}
Fix a signature and a relation $R$ on its labels. An $R$-isomorphism from $G$ to
$H$ consists of bijections $v:G_V\to H_V$ of internal vertices and
$e:G_E\to H_E$ of edges. Internal vertex labels agree, corresponding edge labels
are related by $R$, and mapping every source or target list through the induced
port map gives the corresponding list in $H$. The port map fixes each ordered
input and output position. We write $G\cong_R H$ for the existence of such data.
\end{definition}
\lean{Graph/RelIso}{10}{16}{Net.RelIso}

\begin{lemma}[Identity, inverse, and composition]
If $R$ is reflexive, every graph has an identity $R$-isomorphism. If $R$ is
symmetric, every $R$-isomorphism has an inverse. If $R$ is transitive,
$R$-isomorphisms compose.
\end{lemma}
\begin{proof}
For identity use identity carrier maps and reflexivity of $R$. For inverse use
the inverse carrier bijections. Evaluate the original label equation at the
inverse image, cancel the edge bijection, and use symmetry of $R$. For incidence,
map the original incidence equation through the inverse port map; cancellation
recovers the reversed equation. Composition uses composed carrier bijections,
transitivity of the label relation, and the fact that a list mapped twice is
mapped by the composite. Port maps compose because they act on internal names
and fix both kinds of boundary position.
\end{proof}
\lean{Graph/RelIso}{22}{28}{Net.RelIso.refl}
\lean{Graph/RelIso}{30}{44}{Net.RelIso.symm}
\lean{Graph/RelIso}{46}{62}{Net.RelIso.trans}

\begin{lemma}[Incidence paths and ordinary isomorphisms]
Every $R$-isomorphism preserves and reflects adjacency, hence preserves and
reflects paths under its edge bijection. Every ordinary label-preserving
isomorphism gives an $R$-isomorphism when $R$ is reflexive.
\end{lemma}
\begin{proof}
Map a shared vertex to obtain forward adjacency. For reflection, source and
target preservation express the shared target vertex as images of two original
ports. Injectivity of the port bijection makes those ports equal, giving the
original shared vertex. Apply the existing list-path mapping lemma to this
adjacency equivalence. For an ordinary isomorphism, retain all carrier and
incidence data; equality of its labels and reflexivity give the required
$R$-relation at each edge.
\end{proof}
\lean{Graph/RelIso}{64}{78}{Net.RelIso.adjacent_iff}
\lean{Graph/RelIso}{80}{82}{Net.RelIso.path_iff}
\lean{Graph/RelIso}{86}{93}{Net.RelIso.ofIso}

\subsection{Recursive comparison of box contents}

\begin{definition}[Related boxes]\label{term:related-boxes}
For a relation on bodies of each fixed interface, two boxes are related when
they have the same number of bound variables, the same number of free outputs,
and related bodies at the interface specified by those two numbers. The
constructor $\mathsf{BoxRelated.bodies}$ enforces both metadata equalities; it
does not merely require equality of their sum.
\end{definition}
\lean{Graph/RecursiveIso}{10}{13}{Linear.BoxRelated}

\begin{definition}[Recursive label equivalence]\label{term:recursive-label-equivalence}
At depth zero, label equivalence is literal equality. At positive depth,
primitive labels are equivalent exactly when they are equal; a primitive and a
box are never equivalent. Two box labels are equivalent when their metadata
agree and there exists an isomorphism of their bodies whose corresponding
labels are recursively equivalent at the smaller depth. This is a well-founded
definition because the depth decreases whenever a box body is compared.
\end{definition}
\lean{Graph/RecursiveIso}{17}{23}{Linear.LabelEquiv}

\begin{definition}[Recursive hypernet isomorphism]\label{term:deep-isomorphism}
A recursive isomorphism $G\cong_d H$, represented by $\mathsf{DeepIso}$, is an
isomorphism relative to recursive label equivalence at depth $d$. Its outer
carrier bijections are explicit data. The label proof at a box supplies the
existence of a further body isomorphism, independently of the names chosen at
any other box. Every level preserves its own ordered boundary positions.
\end{definition}
\lean{Graph/RecursiveIso}{27}{27}{Linear.DeepIso}

\begin{lemma}[Equivalence laws for recursive labels]
Recursive label equivalence is reflexive, symmetric, and transitive.
\end{lemma}
\begin{proof}
Induct on depth. At depth zero the claims are equality laws. At positive depth,
primitive cases again reduce to equality and mixed primitive/box cases are
impossible. For box reflexivity, the induction hypothesis supplies reflexivity
of the body-label relation, so the identity relative isomorphism witnesses body
equivalence. For box symmetry, take a witnessing body isomorphism and invert it
using symmetry at the smaller depth. For transitivity, matching the two
box-relation constructors identifies all metadata, so the two body isomorphisms
have a common middle graph and interface. Compose them using transitivity at
the smaller depth. Package each resulting body isomorphism in the box relation.
\end{proof}
\lean{Graph/RecursiveIso}{30}{39}{Linear.labelEquiv_refl}
\lean{Graph/RecursiveIso}{42}{57}{Linear.labelEquiv_symm}
\lean{Graph/RecursiveIso}{60}{82}{Linear.labelEquiv_trans}
Identity, inverse, and composition of recursive graph isomorphisms follow by
instantiating the generic constructions with these three laws. Ordinary
isomorphisms embed into recursive ones by reflexivity of the label relation.
\lean{Graph/RecursiveIso}{87}{87}{Linear.DeepIso.refl}
\lean{Graph/RecursiveIso}{89}{89}{Linear.DeepIso.symm}
\lean{Graph/RecursiveIso}{91}{92}{Linear.DeepIso.trans}
\lean{Graph/RecursiveIso}{95}{95}{Linear.DeepIso.ofIso}

\begin{lemma}[Signatures and operation labels are invariant]
Equivalent labels have equal ordered input signatures, equal ordered output
signatures, and the same operation-label predicate.
\end{lemma}
\begin{proof}
At depth zero or for two primitives, substitute the label equality. Mixed
labels are impossible. For two boxes, the relation constructor identifies the
bound/free metadata, and both signatures depend only on those numbers. Both
box labels fail the operation predicate. No information about body carrier
names is used in these conclusions.
\end{proof}
\lean{Graph/RecursiveIso}{100}{112}{Linear.labelEquiv_input}
\lean{Graph/RecursiveIso}{115}{127}{Linear.labelEquiv_output}
\lean{Graph/RecursiveIso}{130}{142}{Linear.labelEquiv_isOperation}

\Needspace{16\baselineskip}
\subsection{Operation traces and paths}

\begin{theorem}[Recursive renaming preserves and reflects operation paths]
Let $f:G\cong_d H$, with port map $v$ and edge bijection $e$. An operation trace
from $x$ to $y$ along the edge list $\ell$ in $G$ exists exactly when the trace
from $v(x)$ to $v(y)$ along $\mathsf{map}(e,\ell)$ exists in $H$. The analogous
statement holds for forward operation paths.
\end{theorem}
\begin{proof}
First prove forward trace transport by induction on the trace. An empty trace
maps to an empty trace. For a nonempty trace, operation-label invariance supplies
the required label predicate. Source and target preservation map its two
incidence witnesses, and induction supplies the remaining trace.

For reflection transport the target trace through the inverse recursive
isomorphism. The inverse port and edge maps cancel the original maps, recovering
both original endpoints and exactly the original edge list. A forward operation
path is the reverse-oriented trace with its edge list reversed. List mapping
commutes with reversal, so the trace equivalence gives the path equivalence.
In particular, input and output endpoints retain their boundary indices.
\end{proof}
\lean{Graph/RecursiveIso}{145}{156}{Linear.DeepIso.operationTrace}
\lean{Graph/RecursiveIso}{159}{171}{Linear.DeepIso.operationTrace_iff}
\lean{Graph/RecursiveIso}{174}{181}{Linear.DeepIso.operationPath_iff}

\subsection{Transporting occurrences and holes through all boxes}

\begin{definition}[Fibre and hole embeddings]\label{term:hole-embedding}
A fibre embedding maps positions in one label to positions in another,
injectively, and preserves their complete hole observations, including the
absence of a hole. A hole embedding similarly maps graph occurrences
injectively and preserves their observations. These maps concern occurrence
addresses; their definition does not itself assert any incidence property.
\end{definition}
\lean{Graph/RecursiveIsoHoles}{9}{12}{Linear.FiberEmbedding}
\lean{Graph/RecursiveIsoHoles}{15}{18}{Linear.HoleEmbedding}

\begin{lemma}[Assembling fibre embeddings]
An injective map of outer edges and a fibre embedding over each mapped edge
together give a hole embedding of graphs.
\end{lemma}
\begin{proof}
Map $(a,x)$ to the mapped edge paired with the image of $x$ under its fibre
embedding. If two results coincide, projecting their edges and using edge-map
injectivity identifies the original edges. Their remaining fibre positions now
have the same type, and fibre-map injectivity identifies them. Observation
preservation is exactly the preservation property of the corresponding fibre
embedding.
\end{proof}
\lean{Graph/RecursiveIsoHoles}{21}{32}{Linear.HoleEmbedding.ofFibers}

\termdef{all-edge-transport}
A hole embedding also transports the original deep addresses: first use the
occurrence-coordinate bijection of the source, apply the hole embedding, then
use the inverse coordinate bijection of the target. This defines
$\mathsf{HoleEmbedding.allEdge}$.
\lean{Graph/RecursiveIsoHoles}{35}{37}{Linear.HoleEmbedding.allEdge}
The resulting address map is injective. Apply the target coordinate bijection
to an equality of images, cancel its inverse, use embedding injectivity, and
finally use source-coordinate injectivity. It preserves hole lookup because
both coordinate bijections preserve observation and the middle embedding does
so as well.
\lean{Graph/RecursiveIsoHoles}{39}{45}{Linear.HoleEmbedding.allEdge_injective}
\lean{Graph/RecursiveIsoHoles}{47}{53}{Linear.HoleEmbedding.allEdge_hole}

\begin{lemma}[Recursive labels admit fibre embeddings]
If two labels are recursively equivalent, an injective, observation-preserving
map of their fibres exists.
\end{lemma}
\begin{proof}
Induct on depth. Equal primitives use the identity map. Mixed label cases are
impossible. For equivalent boxes, choose a witnessing isomorphism of their
bodies. For each body edge, the induction hypothesis supplies a fibre embedding
from its recursive label-equivalence proof; choose one such map per edge.
Assemble these maps over the body isomorphism's injective edge map to obtain a
body hole embedding, and express it in original deep-address coordinates.

Map the box's shallow position to the shallow position, and each child address
through this body map. The two alternatives remain disjoint. Within the shallow
alternative equality is immediate; within the child alternative injectivity is
the proved injectivity of body-address transport. Shallow observations are empty
on both sides, while child observations agree by the body's transport theorem.
This constructs the desired embedding. Classical choice selects witnesses of
the existential body-isomorphism and embedding statements; it does not supply
any graph invariant as an additional assumption.
\end{proof}
\lean{Graph/RecursiveIsoHoles}{57}{96}{Linear.labelEquiv_hasEmbedding}
The noncomputable function $\mathsf{labelEquiv\_embedding}$ chooses the embedding
whose existence was just proved. Applying it at every outer edge of a recursive
isomorphism, and assembling over its edge bijection, defines
$\mathsf{DeepIso.holeEmbedding}$.
\lean{Graph/RecursiveIsoHoles}{99}{101}{Linear.labelEquiv_embedding}
\lean{Graph/RecursiveIsoHoles}{104}{107}{Linear.DeepIso.holeEmbedding}

\begin{lemma}[Presence transports and uniqueness reflects]
A hole embedding transports every present hole label from source to target. If
the target has unique hole observations, so does the source.
\end{lemma}
\begin{proof}
Map the occurrence witnessing a present label and use preservation of lookup.
For uniqueness, map two source occurrences with equal hole observations. Target
uniqueness makes their images equal, and embedding injectivity makes the
original occurrences equal.
\end{proof}
\lean{Graph/RecursiveIsoHoles}{110}{113}{Linear.HoleEmbedding.hasHole}
\lean{Graph/RecursiveIsoHoles}{116}{122}{Linear.HoleEmbedding.unique}

\begin{theorem}[Hole and context invariance under recursive isomorphism]
Recursively isomorphic graphs have identical hole inventories. One has unique
holes if and only if the other does; one is hole-free if and only if the other
is.
\end{theorem}
\begin{proof}
Construct hole embeddings from the isomorphism and its inverse. Presence
transport gives the two implications for inventory equality. Uniqueness
reflection, in the two directions, gives the context equivalence. Negating hole
presence for every label gives the hole-freeness equivalence. The two chosen
embeddings are not asserted to be mutually inverse; the two implications and
injectivity are sufficient for these results.
\end{proof}
\lean{Graph/RecursiveIsoHoles}{125}{127}{Linear.DeepIso.hasHole_iff}
\lean{Graph/RecursiveIsoHoles}{130}{132}{Linear.DeepIso.uniqueHoles_iff}
\lean{Graph/RecursiveIsoHoles}{134}{137}{Linear.DeepIso.holeFree_iff}

\subsection{Congruence of shallow and deep replacement}

\termdef{relational-plugging}
For a fixed host $C$ and shallow selected edge $h$, an $R$-isomorphism
$f:H\cong_R K$ gives the candidate carrier bijections for the plugged graphs:
use identity on the host summand and $f$ on the replacement summand. For edges,
the host summand is the subtype of edges different from $h$.
\lean{Graph/RelIsoPlug}{16}{17}{Net.RelIsoPlug.vertices}
\lean{Graph/RelIsoPlug}{19}{20}{Net.RelIsoPlug.edges}
The induced port map commutes with host embedding, by a case split on internal,
input, and output ports. It commutes with insertion after the replacement's own
port map: an internal port uses the right summand, and a boundary port reduces
to the host-embedding equation at the selected source or target port.
\lean{Graph/RelIsoPlug}{22}{24}{Net.RelIsoPlug.map_host}
\lean{Graph/RelIsoPlug}{26}{32}{Net.RelIsoPlug.map_insert}

\begin{theorem}[Shallow congruence for related labels]
For reflexive $R$, an $R$-isomorphism $H\cong_R K$ induces
$C[h:=H]\cong_R C[h:=K]$. No non-self-loop condition is required.
\end{theorem}
\begin{proof}
Use the preceding sum bijections. Host internal labels are unchanged and
replacement internal labels are preserved by $f$. Host edge labels satisfy
reflexivity of $R$, and inserted edge labels use the relation supplied by $f$.
For a host incidence list, commute the port map with host embedding pointwise.
For an inserted incidence list, commute it with insertion and then use the
source or target preservation equation of $f$. This establishes every field of
the resulting relative isomorphism without requiring insertion to be globally
injective.
\end{proof}
\lean{Graph/RelIsoPlug}{38}{85}{Net.RelIso.plug}

\begin{lemma}[Box congruence]
Boxing recursively isomorphic bodies with the same metadata gives recursively
isomorphic boxes. Replacing one fixed box edge in a surrounding graph by either
of those bodies also gives recursively isomorphic graphs.
\end{lemma}
\begin{proof}
For a single boxed graph, use identity maps on its empty internal carrier and
one-edge carrier. Its labels are related by the body-isomorphism witness, and
its boundary lists depend only on the common metadata. In a surrounding graph,
use identity maps on both outer carriers. At the selected edge use the related
box-body witness; at every other edge use reflexivity of recursive label
equivalence. The two outer incidence structures are identical, so their
preservation equations reduce to mapping a list by identity.
\end{proof}
\lean{Graph/RecursiveIsoPlug}{10}{19}{Linear.DeepIso.box}
\lean{Graph/RecursiveIsoPlug}{23}{42}{Linear.replaceBox_deepIso}

\begin{theorem}[Congruence at every addressed depth]
For a fixed host $G$ and typed site $s$, if $H\cong_r K$, then
\[
 G[s:=H]\cong_d G[s:=K].
\]
No context-uniqueness, freshness, or non-self-loop assumption is needed.
\end{theorem}
\begin{proof}
Induct on the site. A shallow site uses the generic shallow congruence with the
reflexive recursive-label relation. For a site inside a box, induction gives an
isomorphism between the two changed bodies. Lift it through the selected box
using box congruence. Repeating this step crosses every enclosing box in the
address.
\end{proof}
\lean{Graph/RecursiveIsoPlug}{46}{51}{Linear.plugAt_deepIso}

Two consequences connect this congruence with the preceding milestone. Exact
site freshness is invariant when the replacement is changed by recursive
isomorphism: the retained host predicate stays fixed, and the replacement hole
inventories coincide. Context validity of the two plugged results is also
equivalent, by context invariance applied to their constructed recursive
isomorphism. The latter statement does not assume the host is a context.
\lean{Graph/RecursiveIsoPlug}{54}{58}{Linear.siteFresh_deepIso_iff}
\lean{Graph/RecursiveIsoPlug}{61}{64}{Linear.plugAt_deepIso_unique_iff}

\subsection{Graphs and replacement modulo recursive renaming}

\begin{definition}[Recursive isomorphism classes]\label{term:recursive-quotient}
At fixed depth $d$ and interface $(I,O)$, declare $G\sim H$ when a recursive
isomorphism exists. Identity, inverse, and composition prove that this is an
equivalence relation, packaged as $\mathsf{deepSetoid}$. Its quotient is
$\mathsf{HypernetModuloIso}$, the type of graphs modulo independent carrier
renaming at all levels of their hierarchy. Write $[G]$ for the class of $G$.
\end{definition}
\lean{Graph/RecursiveIsoPlug}{67}{73}{Linear.deepSetoid}
\lean{Graph/RecursiveIsoPlug}{76}{77}{Linear.HypernetModuloIso}

\begin{definition}[Plugging a replacement class]\label{term:quotient-plugging}
For a fixed concrete host and site $s$, define
\[
 \mathsf{plugAtModuloIso}(s,[H])=[G[s:=H]].
\]
\end{definition}
\begin{proof}[Well-definedness]
If $H$ and $K$ represent the same source class, their recursive isomorphism
induces one between the plugged graphs by deep congruence. The results therefore
belong to the same target class. Quotient elimination defines the displayed
operation from this representative-independent function. Its value on a
concrete representative is the displayed equation by the computation rule of
quotient elimination.
\end{proof}
\lean{Graph/RecursiveIsoPlug}{81}{84}{Linear.plugAtModuloIso}
\lean{Graph/RecursiveIsoPlug}{87}{90}{Linear.plugAtModuloIso_mk}
\termdef{quotient-context-validity}
Context validity descends to classes as well: define
$\mathsf{moduloIsoUniqueHoles}([G])$ to mean that $G$ has unique holes. The
isomorphism-invariance theorem makes these propositions equivalent for any two
representatives; propositional extensionality supplies the equality required by
quotient elimination.
\lean{Graph/RecursiveIsoPlug}{93}{94}{Linear.moduloIsoUniqueHoles}

These constructions quotient the replacement argument for a fixed host and
address. Section~\ref{sec:site-transport} extends invariance to corresponding sites in
isomorphic hosts. Section~\ref{sec:composition} proves associativity and units.
Section~\ref{sec:deep-interchange} constructs the quotient of host/site pairs.

\subsection{A strict extension of literal-label isomorphism}

Use the earlier two-hole chain. The permutation $i\mapsto 1-i$ on its two-edge
carrier exchanges edge names. Applying it twice is identity because
$0\leq i<2$. Reindex the chain's edges by this bijection and keep its internal
vertex name unchanged; call the result $H'$. The existing reindexing isomorphism,
embedded into recursive isomorphism, witnesses $H\cong_0 H'$.
\lean{Graph/RecursiveIsoExamples}{9}{13}{Linear.Examples.swapTwo}
\lean{Graph/RecursiveIsoExamples}{15}{16}{Linear.Examples.renamedTwoHoles}
\lean{Graph/RecursiveIsoExamples}{18}{19}{Linear.Examples.twoHoles_deepIso}
The records $H$ and $H'$ are unequal. If they were equal, their edge-label
functions would agree at the edge named zero. Reading the hole name there gives
$0$ in $H$ and $1$ in $H'$, a contradiction.
\lean{Graph/RecursiveIsoExamples}{22}{29}{Linear.Examples.renamedTwoHoles_ne}

Box $H'$ with zero bound variables and one free output. Box congruence relates
this box to the earlier box of $H$.
\lean{Graph/RecursiveIsoExamples}{31}{32}{Linear.Examples.renamedBoxedTwoHoles}
\lean{Graph/RecursiveIsoExamples}{35}{36}{Linear.Examples.boxedTwoHoles_deepIso}
No literal label-preserving isomorphism relates these two boxed graphs. Such an
isomorphism's label equation at the sole source edge would make their box labels
equal. Equality of these labels forces equality of the body records, contrary
to the preceding result. This proves that recursive isomorphism identifies
strictly more graphs than the earlier literal-label notion.
\lean{Graph/RecursiveIsoExamples}{40}{46}{Linear.Examples.boxedTwoHoles_no_literalIso}
The renamed box remains a context by recursive invariance of unique holes. The
two different box records represent the same quotient class by quotient
soundness applied to their recursive isomorphism.
\lean{Graph/RecursiveIsoExamples}{48}{49}{Linear.Examples.renamedBoxedTwoHoles_unique}
\lean{Graph/RecursiveIsoExamples}{52}{55}{Linear.Examples.renamedBoxes_same_class}

For a nontrivial addressed replacement, use the earlier site selecting the first
hole inside a box. Substituting $H$ or $H'$ there gives recursively isomorphic
results by deep congruence. Finally, the primitive operations named
$\mathsf{value}$ and $\mathsf{otherValue}$ are not equivalent: at depth zero
recursive label equivalence is equality, and these are distinct constructors.
Thus carrier renaming does not erase semantic operation labels.
\lean{Graph/RecursiveIsoExamples}{58}{60}{Linear.Examples.deep_plug_renaming}
\lean{Graph/RecursiveIsoExamples}{63}{66}{Linear.Examples.distinct_operations_not_equivalent}

\Needspace{18\baselineskip}
\section{Transporting replacement sites through isomorphisms}
\label{sec:site-transport}

An address denotes an occurrence, not a permanent numerical edge name. If an
isomorphism exchanges two edges, the selected occurrence must move with it.
This section proves existence of that transport at arbitrary box depth,
congruence when both host and replacement change, and invariance of the exact
freshness condition. Depth bounds and ordered outer interfaces remain fixed.

\subsection{Transport of boundary indices and shallow hosts}

\begin{definition}[Boundary transport]\label{term:boundary-transport}
Given equalities $I=I'$ and $O=O'$ of ordered interfaces, transport a port by
fixing its internal name, or by retaining its numerical input or output index
and using the equality to prove the new bound. Transport a graph to the new
interface by keeping its internal vertices, edges, enumerations, and labels,
and mapping every incidence list through this port transport. This graph is
written $H^{I',O'}$ below.
\end{definition}
All graph invariants follow by eliminating the two interface equalities: the
port map is then the identity, and every invariant is exactly the original
one. Reflexive port transport is pointwise the identity, hence is the identity
function by function extensionality. Reflexive graph transport is equal to the
original graph by record equality and the identity law for list mapping.
Keeping carriers definitionally unchanged makes the construction useful even
before the interface equality proofs have been eliminated.
\lean{Graph/TransportPlug}{8}{12}{Port.cast}
\lean{Graph/TransportPlug}{14}{15}{Port.cast_refl}
\lean{Graph/TransportPlug}{17}{19}{Port.cast_refl_fun}
\lean{Graph/TransportPlug}{22}{39}{Net.reboundary}
\lean{Graph/TransportPlug}{41}{43}{Net.reboundary_refl}

\begin{definition}[Complement restriction]\label{term:complement-restriction}
For a bijection $f:A\to B$ and $a\in A$, restrict $f$ to a bijection
$A\setminus\{a\}\to B\setminus\{f(a)\}$. Its inverse is the restriction of
$f^{-1}$.
\end{definition}
Injectivity shows that a surviving source cannot map to $f(a)$. Conversely, if
$f^{-1}(b)=a$, the right inverse equation gives $b=f(a)$, contradicting survival.
Both inverse laws follow from those of $f$, since subtype elements with equal
underlying values are equal.
\lean{Graph/TransportPlug}{46}{51}{Bijection.except}

\begin{lemma}[Transport of attachment positions]
Let $f:C\cong_R D$ and let $e$ be a selected edge. If its corresponding input
interfaces are equal, mapping the $i$th source port of $e$ by $f$ gives the
$i$th source port of $f(e)$. The analogous assertion holds for targets when the
output interfaces are equal.
\end{lemma}
\begin{proof}
Apply optional list lookup at index $i$ to the source-list equation of $f$.
The typing invariant identifies the source-list length with the input-interface
length, so both lookups are in bounds. Lookup commutes with mapping; cancelling
the optional-result constructor gives the port equation. The interface equality
supplies the bound for the target index. Repeat for the target-list equation.
No disjointness between source and target lists is used.
\end{proof}
\lean{Graph/TransportPlug}{58}{65}{Net.RelIso.sourcePort}
\lean{Graph/TransportPlug}{68}{75}{Net.RelIso.targetPort}

\begin{theorem}[Shallow host congruence]
Suppose $f:C\cong_R D$, $R$ is reflexive, and corresponding selected edges have
equal input and output interfaces. For every replacement $A$ of the source
interface,
\[
 C[e:=A]\cong_R D[f(e):=A^{I',O'}].
\]
\end{theorem}
\begin{proof}
On internal vertices use the disjoint sum of the host vertex bijection and the
identity on replacement vertices. On edges use the complement restriction of
the host edge bijection and the identity on replacement edges. These are the
auxiliary maps \texttt{HostPlug.vertices} and \texttt{HostPlug.edges}.
The whole-graph port map commutes with host inclusion by case analysis on a
port. It also commutes with replacement insertion followed by boundary
transport: internal vertices are immediate, and input and output ports use the
attachment-position lemmas just proved.

For surviving edges, label relatedness is supplied by $f$; for inserted edges,
use reflexivity of $R$. Internal vertex labels are preserved by the host
isomorphism or by identity. For incidence of a surviving edge, rewrite with
$f$'s incidence equation, combine successive list maps, and apply commutation
with host inclusion. For an inserted edge, combine the insertion and whole-port
maps and use commutation with replacement insertion. The same argument proves
the target equation. Self-loops are allowed throughout.
\end{proof}
\lean{Graph/TransportPlug}{83}{84}{Net.HostPlug.vertices}
\lean{Graph/TransportPlug}{86}{88}{Net.HostPlug.edges}
\lean{Graph/TransportPlug}{90}{93}{Net.HostPlug.map_host}
\lean{Graph/TransportPlug}{95}{107}{Net.HostPlug.map_insert}
\lean{Graph/TransportPlug}{112}{161}{Net.RelIso.plugHost}

\Needspace{16\baselineskip}
\subsection{Structural correspondence of addresses}

\begin{definition}[Address interface transport]\label{term:site-interface-transport}
Given $s:\mathsf{Site}(G,r,J,K)$ and equalities $J=J'$, $K=K'$, transport its
interface indices to obtain an address of type $\mathsf{Site}(G,r,J',K')$.
The selected occurrence is unchanged. Plugging $A$ at this transported address
equals plugging the boundary-transported $A$ at the original address.
\end{definition}
The definition is equality elimination. For the plugging equation eliminate
the two equalities and use reflexive graph boundary transport; both sides then
become the same substitution.
\lean{Graph/SiteTransport}{10}{12}{Linear.Site.reinterface}
\lean{Graph/SiteTransport}{15}{18}{Linear.plugAt_reinterface}

\begin{lemma}[Corresponding boxes and their replacement]
If $f:G\cong_{d+1}H$ and edge $e$ of $G$ carries box $(n,m,B)$, then $f(e)$
carries a box $(n,m,B')$ with $B\cong_d B'$. Moreover, replacing these two box
bodies by any $A\cong_d A'$ gives recursively isomorphic whole graphs.
\end{lemma}
\begin{proof}
For existence, inspect the target label in the label-relatedness proof for
$e$. A primitive target contradicts the definition of recursive label
equivalence. For a box target, the constructor of the box relation forces
matching $n,m$ and supplies a body isomorphism. This reasoning constructs an
existential statement, so elimination of the propositional box relation is
legitimate.

For replacement retain the vertex and edge bijections of $f$. At $e$ the new
labels are related by the given isomorphism $A\cong_d A'$. At every other edge
injectivity ensures that its target is different from $f(e)$, so both labels
remain unchanged and their original relation applies. Neither replacement
changes incidence or internal vertex labels; the remaining equations are
exactly those of $f$.
\end{proof}
\lean{Graph/SiteTransport}{22}{36}{Linear.DeepIso.boxAt}
\lean{Graph/SiteTransport}{39}{65}{Linear.DeepIso.replaceBox}

\begin{definition}[Corresponding sites]\label{term:site-correspondence}
The proposition $s\simeq t$, implemented by \texttt{SiteIso}, is generated by
two structural rules. A shallow site selecting $e$ corresponds to the site
selecting $f(e)$ for any recursive host isomorphism $f$; transport the latter's
interface indices back along the signature equalities. For sites inside boxes,
use a recursive outer-host isomorphism, corresponding selected box edges with
equal binding/free metadata, and structurally corresponding inner sites.
Both sites have the same local depth and replacement interface.
\end{definition}
This is a relation on addresses themselves, defined before proving anything
about the results of plugging. An inner correspondence may choose any suitable
isomorphism between the box bodies; the outer relation stores existence of a
body isomorphism, not a distinguished naming map for every inner vertex.
\lean{Graph/SiteTransport}{69}{84}{Linear.SiteIso}

\begin{theorem}[Existence of transported sites]
For every recursive isomorphism $f:G\cong_d H$ and every site
$s:\mathsf{Site}(G,r,J,K)$, there exists
$t:\mathsf{Site}(H,r,J,K)$ with $s\simeq t$.
\end{theorem}
\begin{proof}
Induct on $s$. For a shallow site use the shallow correspondence constructor
and the proved equality of corresponding signatures. For an inside site,
extract the corresponding target box and a body isomorphism using the preceding
lemma. Apply the induction hypothesis to the inner site and that body
isomorphism, then assemble the outer correspondence with the inside constructor.
This follows a finite address down the hierarchy and therefore terminates.
\end{proof}
\lean{Graph/SiteTransport}{87}{94}{Linear.Site.transport_exists}

\begin{definition}[Chosen site transport]\label{term:site-transport}
Write $s^f$ for a site chosen from this existence theorem by classical choice.
It has the original local depth and interface in the target host, and satisfies
$s\simeq s^f$ by the specification theorem for choice.
\end{definition}
The construction is noncomputable. No uniqueness, definitional round-trip law,
or canonical choice of inner box isomorphisms is asserted.
\lean{Graph/SiteTransport}{97}{99}{Linear.Site.transport}
\lean{Graph/SiteTransport}{102}{104}{Linear.Site.transport_spec}

\subsection{Substitution and freshness under transport}

\begin{theorem}[Simultaneous host and replacement congruence]
If $s\simeq t$ and $A\cong_r B$, then
\[
 \mathsf{plugAt}(s,A)\cong_d\mathsf{plugAt}(t,B).
\]
\end{theorem}
\begin{proof}
First take the same replacement on both sides and induct on the structural
correspondence. In the shallow case rewrite plugging at a transported interface
using its equation, then apply shallow host congruence. Recursive label
equivalence supplies the signature equalities and reflexivity required there.
In the inside case the induction hypothesis supplies an isomorphism of the
replaced inner bodies. Lift it through the corresponding outer boxes using
box replacement congruence. This proves \texttt{SiteIso.plug}.
For distinct isomorphic replacements, compose that isomorphism with the
previously proved congruence in the replacement at the fixed target site $t$.
\end{proof}
\lean{Graph/SiteTransport}{107}{117}{Linear.SiteIso.plug}
\lean{Graph/SiteTransport}{120}{125}{Linear.SiteIso.plug_congr}

\begin{lemma}[Quotient substitution agrees after host transport]
For every replacement isomorphism class $[A]$,
$\mathsf{plugAtModuloIso}(s,[A])=
\mathsf{plugAtModuloIso}(s^f,[A])$.
\end{lemma}
\begin{proof}
Use quotient induction to choose a representative $A$. The specification of
$s^f$ gives a structural correspondence; substitution congruence gives a
recursive isomorphism between the two results. Quotient soundness identifies
their classes.
\end{proof}
This is invariance under transporting the host/site pair. A separate quotient
type of host/site pairs has not been constructed. Section~\ref{sec:composition}
proves the nested composition and unit laws.
\lean{Graph/SiteTransport}{128}{132}{Linear.plugAtModuloIso_transport}

\begin{lemma}[Invariance of retained holes]
A recursive host isomorphism preserves and reflects the holes outside a
selected outer edge. Structurally corresponding sites retain exactly the same
hole labels at every depth.
\end{lemma}
\begin{proof}
An outer retained occurrence consists of an edge $x\ne e$ and a hole-observing
position in its label. Map $x$ by the host edge bijection and the position by
the previously constructed label-fibre embedding. Injectivity preserves
$x\ne e$, and the embedding preserves its hole observation. For reflection,
use the inverse edge bijection and the inverse isomorphism's fibre embedding;
the right inverse equation ensures that the preimage survives.

Interface transport does not change retention: eliminate the interface
equalities. Now induct on structural site correspondence. The shallow case
uses this equality elimination and the outer retained-hole equivalence. The
inside case uses the definition of deep retention as a disjunction: either
the occurrence lies under a surviving outer edge, or it is retained by the
inner site. Apply the outer equivalence and the induction hypothesis to the
two disjuncts.
\end{proof}
\lean{Graph/SiteTransport}{135}{147}{Linear.DeepIso.retainedHole_iff}
\lean{Graph/SiteTransport}{150}{153}{Linear.SiteRetainedHole_reinterface}
\lean{Graph/SiteTransport}{156}{164}{Linear.SiteIso.retainedHole_iff}

\begin{lemma}[Invariance of freshness and result validity]
For $s\simeq t$ and $A\cong_r B$, exact freshness holds for $(s,A)$ if and
only if it holds for $(t,B)$. The two substituted results satisfy global
unique-hole validity together, whether or not either original host is valid.
\end{lemma}
\begin{proof}
Freshness forbids a label occurring both in the retained host and in the
replacement. Transport the first occurrence with retained-hole equivalence and
the second with the replacement isomorphism's hole-inventory equivalence. This
proves each direction. For result validity, simultaneous substitution
congruence supplies a recursive isomorphism; apply its preservation and
reflection of unique holes. No original-host validity hypothesis is needed
for this invariance statement.
\end{proof}
\lean{Graph/SiteTransport}{167}{173}{Linear.SiteIso.fresh_iff}
\lean{Graph/SiteTransport}{176}{181}{Linear.SiteIso.plug_unique_iff}

\subsection{A moved address and a wrong-address counterexample}

Use the two-hole chain and its edge swap from the preceding section. The first
hole, formerly edge zero, is edge one after renaming. Define the target deep
site to enter the renamed box and select edge one. The outer boxed isomorphism
and the shallow correspondence for the edge swap prove structural
correspondence with the original deep site.
\lean{Graph/SiteTransportExamples}{9}{11}{Linear.Examples.renamedFirstSite}
\lean{Graph/SiteTransportExamples}{14}{16}{Linear.Examples.firstSites_correspond}
Simultaneous congruence now allows both the host body and the inserted chain to
be renamed independently. The replacement consisting of hole one still
collides with the retained hole one: freshness invariance transports the
previously proved failure of freshness to the new site.
\lean{Graph/SiteTransportExamples}{19}{21}{Linear.Examples.simultaneous_deep_renaming}
\lean{Graph/SiteTransportExamples}{24}{26}{Linear.Examples.renamed_collision_not_fresh}

For contrast within this example, define the wrong target site by keeping the
old numerical name zero. It selects hole one instead of hole zero. After
inserting the one-hole replacement, the wrong result retains hole zero,
witnessed by the surviving edge one inside the renamed box. The exact deep
inventory theorem therefore proves that hole zero is present. In the original
substitution, hole zero is consumed and absent, as already proved. A recursive
isomorphism would preserve hole presence, contradicting these two facts.
Thus the wrong-address substitution is not recursively isomorphic to the
original substitution; moving the address is essential.
\lean{Graph/SiteTransportExamples}{29}{31}{Linear.Examples.wrongRenamedSite}
\lean{Graph/SiteTransportExamples}{35}{42}{Linear.Examples.wrong_address_changes_result}

\Needspace{18\baselineskip}
\section{Associativity, units, and exact context composition}
\label{sec:composition}

We now compose addressed substitutions, rather than merely compare renamed
versions of one substitution. If $s$ selects an occurrence in $G$, inserting
$H$ there creates a copy of every occurrence of $H$. A site $t$ in $H$ thus
induces a site in $\mathsf{plugAt}(s,H)$. The main theorem identifies substitution
at that induced site with first substituting within $H$ and then inserting the
result at $s$. The identification is a concrete recursive isomorphism; on
isomorphism classes it is equality. We also prove the corolla unit laws and
an exact context-validity criterion for the composite, including cases where
an intermediate collision is removed by the final substitution.

\subsection{Shallow associativity with explicit carrier bijections}

Let $C$ be a graph, $e$ a selected edge, $H$ a replacement of its interface,
$f$ an edge of $H$, and $A$ a replacement of the interface of $f$. Compare
\[
 L=(C[e:=H])[\mathsf{inr}(f):=A],\qquad
 R=C[e:=H[f:=A]].
\]
These have the same outer interface, but different carrier records.

\begin{definition}[Reassociation bijections]\label{term:reassociation}
The internal-vertex bijection from $L$ to $R$ is
\[
 (C_V+H_V)+A_V\longrightarrow C_V+(H_V+A_V),
\]
which keeps each original name and changes only its sum tags. Its three cases
send a host vertex to the host summand, a middle vertex to the middle summand,
and a final vertex to the final summand. The edge bijection does the same for
surviving host edges, surviving middle edges, and final inserted edges.
\end{definition}
More explicitly, the left surviving-edge carrier consists of edges of
$(C_E\setminus\{e\})+H_E$ other than $\mathsf{inr}(f)$. A surviving host
edge maps to the host summand on the right. An edge $\mathsf{inr}(y)$ maps to
the surviving-middle summand, with $y\ne f$ obtained by injectivity of the sum
constructor. Final edges map to final edges. The inverse maps reverse these
three cases; inverse equations are proved by cases, with proof irrelevance
identifying the subtype certificates.
\lean{Graph/PlugAssociativity}{13}{18}{Net.PlugAssoc.vertices}
\lean{Graph/PlugAssociativity}{21}{32}{Net.PlugAssoc.edges}

\begin{lemma}[Attachment and port commutation]
The source and target positions of the inserted edge $\mathsf{inr}(f)$ are the
images, under insertion into $C$, of the corresponding positions of $f$ in $H$.
The reassociation port map commutes with inclusion of original host ports,
inclusion of middle ports, and insertion of final replacement ports.
\end{lemma}
\begin{proof}
The attachment assertions unfold positional lookup in the inserted incidence
lists and use lookup through a list map. For host ports, case analysis on
internal, input, and output ports proves the equation directly. For middle
ports, internal vertices are immediate; an input or output is attached to a
host port, so use the host equation. For final ports, internal vertices are
again immediate. A boundary port attaches at $f$; rewrite that attachment by
the positional lookup assertion and apply the middle-port equation.
No argument assumes that an input attachment differs from an output attachment.
\end{proof}
\lean{Graph/PlugAssociativity}{35}{37}{Net.PlugAssoc.sourcePort_inserted}
\lean{Graph/PlugAssociativity}{40}{42}{Net.PlugAssoc.targetPort_inserted}
\lean{Graph/PlugAssociativity}{45}{48}{Net.PlugAssoc.map_host}
\lean{Graph/PlugAssociativity}{51}{58}{Net.PlugAssoc.map_middle}
\lean{Graph/PlugAssociativity}{61}{75}{Net.PlugAssoc.map_insert}

\begin{theorem}[Shallow associativity]
There is a literal-label graph isomorphism $L\cong R$.
\end{theorem}
\begin{proof}
Use the two reassociation bijections. Vertex and edge labels agree by the
three carrier cases. For sources, split an edge into surviving host, surviving
middle, or final replacement. In the host case the left source list maps the
original host list through two inclusions and reassociation; combine list maps
and apply host commutation. The middle case uses middle commutation, and the
final case uses final insertion commutation. These give exactly the three
right-hand source lists. The target proof is identical with target lists.
The isomorphism preserves literal labels and is valid for self-loops and cycles.
\end{proof}
\lean{Graph/PlugAssociativity}{80}{118}{Net.plug_assoc}

\subsection{Composition at arbitrary box depth}

\begin{lemma}[Overwrite and insertion of box bodies]
Replacing the same box body twice has exactly the effect of the last replacement.
Changing the body of an inserted box commutes, up to literal-label isomorphism,
with changing that body before shallow insertion.
\end{lemma}
\begin{proof}
For overwrite, the two graphs have the same carriers, incidence, and internal
labels. At the selected edge both final labels contain the last body. At any
other edge both conditional updates keep the original label. Extensionality of
the label function and proof irrelevance give equality of the graph records.

For commutation with insertion use identity carrier bijections. Surviving host
labels are unchanged. On inserted edges the two tests for being the selected
edge agree, because $\mathsf{inr}(x)=\mathsf{inr}(f)$ exactly when $x=f$.
Incidence is unchanged by box replacement. The insertion functions also agree:
internal ports have the same tags, and boundary ports attach to the same outer
host positions. List mapping then proves both incidence equations.
\end{proof}
\lean{Graph/DeepAssociativity}{10}{20}{Linear.replaceBox_overwrite}
\lean{Graph/DeepAssociativity}{23}{62}{Linear.plug_replaceBox_iso}

\begin{definition}[Composed site]\label{term:composed-site}
Let $s:\mathsf{Site}(G,r,J,K)$ and $t:\mathsf{Site}(H,q,U,V)$, where $H$ has
the replacement interface of $s$. Define
\[
 s\star t:\mathsf{Site}(\mathsf{plugAt}(s,H),q,U,V).
\]
If $s$ is shallow and $t$ is shallow, select the inserted edge
$\mathsf{inr}(f)$. If $s$ is shallow and $t$ enters a box, enter the corresponding
inserted box and retain $t$'s inner address. If $s$ enters a box, rebuild that
box with the first replacement already performed, then descend to the recursively
composed inner site.
\end{definition}
This is recursion on the outer address. The new enclosing box carries exactly
the body produced by the first substitution, so the box-label equation is the
selected branch of the update. Local depth and interface indices are preserved
by construction. Although the surrounding graph operations use classical
finite-subset constructions, no isomorphism or target address is selected by
an existence theorem here.
\lean{Graph/DeepAssociativity}{66}{77}{Linear.Site.comp}

\begin{theorem}[Deep associativity]
For every compatible $s,t,A$,
\[
 \mathsf{plugAt}(s\star t,A)\cong_d
 \mathsf{plugAt}(s,\mathsf{plugAt}(t,A)).
\]
No uniqueness, freshness, or non-self-loop premise is required.
\end{theorem}
\begin{proof}
Induct on $s$. When $s$ and $t$ are shallow use the explicit shallow associator
and embed its literal-label isomorphism into recursive isomorphism. When $s$
is shallow and $t$ enters a box, use the box/insertion commutation lemma.
When $s$ enters a box, the left expression first updates that box to contain
$\mathsf{plugAt}(s_{\rm inner},H)$ and later updates it again after following
the composed inner address. Overwrite reduces these two updates to one. The
induction hypothesis compares its final body to the body obtained by performing
the inner substitution before the outer one. Box replacement congruence lifts
that body isomorphism through the enclosing box.
\end{proof}
\lean{Graph/DeepAssociativity}{80}{94}{Linear.plugAt_assoc}

\begin{lemma}[Quotient composition and three-stage coherence]
On recursive isomorphism classes,
\[
 P_{s\star t}([A])=P_s(P_t([A])),\qquad
 P_s=\mathsf{plugAtModuloIso}(s,-).
\]
For a third compatible site $u$, the two functions associated with
$(s\star t)\star u$ and $s\star(t\star u)$ also agree on every replacement class.
\end{lemma}
\begin{proof}
For the binary equation use quotient induction and quotient soundness applied
to deep associativity on representatives. For three stages, apply the binary
equation twice to each side; both reduce to $P_s(P_t(P_u([A])))$.
The ambient concrete hosts of the two composite addresses can differ, so this
is a statement about the induced quotient results, not literal equality of
address records in different hosts.
\end{proof}
\lean{Graph/DeepAssociativity}{97}{102}{Linear.plugAtModuloIso_comp}
\lean{Graph/DeepAssociativity}{205}{210}{Linear.plugAtModuloIso_comp_three}

The selected-hole observation of $s\star t$ equals that of $t$. Induct on $s$;
for shallow $s$, either form of $t$ retains the same selected label, and for
inside $s$ the definition discards the enclosing box and the induction hypothesis
applies.
\lean{Graph/DeepAssociativity}{105}{110}{Linear.Site.selectedHole_comp}

\subsection{Exact inventories and context validity of composites}

Write $R_s(h)$ for $\mathsf{SiteRetainedHole}(s,h)$ and $H_A(h)$ for
$\mathsf{HasHole}(A,h)$.

\begin{lemma}[Retained inventories compose by union]
Removing an inserted shallow edge $\mathsf{inr}(f)$ retains exactly the holes
under surviving original-host edges and surviving middle-graph edges. More
generally,
\[
 R_{s\star t}(h)\quad\Longleftrightarrow\quad R_s(h)\lor R_t(h).
\]
Consequently the final hole inventory is exactly
$R_s(h)\lor R_t(h)\lor H_A(h)$.
\end{lemma}
\begin{proof}
For shallow inserted-edge removal, decompose a retained occurrence by its outer
edge tag. A surviving host tag already supplies its inequality with the original
selected edge. An inserted tag survives precisely when its middle edge differs
from $f$. Its fibre position and hole observation are unchanged. Conversely,
construct the appropriate tagged edge from either of these witnesses.

For the general statement induct on $s$. A shallow $s$ and shallow $t$ give
the preceding result directly. If $t$ enters a box, append its retained inner
inventory and reassociate disjunction. For inside $s$, replacement of its
selected box preserves all retained outer-edge fibres; apply that fact and
the induction hypothesis to the inner composite, then reassociate disjunction.
Finally apply the exact inventory theorem for replacement at $s\star t$.
\end{proof}
\lean{Graph/DeepAssociativity}{114}{124}{Linear.retainedHole_inserted}
\lean{Graph/DeepAssociativity}{127}{143}{Linear.siteRetainedHole_comp}
\lean{Graph/DeepAssociativity}{147}{153}{Linear.plugAt_comp_hasHole_iff}

\begin{lemma}[Composite freshness]
Exact freshness for $(s\star t,A)$ is equivalent to both
\[
 \forall h,\ \neg(R_s(h)\land H_A(h))
 \quad\text{and}\quad
 \mathsf{SiteFresh}(t,A).
\]
\end{lemma}
\begin{proof}
Use the retained-union theorem. In the forward direction include either retained
inventory into the union and apply composite freshness. In the reverse direction,
split a retained witness according to the union and apply the corresponding
freshness assumption.
\end{proof}
\lean{Graph/DeepAssociativity}{156}{168}{Linear.siteFresh_comp_iff}

\begin{theorem}[Sharp context closure for composition]
Suppose $G$ and $H$ have unique holes. Then
$\mathsf{plugAt}(s\star t,A)$ has unique holes if and only if all three conditions
hold:
\begin{enumerate}
\item $A$ has unique holes;
\item $A$ is fresh for the composite site $s\star t$;
\item $\forall h,\ \neg(R_s(h)\land R_t(h))$.
\end{enumerate}
Thus the three inventories $R_s,R_t,H_A$ must be pairwise disjoint.
The intermediate graph $\mathsf{plugAt}(s,H)$ need not have unique holes.
\end{theorem}
\begin{proof}
By deep associativity it suffices to check uniqueness of
$\mathsf{plugAt}(s,\mathsf{plugAt}(t,A))$. Apply sharp context closure first
with host $G$, then with host $H$. This gives uniqueness of $A$, freshness for
$(t,A)$, and disjointness of $R_s$ from the hole inventory of
$\mathsf{plugAt}(t,A)$. The latter inventory is $R_t\cup H_A$ by the exact
inventory theorem. Splitting that union gives disjointness of $R_s$ from $R_t$
and from $H_A$. Combine the last condition with freshness for $(t,A)$ using
the composite-freshness lemma. Conversely, reverse each implication: the
pairwise exclusions make the inner result fresh for $s$, and both applications
of sharp closure establish the final result's uniqueness.
\end{proof}
\lean{Graph/DeepAssociativity}{173}{191}{Linear.plugAt_comp_unique_iff}

\begin{definition}[Certified composed plugging]\label{term:certified-composition}
Given contexts $G,H,A$, compatible sites $s,t$, composite freshness, and
disjoint retained inventories $R_s,R_t$, package
$\mathsf{plugAt}(s\star t,A)$ with the uniqueness proof just obtained.
This defines \texttt{Context.plugAtComp}; it requires no certificate that the
intermediate plugged graph is a context.
\end{definition}
\lean{Graph/DeepAssociativity}{195}{201}{Linear.Context.plugAtComp}

\subsection{Corolla units}

\begin{definition}[Substitution unit]\label{term:substitution-unit}
For a signature label $o$, let $\mathbf c_o$ be its one-edge corolla, with its
sole edge selected. This is a substitution unit of interface
$(\sigma_{\rm in}(o),\sigma_{\rm out}(o))$. It is an addressed graph, with no uniqueness premise on its label;
no free wire identifying an input vertex with an output vertex is introduced.
\end{definition}
We prove $\mathbf c_o[0:=H]\cong H$ and
$C[e:=\mathbf c_{\ell(e)}]\cong C$ using literal-label isomorphisms.

A corolla has no surviving edge after removing zero: an element of
$\{i:\mathsf{Fin}(1)\mid i\ne0\}$ contradicts the fact that every index below
one is zero. Its $i$th source and target attachment positions are its input
and output ports at $i$, respectively. These equations follow by unfolding
lookup in the mapped finite-range lists.
\lean{Graph/PlugUnits}{10}{10}{Net.PlugUnit.noEdge}
\lean{Graph/PlugUnits}{13}{14}{Net.PlugUnit.corolla_sourcePort}
\lean{Graph/PlugUnits}{17}{18}{Net.PlugUnit.corolla_targetPort}

\begin{lemma}[Left unit]
Replacing the sole edge of $\mathbf c_o$ by $H$ yields a graph isomorphic to $H$.
\end{lemma}
\begin{proof}
The vertex bijection discards the empty left summand in $\varnothing+H_V$.
The edge bijection discards the impossible surviving-corolla summand and retains
$H_E$. Their inverses include into the right summands. Both inverse laws follow
by cases, eliminating impossible elements.

Mapping an inserted port through the vertex projection gives the original port:
internal ports are immediate; input and output ports use the corolla attachment
equations. For each inserted edge, combine the incidence list maps and apply
this port equation, leaving the original source or target list. Surviving
corolla edges are impossible. All labels agree directly on the remaining
summands.
\end{proof}
\lean{Graph/PlugUnits}{21}{25}{Net.PlugUnit.leftVertices}
\lean{Graph/PlugUnits}{28}{32}{Net.PlugUnit.leftEdges}
\lean{Graph/PlugUnits}{35}{40}{Net.PlugUnit.left_insert}
\lean{Graph/PlugUnits}{88}{111}{Net.corolla_plug_iso}

\begin{lemma}[Right unit]
Replacing $e$ in $C$ by its own labelled corolla yields a graph isomorphic to $C$.
\end{lemma}
\begin{proof}
The vertex bijection discards the empty right summand. The edge bijection maps
each surviving host edge to itself and the inserted corolla edge to $e$.
Its inverse tests whether an edge equals $e$: if so use the sole corolla edge;
otherwise use the surviving-host subtype. The inverse laws follow from this
test and uniqueness of an element of $\mathsf{Fin}(1)$.

The vertex projection cancels inclusion of host ports. Enumerating all source
positions of $e$ reconstructs $C$'s source list: the lengths agree by typing,
and each indexed element agrees by lookup in the finite-range map. The target
list is reconstructed in the same way. Thus surviving-edge incidence agrees
by cancellation of host inclusion; the inserted corolla's full boundary list
maps to exactly the source or target list of $e$. Labels agree on both edge
cases. This proof also covers an edge whose source and target share vertices.
\end{proof}
\lean{Graph/PlugUnits}{43}{47}{Net.PlugUnit.rightVertices}
\lean{Graph/PlugUnits}{50}{63}{Net.PlugUnit.rightEdges}
\lean{Graph/PlugUnits}{66}{68}{Net.PlugUnit.right_host}
\lean{Graph/PlugUnits}{71}{76}{Net.PlugUnit.sourcePort_list}
\lean{Graph/PlugUnits}{79}{84}{Net.PlugUnit.targetPort_list}
\lean{Graph/PlugUnits}{114}{145}{Net.plug_corolla_iso}

\Needspace{16\baselineskip}
\subsection{Restoration and unit equations at deep sites}

\begin{definition}[Original selected fragment]\label{term:original-fragment}
The original fragment at a site $s$ is the corolla carrying the label of its
selected edge. For an inside site descend to its inner site. If the selected
edge is itself a box, its original fragment includes that complete box label
and its original body. The fragment has exactly the local depth and interface
required by $s$.
\end{definition}
\lean{Graph/DeepUnits}{11}{15}{Linear.Site.original}

\begin{lemma}[Deep restoration]
Substituting the original fragment at any site restores the entire host up to
recursive isomorphism.
\end{lemma}
\begin{proof}
First note that reinstalling a box's original body gives a literal-label
isomorphism with identity carriers: at the selected edge use its original
box-label equality, and elsewhere no label changes. Incidence and internal
labels are untouched.

Now induct on the site. For a shallow site apply the right unit. For an inside
site, the induction hypothesis restores its inner host up to recursive
isomorphism. Box congruence transports that isomorphism to the surrounding
updated box. Compose it with the just-proved isomorphism for reinstalling the
original body.
\end{proof}
\lean{Graph/DeepUnits}{18}{33}{Linear.replaceBox_original_iso}
\lean{Graph/DeepUnits}{37}{43}{Linear.plugAt_original}

Consequently, plugging the original fragment class at $s$ returns the original
host class, by quotient soundness. A selected one-edge corolla acts as identity
on every replacement class of its interface: quotient induction reduces this
to the left-unit isomorphism. Together with the composition equations above,
these are unit and associativity laws for the induced substitution operations
on recursive isomorphism classes. The quotient of pointed hosts and the
independent-site interchange law are developed in Section~\ref{sec:deep-interchange}.
\lean{Graph/DeepUnits}{46}{49}{Linear.plugAtModuloIso_original}
\lean{Graph/DeepUnits}{52}{57}{Linear.plugAtModuloIso_corolla}

\subsection{Regression proofs: self-loops, mixed depth, and repaired collisions}

For the concrete one-edge self-loop, instantiate shallow associativity with
two unary corollas. Instantiate the right unit with the loop's own corolla as
well. Both retain the loop's shared source/target vertex, demonstrating that
neither theorem relies on a non-self-loop hypothesis.
\lean{Graph/CompositionExamples}{8}{11}{Graph.Examples.loop_associativity}
\lean{Graph/CompositionExamples}{14}{15}{Graph.Examples.loop_unit}

Define a unary variable fragment $V$ as a corolla of the variable label. It has
interface $*\to *$ and no holes: any alleged hole observation would equate
\texttt{none} with \texttt{some}. Let $t$ select the sole edge of the earlier
replacement containing hole one. Its retained inventory is empty, because
there is no other edge in its one-element carrier.
\lean{Graph/CompositionExamples}{22}{23}{Linear.Examples.variableFragment}
\lean{Graph/CompositionExamples}{25}{27}{Linear.Examples.variableFragment_holeFree}
\lean{Graph/CompositionExamples}{30}{30}{Linear.Examples.repeatedSite}
\lean{Graph/CompositionExamples}{32}{34}{Linear.Examples.repeatedSite_retained_absent}

Take the earlier boxed two-hole host and let $s$ select hole zero inside it.
Inserting the one-hole replacement carrying label one gives the already proved
collision with the surviving hole one. This intermediate graph is invalid.
Now fill the inserted occurrence at $s\star t$ with $V$. The composite criterion
proves uniqueness: $V$ is hole-free, hence unique and fresh; and the empty
retained inventory of $t$ is disjoint from the retained host inventory.
\lean{Graph/CompositionExamples}{38}{46}{Linear.Examples.intermediate_collision_repaired}
The final graph still contains the host's hole one, witnessed by the surviving
inner edge. Hole zero is absent: it cannot be retained by $s$ because it was
the selected unique hole; it cannot be retained by $t$ because that inventory
is empty; and it cannot occur in $V$ because $V$ is hole-free. The exact composite
inventory theorem exhausts these three possibilities.
\lean{Graph/CompositionExamples}{49}{60}{Linear.Examples.repaired_inventory}

For a mixed-depth example, form a one-edge shell whose label is the box label
of the boxed two-hole host, and select that shell's sole edge. Insert the boxed
host, then follow its deep first-hole site. Deep associativity compares this
with filling the boxed host first and then inserting it into the shell.
\lean{Graph/CompositionExamples}{63}{64}{Linear.Examples.boxShell}
\lean{Graph/CompositionExamples}{66}{66}{Linear.Examples.shellSite}
\lean{Graph/CompositionExamples}{68}{71}{Linear.Examples.shallow_then_deep_associativity}
A three-stage instance adds the sole site of the repeated-hole replacement and
then the variable fragment; the two site-composition bracketings agree on the
resulting quotient class by three-stage coherence. Finally, restoring the
original fragment at the boxed first-hole site recovers the original boxed
host by deep restoration.
\lean{Graph/CompositionExamples}{74}{79}{Linear.Examples.three_stage_composition}
\lean{Graph/CompositionExamples}{82}{83}{Linear.Examples.deep_restoration}

\Needspace{18\baselineskip}
\section{Concrete focuses and the boundary argument of Lemma 39(1)}
\label{sec:concrete-focus}

This section returns to the proof order of Appendix~F. We formalize the
boundary argument in Proof~F.4: replacing focus-free fragments preserves the
context's focus mode and whether its source is an input. The result holds for
arbitrarily deep sites and finite sequences of fillings. It does not use
output-closure, rootedness, determinism, or transitions. It supplies the
initial/final agreement needed in Lemma~39(1), for the present substitution
representation. Connecting it to the paper's full specimen definition is still
an obligation; no concrete \texttt{SpecimenSystem} is constructed here.

\subsection{Global focus freedom and a unique shallow focus}
\begin{definition}[Focus labels and focus freedom]\label{term:focus-free}
For a mode $m$, the symbol $\mathsf{focusLabel}(d,m)$ is the primitive focus
symbol in the alphabet of depth $d$. At depth zero it uses the universe-lift
wrapper; at a successor depth it uses the primitive summand.
A label is \emph{focus-free} if it is a primitive other than every focus symbol,
or a box whose body has only focus-free labels. A graph is focus-free when all
its edge labels satisfy this recursively defined condition. Thus the condition
excludes focuses inside boxes at every depth, not merely at the outer layer.
\end{definition}
\lean{Graph/Focused}{13}{15}{Linear.focusLabel}
\lean{Graph/Focused}{18}{21}{Linear.LabelFocusFree}
\lean{Graph/Focused}{24}{25}{Linear.FocusFree}

\begin{definition}[Focussed hypernet]\label{term:focussed-hypernet}
A focussed hypernet consists of a hypernet $G$, a distinguished shallow edge
$f$, its mode $m$, and the following certificates. The label of $f$ is
$\mathsf{focusLabel}(d,m)$; every other edge has a focus-free label; the source
and target lists of $f$ are respectively $[s]$ and $[t]$; $s\ne t$; and it is
not the case that $s$ is an input and $t$ an output.
The label of $f$ is primitive, so has no box children. Consequently these
certificates enforce exactly one focus globally, with the shallowness,
self-acyclicity, and non-exposure required in \S6.2. The explicit singleton
incidence certificates record the endpoints without choosing positional casts.
\end{definition}
\lean{Graph/Focused}{29}{39}{Linear.Focussed}

\begin{lemma}[Focus-free addressed substitution]
If $G$ and $H$ are focus-free, so is $\mathsf{plugAt}(s,H)$ for every site $s$ in $G$
with the interface of $H$.
\end{lemma}
\begin{proof}
Induct on the site. At a shallow site, every surviving edge retains its old
label, and every inserted edge has a label of $H$, so both kinds are focus-free.
Inside a box, focus freedom of the host implies focus freedom of that box's
body. Apply the induction hypothesis to the body. Replacing its label preserves
focus freedom of that edge; every other outer label remains unchanged.
\end{proof}
\lean{Graph/Focused}{42}{63}{Linear.focusFree_plugAt}

\begin{definition}[Boundary initiality and finality]\label{term:boundary-finality}
For a focussed hypernet, define $\mathsf{Initial}$ to mean that the mode is
search and the focus source is an input. Define $\mathsf{Final}$ by replacing
search with backtrack. These predicates are defined on any interface for use
in the preservation lemmas. On states of type $*\to\epsilon$, there is exactly
one input, giving the boundary characterizations used in \S6.4 and Proof~F.4.
Finality is a syntactic predicate here; absence of outgoing transitions from
final states has not yet been proved, because transitions on this explicit focus-edge state representation are not yet
defined. Section~\ref{sec:search} proves the search-only boundary laws in coordinates.
\end{definition}
\lean{Graph/Focused}{69}{70}{Linear.Focussed.Initial}
\lean{Graph/Focused}{73}{74}{Linear.Focussed.Final}

\subsection{Preservation through one filling}
\begin{lemma}[Shallow filling preserves the focussed structure]
Let $G$ be focussed at $f$ and let $e\ne f$. Replacing $e$ by a focus-free $H$
produces a focussed hypernet with the same mode.
\end{lemma}
\begin{proof}
The new focus is the retained edge $\mathsf{inl}(f)$ with unchanged label.
Every other retained edge inherits focus freedom from $G$; every inserted edge
inherits it from $H$. Its endpoints are the host images of $s,t$, and its
incidence lists are their singleton lists by functoriality of list mapping.
The host map is injective, so these endpoints remain distinct. It maps internal
ports to internal ports and fixes input and output constructors. Therefore
exposure cannot be introduced.
\end{proof}
\lean{Graph/Focused}{77}{107}{Linear.Focussed.plug}

\begin{lemma}[Shallow boundary preservation and agreement]
A shallow filling as above is initial exactly when its host is initial, and
final exactly when its host is final. Any two focus-free fillings of the same
selected edge therefore agree on initiality and finality.
\end{lemma}
\begin{proof}
The focus mode is unchanged. The host image of its source is an input if and
only if the original source is an input, by case analysis on the three port
constructors. Substitute these facts into the two boundary predicates. For two
fillings, compose each equivalence with the inverse equivalence for the other
filling. This is exactly the argument of Proof~F.4.
\end{proof}
\lean{Graph/Focused}{111}{116}{Linear.Focussed.plug_initial_iff}
\lean{Graph/Focused}{118}{123}{Linear.Focussed.plug_final_iff}
\lean{Graph/Focused}{127}{135}{Linear.Focussed.specimen_boundary_agreement}

\begin{lemma}[Replacing a box body preserves an outer focus]
Replacing a box body by a focus-free body of the same interface preserves a
surrounding focussed hypernet.
\end{lemma}
\begin{proof}
The box edge cannot be the focus edge: their labels belong to distinct summands
of the alphabet. Keep the focus, its mode, and its two endpoints. All incidence
lists are unchanged by box-body replacement, so the self-acyclicity and
non-exposure certificates remain valid. The changed label is focus-free by the
hypothesis on the new body, and all other non-focus labels are unchanged.
\end{proof}
\lean{Graph/Focused}{138}{164}{Linear.Focussed.replaceBox}

\begin{definition}[A site avoiding the focus]\label{term:focus-avoiding-site}
A shallow site avoids the focus when its selected edge differs from the focus
edge. A site strictly inside a box automatically avoids the focus, since the
focus is shallow and primitive. This condition does not assert that the
selected edge is a hole: the preservation results intentionally apply to a
larger class of replacements than paper specimens require.
\end{definition}
\lean{Graph/Focused}{168}{172}{Linear.Focussed.Avoids}

\begin{lemma}[Deep focussed filling and boundary agreement]
A focus-free filling of any site avoiding the focus preserves the focussed
structure and both boundary predicates. Consequently two such fillings of one
site agree on initiality and finality.
\end{lemma}
\begin{proof}
For a shallow site use the shallow filling construction. For a site inside a
box, the box is distinct from the focus, hence its old body is focus-free.
Focus-free addressed substitution proves the modified body focus-free. Apply
the box-body preservation construction. In this second case the outer mode and
endpoints remain literally the same, so both boundary equivalences are
reflexive. Finally compose the equivalences for two fillings through their
common host.
\end{proof}
\lean{Graph/Focused}{175}{190}{Linear.Focussed.plugAt}
\lean{Graph/Focused}{193}{200}{Linear.Focussed.plugAt_boundary_iff}
\lean{Graph/Focused}{204}{211}{Linear.Focussed.deep_specimen_boundary_agreement}

\subsection{Finite fillings and the state space}
\begin{definition}[Finite filling sequence]\label{term:focus-fillings}
A filling sequence starts with a focussed graph and consists of finitely many
addressed substitutions of focus-free graphs at sites avoiding the current
focus. The constructors are the empty sequence and extension by one filling.
Sites may occur at any box depth. Each intermediate focus is the certified
focus transported by the preceding construction. In particular, successive
fillings of a context's holes are included whenever their sites are supplied.
This representation does not yet define a simultaneous tuple-substitution API
or identify all such sequences modulo permutation.
\end{definition}
\lean{Graph/Focused}{215}{221}{Linear.Focussed.Fillings}

\begin{theorem}[Boundary agreement for finite fillings]
Every finite filling sequence preserves initiality and finality. If two
sequences start from the same focussed context, their endpoints agree on both
predicates.
\end{theorem}
\begin{proof}
Induct on the sequence. The empty case is reflexivity. In the extension case,
compose the one-step boundary equivalence with the induction hypothesis. For
two sequences, compose the resulting equivalences through their common initial
context. No template relation between replacements is necessary. Thus once
concrete specimens supply these filling witnesses, the initial/final conclusion
of Lemma~39(1) follows without an extra assumed locality law.
\end{proof}
\lean{Graph/Focused}{224}{231}{Linear.Focussed.fillings_boundary_iff}
\lean{Graph/Focused}{235}{241}{Linear.Focussed.fillings_boundary_agreement}

\begin{definition}[Concrete state]\label{term:concrete-state}
A state packages a hypernet of type $*\to\epsilon$, a focussed structure on that
hypernet, and a proof that it contains no holes at any depth. This is the state
space of \S6.4 in our graph representation. It is not yet equipped with a
\texttt{Machine}: the coordinate search relation of Section~\ref{sec:search} still needs a realization
bridge, and substitution and compute transitions remain open.
\end{definition}
\lean{Graph/Focused}{247}{250}{Linear.State}

\subsection{Nonvacuous checks}
The regression graph has an input, an internal vertex, and an output, with a
focus edge from the input to the internal vertex followed by a unary hole.
It is obtained by relabelling the existing two-edge chain, preserving all
signatures. Its focus is certified directly: the only other edge is the hole,
the endpoints differ, and its target is internal, so the focus is not exposed.
\lean{Graph/FocusedExamples}{9}{13}{Linear.Examples.focusChain}
\lean{Graph/FocusedExamples}{15}{34}{Linear.Examples.focusChain_focussed}
The unary variable corolla is focus-free because a variable label cannot equal
a focus label. Filling the second edge of the search chain by this corolla
remains initial by shallow boundary preservation. The backtracking chain is
final, whereas the search chain is not, since the two mode constructors differ.
\lean{Graph/FocusedExamples}{36}{38}{Linear.Examples.variableFragment_focusFree}
\lean{Graph/FocusedExamples}{41}{46}{Linear.Examples.search_filling_initial}
\lean{Graph/FocusedExamples}{49}{51}{Linear.Examples.boundary_modes_distinct}
Finally, boxing the search chain gives a graph whose outer edge is a box but
whose body contains a focus. If it were focus-free, applying that predicate to
the outer box edge and then the first body edge would say that the search focus
label differs from itself, a contradiction. This checks that the definition
really inspects nested boxes.
\lean{Graph/FocusedExamples}{55}{60}{Linear.Examples.boxed_focus_not_free}

\Needspace{18\baselineskip}
\section{The five search rules on underlying graph coordinates}
\label{sec:search}

The next operational dependency is Fig.~9(a--e). We now define these five rules
on actual ordered incidences of a fixed hypernet and prove determinism and
preservation under shallow hole filling. The token is represented by a mode
and a vertex of the underlying graph. This follows the paper's observation
that search leaves the underlying graph unchanged. The relationship to the
explicit focus-edge representation is still an adequacy obligation: explicit insertion and incidence recovery are now proved in
Section~\ref{sec:realization}, and Section~\ref{sec:erasure} supplies general
erasure and search on explicit states. Equivalence with the paper's contextual
graph rewrites remains open. The relation is
defined even on graphs with holes; its intended state interpretation restricts
the underlying graph to be focus-free, hole-free, and of type $*\to\epsilon$,
with a token on a term vertex. The stronger structural lemmas below do not
need these restrictions.

\subsection{Operation kinds and certified incidence views}
\begin{definition}[Operation kind]\label{term:operation-kind}
The two operation kinds are passive and active. A function from operations to
these kinds specifies the partition from \S5.1. Finishing a passive operation
produces backtrack mode; finishing an active operation produces rewrite mode.
Neither finishing mode is search, by distinction of constructors.
\end{definition}
\lean{Graph/Search}{10}{13}{Linear.OperationKind}
\lean{Graph/Search}{16}{18}{Linear.OperationKind.finish}
\lean{Graph/Search}{20}{21}{Linear.OperationKind.finish_ne_search}
The embedding of primitive labels into each depth alphabet is the universe lift
at zero and the left summand at successor depth. It is injective: project the
lift at depth zero, or use injectivity of the left injection at successor depth.
\lean{Graph/Search}{24}{26}{Linear.primitiveLabel}
\lean{Graph/Search}{28}{32}{Linear.primitiveLabel_injective}

\begin{definition}[Search operation view]\label{term:search-operation}
A search operation view records an actual edge $e$ with operation label $o$,
its source $s$, and lists $A,T$ of arguments and thunks. Its certificates assert
\[
 \mathsf{source}(e)=[s],\qquad
 \mathsf{target}(e)=A\mathbin{+\!+}T,\qquad
 |A|=\mathsf{arity}(o).
\]
Thus $A$ is exactly the ordered term-argument prefix. The ambient graph's
signature typing already specifies the types of the remaining thunk targets.
This view exposes certified data; it adds no machine-law assumption.
\end{definition}
\lean{Graph/Search}{36}{45}{Linear.SearchOperation}

\begin{lemma}[Existence and uniqueness of operation data]
Every shallow operation edge has such a view. Two views of the same edge have
the same operation, source, and argument list. Every argument is a target of
the edge, and the argument list has no repeated vertices.
\end{lemma}
\begin{proof}
Signature typing gives a singleton source list and a target list at least as
long as the arity. Choose its unique source; take the first arity-many targets
for $A$ and drop that prefix for $T$. The take/drop identity gives the target
certificate, and the length bound gives the arity certificate.
For uniqueness, the edge label determines the operation by injectivity of
primitive labels. Equality of singleton source lists determines the source.
The two argument lists have the same length; taking that many entries from the
equal target lists identifies them. Argument membership follows by inclusion
of the left append summand. No repetitions follows by restriction of the
whole target list's no-repetition certificate.
\end{proof}
\lean{Graph/Search}{52}{68}{Linear.SearchOperation.exists_view}
\lean{Graph/Search}{71}{82}{Linear.SearchOperation.data_unique}
\lean{Graph/Search}{84}{86}{Linear.SearchOperation.argument_mem}
\lean{Graph/Search}{88}{91}{Linear.SearchOperation.arguments_nodup}

\subsection{Local outcomes and the search relation}
\begin{definition}[Entry and return outcomes]\label{term:search-outcomes}
Entering an operation with argument list $A$ yields search at its first
argument if $A$ is nonempty. If $A$ is empty, it yields the operation's finishing
mode at its source. Returning from argument index $i$ yields search at argument
$i+1$ if that entry exists; otherwise it yields the finishing mode at the
source. Indices are zero-based. Return transitions additionally require that
argument $i$ exists, so the last case cannot be reached by returning from an
out-of-range index.
\end{definition}
\lean{Graph/Search}{94}{97}{Linear.SearchOperation.enter}
\lean{Graph/Search}{100}{103}{Linear.SearchOperation.resume}
These outcomes are independent of the chosen view of an edge. Substitute the
operation, source, and argument equalities from data uniqueness into their
definitions; the resulting expressions are identical.
\lean{Graph/Search}{105}{108}{Linear.SearchOperation.enter_unique}
\lean{Graph/Search}{110}{113}{Linear.SearchOperation.resume_unique}

\begin{definition}[Search position and step]\label{term:search-step}
A search position is a pair $(m,v)$ of mode and outer graph vertex. The search
relation has three constructors covering the five rules:
\begin{enumerate}
\item At a variable edge's singleton source $v$, search becomes rewrite at $v$
(Fig.~9a).
\item At an operation's source, search yields its entry outcome
(Fig.~9b for nonempty arguments and Fig.~9e for empty arguments).
\item Backtracking at argument $i$ yields its return outcome
(Fig.~9c when another argument exists and Fig.~9d when it does not).
\end{enumerate}
Every edge mentioned is shallow. Box contents and thunk targets are not searched
by these rules. No rule changes the underlying graph, and this relation contains
neither substitution nor extrinsic computation steps.
\end{definition}
\lean{Graph/Search}{151}{151}{Linear.SearchPosition}
\lean{Graph/Search}{155}{163}{Linear.SearchStep}
The four named operation rules specialize the two grouped constructors by
rewriting the argument list or the lookup of the next argument. Together with
the variable constructor, they expose precisely the five cases of Fig.~9(a--e).
\lean{Graph/Search}{169}{172}{Linear.SearchStep.first_argument}
\lean{Graph/Search}{175}{177}{Linear.SearchStep.nullary}
\lean{Graph/Search}{180}{183}{Linear.SearchStep.next_argument}
\lean{Graph/Search}{186}{189}{Linear.SearchStep.last_argument}

\subsection{Determinism and initial/final boundaries}
\begin{theorem}[Search determinism]
Two search steps with the same input coordinates have the same output
coordinates.
\end{theorem}
\begin{proof}
Case-split on their rule constructors. Search and backtrack inputs cannot agree.
For two variable steps the result is rewrite at the same vertex. For a variable
step and an operation-entry step, uniqueness of the edge having the given
source forces their edges equal, contradicting their distinct primitive labels.
For two operation-entry steps, source incidence uniqueness again identifies the
edges; entry-outcome uniqueness identifies the results.
For two return steps, both starting vertices occur in their operation target
lists. Target incidence uniqueness identifies the edges and hence their
argument lists. Since these lists have no repetitions, equal arguments force
equal indices. Return-outcome uniqueness then identifies the results.
The more general helper permits two separately named input coordinates with an
explicit equality; the deterministic theorem specializes that equality to
reflexivity. No determinism hypothesis is assumed.
\end{proof}
\lean{Graph/Search}{192}{232}{Linear.SearchStep.result_unique}
\lean{Graph/Search}{235}{237}{Linear.SearchStep.deterministic}

\begin{lemma}[Search boundary laws]
A rewrite-mode coordinate has no search step. A backtracking coordinate at an
input has no search step. No search step yields search mode at an input.
\end{lemma}
\begin{proof}
No constructor starts in rewrite mode. A step from backtrack must return from
an argument, which is a target of an edge. The graph's incidence conditions
forbid an input vertex from being a target.
For the last assertion, a variable step ends in rewrite, and finishing an
operation never produces search. Therefore a search-mode result must be an
argument selected on entry or on return. It is again an edge target and cannot
be an input. Apply this stronger non-input assertion to rule out any proposed
step into an initial coordinate.
\end{proof}
\lean{Graph/Search}{240}{242}{Linear.SearchStep.no_rewrite}
\lean{Graph/Search}{245}{249}{Linear.SearchStep.final_no_step}
\lean{Graph/Search}{252}{265}{Linear.SearchStep.search_target_not_input}
\lean{Graph/Search}{268}{270}{Linear.SearchStep.no_initial_target}
The finality assertion here concerns search steps only. It is not yet the
paper's assertion that final states have no transition of any kind; that also
requires the substitution and compute rules and the focus-edge bridge.

\subsection{Preservation under context filling}
\begin{lemma}[Transport of retained operation data]
Replacing an edge other than an operation edge retains a view of that operation.
Entry and return outcomes commute with the host port map.
\end{lemma}
\begin{proof}
Keep the operation label and use the surviving-edge injection. Map the source,
argument list, and thunk list through the host map. Mapping a singleton and an
append gives the two incidence certificates, and mapping preserves argument
length. Entry commutation follows by the empty/nonempty cases of the argument
list. Return commutation follows because lookup in a mapped list is the mapped
optional lookup; split according to whether the next argument exists. Modes
are unchanged by the port map.
\end{proof}
\lean{Graph/Search}{116}{127}{Linear.SearchOperation.plug}
\lean{Graph/Search}{129}{136}{Linear.SearchOperation.plug_enter}
\lean{Graph/Search}{138}{145}{Linear.SearchOperation.plug_resume}

\begin{theorem}[Existing context search steps survive filling]
Suppose $G$ has a shallow hole edge $h$ and a search step $p\to q$. For every
interface-compatible filling $H$, the graph $G[h:=H]$ has a search step from
the host image of $p$ to the host image of $q$.
\end{theorem}
\begin{proof}
In every rule, the interacting edge is a variable or operation and therefore
cannot equal the hole edge. In the variable case, retain that edge, its label,
and its mapped singleton source to apply the same rule. In the entry case,
use the retained operation view and entry commutation. In the return case,
lookup commutes with mapping, so the old argument witness supplies the mapped
argument witness; apply return commutation to identify the result.
\end{proof}
\lean{Graph/Search}{274}{308}{Linear.SearchStep.plug}
This theorem starts with a step of the context itself. Lemma~39(3) instead
starts with a step of a filled specimen and must establish that it comes from
the context. Extracting that step requires the entering/exiting analysis,
rootedness, output-closure, and the focus-edge bridge. The preservation theorem
proves one needed direction of locality, not the whole lemma.

\Needspace{24\baselineskip}
\subsection{Concrete checks of rule selection}
The fixture has one operation with two term arguments and one thunk. Its graph
is a corolla with targets $0,1,2$ in that order; the operation view certifies
that the argument list is $[0,1]$ and the thunk list is $[2]$. The incidence
certificates are checked by computation.
\lean{Graph/SearchExamples}{9}{9}{Linear.SearchExamples.ops}
\lean{Graph/SearchExamples}{11}{12}{Linear.SearchExamples.graph}
\lean{Graph/SearchExamples}{14}{23}{Linear.SearchExamples.operation}
Entry selects target $0$ and return from target $0$ selects target $1$. Return
from target $1$ finishes at the source in the declared mode, for either operation
kind. These are direct applications of the corresponding rule theorems.
\lean{Graph/SearchExamples}{26}{28}{Linear.SearchExamples.enter_first}
\lean{Graph/SearchExamples}{31}{33}{Linear.SearchExamples.return_next}
\lean{Graph/SearchExamples}{36}{38}{Linear.SearchExamples.finish_skips_thunk}
A hypothetical return step to search at thunk target $2$ would, by determinism,
have the same result as the finishing step. It would identify the finishing
mode with search, which is impossible. This checks that the full target list
has not accidentally been used as the argument list.
\lean{Graph/SearchExamples}{41}{45}{Linear.SearchExamples.cannot_search_thunk}
The existing nullary value corolla admits the empty argument view and finishes
by backtracking when declared passive. The unary variable corolla uses the
variable constructor and triggers rewrite without moving the coordinate.
\lean{Graph/SearchExamples}{48}{54}{Linear.SearchExamples.value_finishes}
\lean{Graph/SearchExamples}{57}{60}{Linear.SearchExamples.variable_triggers}

\Needspace{32\baselineskip}
\section{Explicit focus insertion and realization of typed search positions}
\label{sec:realization}

This section advances the operational obligation in the remaining-work list.
A term-vertex coordinate now determines an actual focussed hypernet, and on a
hole-free graph of type $*\to\epsilon$ it determines an actual state. The
construction proves all incidence and typing laws, rather than assuming a
realization interface. Contracting the inserted focus recovers the original
ordered incidence lists. Search preserves the term-vertex condition, so every
typed coordinate search step has two such state realizations.

This is the insertion direction of the bridge. It does not yet erase an
arbitrary pre-existing focussed state, establish inverse laws modulo graph
isomorphism, or identify the coordinate relation with the paper's contextual
interaction rules. The endpoint theorem below must not be read as a proof of
those still-open adequacy statements. Substitution, compute transitions, and
the concrete specimen interpretation also remain open.

\Needspace{16\baselineskip}
\subsection{Splitting a vertex}
\begin{definition}[Focus insertion carriers and maps]\label{term:focus-insertion}
Let $G$ be a hypernet and let $v$ be a term vertex that is not an output. Extend
the internal-vertex carrier from $V$ to $V+1$ and the edge carrier from $E$ to
$E+1$. Write $z$ for the fresh internal vertex, $\iota$ for the inclusion of old
ports, and $f$ for the fresh edge. Boundaries are unchanged. Define
\[
 S_v(w)=\begin{cases}z&w=v,\\\iota(w)&w\ne v,\end{cases}
 \qquad
 C_v(\iota(w))=w,\quad C_v(z)=v.
\]
Here $S_v$ redirects old source occurrences of $v$, while $C_v$ contracts the
new focus. Old internal vertices retain their wire labels and $z$ has label
term. The use of a singleton carrier gives explicit finite enumerations.
\end{definition}
\lean{Graph/FocusInsertion}{10}{10}{Linear.FocusInsertion.Vertex}
\lean{Graph/FocusInsertion}{12}{12}{Linear.FocusInsertion.host}
\lean{Graph/FocusInsertion}{14}{14}{Linear.FocusInsertion.fresh}
\lean{Graph/FocusInsertion}{16}{20}{Linear.FocusInsertion.collapse}
\lean{Graph/FocusInsertion}{30}{32}{Linear.FocusInsertion.split}
\lean{Graph/FocusInsertion}{63}{65}{Linear.FocusInsertion.internalLabel}

\begin{lemma}[Map properties]
The inclusion $\iota$ is injective and never reaches $z$. The contraction
satisfies $C_v\iota=\mathrm{id}$ and $C_vS_v=\mathrm{id}$; consequently $S_v$
is injective. No image of $S_v$ equals $\iota(v)$. Every new port is either an
old port's image or $z$. Both inclusion and redirection preserve wire labels.
\end{lemma}
\begin{proof}
Case analysis on internal, input, and output constructors proves the inclusion
claims; on internals use injectivity and disjointness of sum tags. For
$C_vS_v$, split on $w=v$: the fresh port contracts to $v$ in the equality case,
and otherwise contraction undoes inclusion. Applying $C_v$ to equal redirection
images proves injectivity. An image cannot equal $\iota(v)$: in the equality
case it is fresh; otherwise inclusion injectivity would force $w=v$.
The port classification uses the two internal sum cases and the boundary
constructors; every inhabitant of the fresh singleton carrier equals its sole
member. Inclusion preserves labels by cases. Redirection preserves labels
because only $v$ changes image, and both $v$ and $z$ have term label.
\end{proof}
\lean{Graph/FocusInsertion}{22}{22}{Linear.FocusInsertion.collapse_host}
\lean{Graph/FocusInsertion}{24}{24}{Linear.FocusInsertion.host_ne_fresh}
\lean{Graph/FocusInsertion}{26}{27}{Linear.FocusInsertion.host_injective}
\lean{Graph/FocusInsertion}{34}{38}{Linear.FocusInsertion.collapse_split}
\lean{Graph/FocusInsertion}{40}{42}{Linear.FocusInsertion.split_injective}
\lean{Graph/FocusInsertion}{44}{49}{Linear.FocusInsertion.split_ne_host}
\lean{Graph/FocusInsertion}{51}{61}{Linear.FocusInsertion.vertex_cases}
\lean{Graph/FocusInsertion}{67}{68}{Linear.FocusInsertion.host_label}
\lean{Graph/FocusInsertion}{70}{75}{Linear.FocusInsertion.split_label}

\begin{definition}[Inserted incidence lists]
For every old edge $e$, use $S_v(\mathsf{source}_G(e))$ for its source list and
$\iota(\mathsf{target}_G(e))$ for its target list, mapping entry by entry.
The fresh edge has source $[\iota(v)]$, target $[z]$, and the chosen focus mode
as label. Old edge labels, including whole boxed bodies, are unchanged.
\end{definition}
\lean{Graph/FocusInsertion}{77}{79}{Linear.FocusInsertion.source}
\lean{Graph/FocusInsertion}{81}{83}{Linear.FocusInsertion.target}

\begin{lemma}[Incidence uniqueness and coverage]
The inserted incidence lists satisfy source and target uniqueness. A port
occurs in a source list exactly when it is not an output, and in a target list
exactly when it is not an input.
\end{lemma}
\begin{proof}
For two old source edges, pull equality of mapped vertices back along the
injective $S_v$ and apply source uniqueness in $G$. An old source cannot meet
the fresh edge's source because $S_v$ avoids $\iota(v)$. There is only one
fresh edge. The target proof is identical with inclusion in place of
redirection and disjointness from $z$ in place of avoidance of $\iota(v)$.

For source coverage, all mapped old sources are non-outputs: a redirected one
is internal, and any other one retains its old boundary kind. The fresh edge's
source is a non-output by the hypothesis on $v$. Conversely, an old image
$\iota(w)$ with $w=v$ is sourced by the fresh edge. If $w\ne v$, use coverage
in $G$ and its retained source edge. The fresh port $z$ is sourced by the old
edge that sourced $v$, which exists because $v$ was not an output.
For target coverage, retained targets remain non-inputs and the fresh target
is internal. Conversely an old non-input image inherits a target edge from
$G$, and $z$ is targeted by the fresh focus edge.
\end{proof}
\lean{Graph/FocusInsertion}{85}{104}{Linear.FocusInsertion.source_unique}
\lean{Graph/FocusInsertion}{106}{125}{Linear.FocusInsertion.target_unique}
\lean{Graph/FocusInsertion}{127}{154}{Linear.FocusInsertion.source_coverage}
\lean{Graph/FocusInsertion}{156}{174}{Linear.FocusInsertion.target_coverage}

\begin{theorem}[Well-formed explicit focus insertion]
The construction gives a finite typed hypernet with the same external
interface as $G$.
\end{theorem}
\begin{proof}
Enumerate each extended carrier by the sum of the old enumeration and the
singleton enumeration. Old incidence lists retain their no-repetition property
because both maps are injective; the new lists are singletons. Use the proved
uniqueness and coverage lemmas. For typing, mapping a port list and then reading
labels is the same as reading the old labels, by the label-preservation laws.
The new source and target both have term label, exactly the focus signature.
These checks populate every field of the concrete net record.
\end{proof}
\lean{Graph/FocusInsertion}{177}{226}{Linear.FocusInsertion.insert}

\Needspace{16\baselineskip}
\subsection{Focus certification, contraction, and holes}
\begin{theorem}[Unique focus and contraction]
If $G$ is globally focus-free, the inserted graph is focussed in the sense of
\S6.2. Contracting the new focus restores each original source and target list.
\end{theorem}
\begin{proof}
Choose the fresh edge as focus. Every old label remains globally focus-free;
the singleton fresh-edge carrier has no second member. The focus endpoints are
$\iota(v)$ and $z$, which are distinct. Its target is internal, so the focus is
not exposed. Its incidence lists and mode are those specified by construction.
To recover source lists, compose their map by $S_v$ with the map by $C_v$ and
use $C_vS_v=\mathrm{id}$. For target lists use $C_v\iota=\mathrm{id}$.
These statements concern the retained edges; the focus itself is deleted when
interpreting contraction. They establish recovery for the constructed family,
not a general erasure operation on arbitrary focussed graphs.
\end{proof}
\lean{Graph/FocusInsertion}{229}{245}{Linear.FocusInsertion.focussed}
\lean{Graph/FocusInsertion}{248}{255}{Linear.FocusInsertion.collapse_sources}
\lean{Graph/FocusInsertion}{257}{264}{Linear.FocusInsertion.collapse_targets}

\begin{lemma}[Hole freedom and boundary observations]
If $G$ is hole-free at every depth, so is its focus insertion. The explicit
focus is initial exactly when its mode is search and $v$ is an input; it is
final exactly when its mode is backtrack and $v$ is an input.
\end{lemma}
\begin{proof}
Represent a hypothetical hole occurrence by an outer edge and a label-local
position. If the edge is retained, its unchanged label gives a hole occurrence
in $G$, a contradiction. If it is the fresh edge, its focus label is not a hole
and has no children. This covers zero and successor depth alphabets.
For boundary observations, the mode is unchanged and the source is
$\iota(v)$. Inclusion preserves input status, by cases on the port constructor.
\end{proof}
\lean{Graph/FocusInsertion}{267}{286}{Linear.FocusInsertion.holeFree}
\lean{Graph/FocusInsertion}{289}{293}{Linear.FocusInsertion.boundary_iff}

\Needspace{16\baselineskip}
\subsection{Search preserves term typing}
\begin{lemma}[Typed operation data]
An operation view's source has term label. Mapping its argument list through
vertex labels gives exactly the list of arity-many term labels. Consequently
every argument vertex has term label.
\end{lemma}
\begin{proof}
Rewrite the operation's source-typing certificate by its singleton source and
operation label; the singleton signature forces the source label to be term.
For arguments, rewrite the target-typing certificate using the argument/thunk
append decomposition. Take an argument-length prefix on both sides. The
arity certificate ensures that this is exactly the entire argument-label list
on one side and the term prefix of the signature on the other. Finally map
membership of an argument into this repeated-term list: every member equals
term.
\end{proof}
\lean{Graph/SearchRealization}{14}{18}{Linear.SearchOperation.source_term}
\lean{Graph/SearchRealization}{21}{35}{Linear.SearchOperation.arguments_typed}
\lean{Graph/SearchRealization}{37}{41}{Linear.SearchOperation.argument_term}

\begin{theorem}[Term typing preservation]
If a search step starts at a term vertex, its result is at a term vertex.
\end{theorem}
\begin{proof}
A variable trigger does not move, so inherits the premise. Operation entry
selects either the first argument or the source when there are no arguments.
Operation return selects either the next argument or the source on completion.
Apply the operation data lemmas in each case. Neither operation case needs the
premise on the old vertex; its result's type follows from graph typing alone.
\end{proof}
\lean{Graph/SearchRealization}{46}{60}{Linear.SearchStep.preserves_term}

\Needspace{16\baselineskip}
\subsection{Actual states from typed coordinates}
\begin{definition}[Typed position and realization]\label{term:typed-position}
For a graph $G:*\to\epsilon$, a typed position is a search coordinate whose
vertex is certified to have term label. Such a vertex cannot be an output,
since the output index type is empty. If $G$ is globally focus-free and
hole-free, realize a typed position by inserting its mode at its vertex and
packaging the resulting graph, unique-focus certificate, and hole-freedom
certificate into the concrete state record.
\end{definition}
\lean{Graph/SearchRealization}{63}{64}{Linear.TypedPosition}
\lean{Graph/SearchRealization}{70}{74}{Linear.TypedPosition.not_output}
\lean{Graph/SearchRealization}{77}{81}{Linear.TypedPosition.realize}

\begin{definition}[Initial attachment]\label{term:initial-attachment}
The initial typed position is search at the sole input, whose wire label is
term by the interface. Realizing it is an initial state: the constructed focus
has search mode and its source is that input. More generally, the realized
state's initial and final predicates coincide with their coordinate boundary
predicates, by the insertion boundary lemma.
\end{definition}
\lean{Graph/SearchRealization}{84}{85}{Linear.TypedPosition.initial}
\lean{Graph/SearchRealization}{88}{89}{Linear.TypedPosition.initial_realize}
\lean{Graph/SearchRealization}{92}{95}{Linear.TypedPosition.realize_boundary}
This discharges the graph-construction part of initial attachment. Proposition~38
still needs the concrete observation interface and specimen witnesses required
by \texttt{initial\_lift}.

\begin{definition}[Advancing a typed position]
Given a search step from a typed position, equip its result coordinate with
the term-typing proof supplied by preservation. This gives the next typed
position on the same underlying graph.
\end{definition}
\lean{Graph/SearchRealization}{98}{100}{Linear.TypedPosition.advance}

\begin{theorem}[Realization of search endpoints]
For every search step from a typed position on a focus-free, hole-free graph
$G:*\to\epsilon$, realize the initial position and its typed successor as
states $s,t$. Their explicit focus modes are the corresponding coordinate
modes. At each endpoint, collapsing the inserted focus recovers every
original edge's source and target list in $G$.
\end{theorem}
\begin{proof}
Use the realization constructor for the input and for the typed successor.
Their focus modes are correct by the definition of insertion. Apply the two
incidence-recovery theorems at each coordinate. The common underlying graph
is unchanged; only the position and mode of the inserted focus may differ.
The statement supplies actual state endpoints and recovery certificates, not
a new definition of a contextual graph-rewrite step between arbitrary states.
\end{proof}
\lean{Graph/SearchRealization}{104}{122}{Linear.TypedPosition.step_realizes}

\Needspace{22\baselineskip}
\Needspace{16\baselineskip}
\subsection{Boundary and interior regression checks}
The existing value corolla is focus-free because its operation label cannot
be a focus label. Realizing its initial typed position gives the explicit
state \texttt{initialValue}. Advancing by its passive nullary search step and
realizing the result gives \texttt{finalValue}. Their initial and final
certificates follow from their modes and common input source. The initial
state's distinct focus endpoints are guaranteed by its focussed certificate.
\lean{Graph/RealizationExamples}{8}{10}{Linear.Examples.value_focusFree}
\lean{Graph/RealizationExamples}{13}{14}{Linear.Examples.initialValue}
\lean{Graph/RealizationExamples}{17}{19}{Linear.Examples.finalValue}
\lean{Graph/RealizationExamples}{21}{22}{Linear.Examples.value_boundary_states}
\lean{Graph/RealizationExamples}{25}{27}{Linear.Examples.initial_value_has_distinct_endpoints}
The two-hole chain is also globally focus-free, because neither hole label can
be a focus. Insert search at its internal vertex: the first edge still targets
the old vertex, the second edge now starts at the fresh vertex, and the new
focus connects the old vertex to the fresh one. Unfolding the incidence maps
proves all four equalities. This focus is not initial because its source is
internal. Unlike the value example, this test deliberately retains holes and
therefore tests focussed hypernets rather than closed states.
\lean{Graph/RealizationExamples}{29}{31}{Linear.Examples.twoHoles_focusFree}
\lean{Graph/RealizationExamples}{35}{42}{Linear.Examples.interior_insertion}
\lean{Graph/RealizationExamples}{45}{49}{Linear.Examples.interior_not_initial}

\addtocontents{toc}{\protect\newpage}
\Needspace{32\baselineskip}
\section{Focus insertion, erasure, and recovery}
\label{sec:erasure}

We now erase the focus of an arbitrary focussed graph with no outputs and
prove that reinsertion recovers the original graph and focus up to isomorphism.
Unlike the previous section's recovery result, the input need not have been
constructed by our insertion function. The converse round trip, insertion
followed by erasure, recovers the underlying graph. These constructions
support the state-search semantics developed in Section~\ref{sec:state-search}.

\Needspace{16\baselineskip}
\subsection{Contraction maps and the retained carriers}
Fix a focussed graph $G:I\to\epsilon$, focus edge $e_f$, source $s$, and an
internal vertex $t$ equal to the focus target. We first work with this explicit
witness, then prove that every such graph supplies it.

\begin{definition}[General focus erasure maps]\label{term:focus-erasure}
Remove $t$ from the internal carrier and $e_f$ from the edge carrier, using
subtypes of names unequal to the deleted names. Let $j$ include the remaining
ports into the old ports. If $w\ne t$, let $k(w)$ keep $w$ with its non-deletion
certificate. Self-acyclicity ensures $s\ne t$, so $p=k(s)$ is defined. Set
\[
 c(w)=\begin{cases}p&w=t,\\k(w)&w\ne t,\end{cases}
 \qquad
 a(x)=\begin{cases}t&j(x)=s,\\j(x)&j(x)\ne s.\end{cases}
\]
The map $c$ contracts the focus. The map $a$ recovers the appropriate old
\emph{source} representative, choosing $t$ at the joined port; $j$ recovers
the old \emph{target} representative. Boundary ports and their order are
unchanged.
\end{definition}
\lean{Graph/FocusErasure}{11}{11}{Linear.FocusErasure.Vertex}
\lean{Graph/FocusErasure}{13}{13}{Linear.FocusErasure.lift}
\lean{Graph/FocusErasure}{15}{19}{Linear.FocusErasure.keep}
\lean{Graph/FocusErasure}{38}{38}{Linear.FocusErasure.position}
\lean{Graph/FocusErasure}{40}{42}{Linear.FocusErasure.contract}
\lean{Graph/FocusErasure}{62}{64}{Linear.FocusErasure.expand}

\begin{lemma}[Port-map identities]
The map $j$ is injective, avoids $t$, and satisfies $j(k(w))=w$.
Moreover $j(p)=s$, $c(t)=c(s)=p$, $j(c(w))=w$ for $w\ne t$, and $c(j(x))=x$.
The expansion $a(x)$ never equals $s$; it satisfies $a(c(w))=w$ for $w\ne s$
and $c(a(x))=x$.
\end{lemma}
\begin{proof}
The inclusion and keeping claims follow by port-constructor cases; subtype
extensionality gives injectivity on internal ports, and the subtype certificate
excludes $t$. The inequality $s\ne t$ is exactly focus self-acyclicity rewritten
by the chosen target witness. The identities for $p,c$ follow by unfolding and
splitting on whether the argument is $t$. Apply inclusion injectivity to derive
$c(j(x))=x$ from $j(c(j(x)))=j(x)$.
For expansion, when $j(x)=s$ its result is $t\ne s$; otherwise it is already
different from $s$. To prove $a(c(w))=w$, split on $w=t$. At $t$, contraction
lands at $p$ and expansion returns $t$. Otherwise inclusion after contraction
returns $w\ne s$, so expansion leaves it unchanged. For $c(a(x))=x$, split
on $j(x)=s$. In that case $x=p$ by inclusion injectivity; otherwise use
$c(j(x))=x$.
\end{proof}
\lean{Graph/FocusErasure}{21}{22}{Linear.FocusErasure.lift_keep}
\lean{Graph/FocusErasure}{24}{25}{Linear.FocusErasure.lift_injective}
\lean{Graph/FocusErasure}{27}{31}{Linear.FocusErasure.lift_ne_target}
\lean{Graph/FocusErasure}{33}{35}{Linear.FocusErasure.source_ne_target}
\lean{Graph/FocusErasure}{44}{45}{Linear.FocusErasure.lift_position}
\lean{Graph/FocusErasure}{47}{48}{Linear.FocusErasure.contract_target}
\lean{Graph/FocusErasure}{50}{52}{Linear.FocusErasure.lift_contract}
\lean{Graph/FocusErasure}{54}{56}{Linear.FocusErasure.contract_lift}
\lean{Graph/FocusErasure}{58}{59}{Linear.FocusErasure.contract_source}
\lean{Graph/FocusErasure}{66}{71}{Linear.FocusErasure.expand_ne_source}
\lean{Graph/FocusErasure}{73}{78}{Linear.FocusErasure.expand_contract}
\lean{Graph/FocusErasure}{80}{86}{Linear.FocusErasure.contract_expand}

\begin{lemma}[Injectivity on retained incidence lists]
A source of a retained edge cannot equal $s$, and a target of a retained edge
cannot equal $t$. Contraction is injective on vertices other than $s$, and also
on vertices other than $t$.
\end{lemma}
\begin{proof}
If a retained edge sourced $s$, source uniqueness with the focus edge would
identify the two edges, contrary to retention. Target uniqueness proves the
other exclusion. For injectivity off $s$, apply $a$ to equality of contracted
ports and use $a(c(w))=w$. For injectivity off $t$, apply $j$ and use
$j(c(w))=w$. These are restricted injectivity assertions: $c$ itself identifies
the two distinct focus endpoints.
\end{proof}
\lean{Graph/FocusErasure}{88}{91}{Linear.FocusErasure.source_avoids}
\lean{Graph/FocusErasure}{93}{97}{Linear.FocusErasure.target_avoids}
\lean{Graph/FocusErasure}{99}{102}{Linear.FocusErasure.source_injective}
\lean{Graph/FocusErasure}{104}{107}{Linear.FocusErasure.target_injective}

\begin{lemma}[Label preservation]
Both focus endpoints have term label. Keeping an undeleted port preserves its
label, and contraction preserves every port's label.
\end{lemma}
\begin{proof}
Rewrite the focus source and target typing certificates using their singleton
incidences and focus label; its signature is term-to-term. Keeping preserves
labels by cases on internal and boundary ports. For contraction, all ports
other than $t$ use keeping. At $t$, its image is the kept source and both
endpoints have the same term label.
\end{proof}
\lean{Graph/FocusErasure}{109}{112}{Linear.FocusErasure.source_term}
\lean{Graph/FocusErasure}{114}{118}{Linear.FocusErasure.target_term}
\lean{Graph/FocusErasure}{120}{122}{Linear.FocusErasure.keep_label}
\lean{Graph/FocusErasure}{124}{132}{Linear.FocusErasure.contract_label}

\Needspace{16\baselineskip}
\subsection{A valid focus-free underlying graph}
\begin{theorem}[General erasure construction]
Delete $e_f$ and $t$, retain every other edge label, and map all retained source
and target lists through $c$. The result is a finite typed hypernet of the same
interface $I\to\epsilon$.
\end{theorem}
\begin{proof}
Finite subtype enumerations supply the carriers. Restricted contraction
injectivity on retained sources and targets preserves no-repetition. Given a
shared contracted source of two retained edges, restricted injectivity recovers
a shared old source, so old source uniqueness identifies the edge names.
The target proof is the same with the target restriction.

For source coverage, there are no output ports. Given an erased port $x$, use
the old source coverage at $a(x)$. Since $a(x)\ne s$, its source edge cannot be
the focus; contraction of that source gives $c(a(x))=x$.
For target coverage, a retained target $w$ differs from $t$, so
$j(c(w))=w$ and its input status is unchanged. Conversely, if $x$ is not an
input, old target coverage at $j(x)$ gives an edge. It is not the focus because
$j(x)\ne t$. Its contracted target is $c(j(x))=x$.
Finally, mapping contracted lists through labels equals reading the old labels,
by label preservation, so both signature-typing certificates follow from $G$.
\end{proof}
\lean{Graph/FocusErasure}{135}{210}{Linear.FocusErasure.erase}

\begin{lemma}[Erased invariants]
The erased graph is globally focus-free, its distinguished port $p$ has term
label, and erasure preserves global hole-freedom.
\end{lemma}
\begin{proof}
Every retained edge was different from the unique focus, so its entire label,
including any box body, is focus-free. The port $p$ keeps the term-labelled
source, hence has term label. A hypothetical hole in the erased graph lies at
a retained edge with unchanged label and label-local position; the same
occurrence would be a hole in the original graph. Thus a hole-free original
has a hole-free erasure.
\end{proof}
\lean{Graph/FocusErasure}{212}{213}{Linear.FocusErasure.erased_focusFree}
\lean{Graph/FocusErasure}{215}{217}{Linear.FocusErasure.position_term}
\lean{Graph/FocusErasure}{219}{222}{Linear.FocusErasure.erased_holeFree}

\Needspace{16\baselineskip}
\subsection{Reinsertion recovers an arbitrary focussed graph}
\begin{definition}[Restoring a deleted name]\label{term:focus-round-trip}
For any type $A$ and distinguished $b\in A$, the restoration bijection
$(A\setminus\{b\})+1\to A$ sends retained names to themselves and the fresh
singleton name to $b$. Its inverse tests whether a name equals $b$.
Applying this to $t$ and $e_f$ gives the vertex and edge bijections needed for
reinsertion. Reinsert the original focus mode at $p$ in the erased graph.
The port $p$ is a non-output because the output interface is empty, and its
term-typing certificate was proved above.
\end{definition}
\lean{Graph/FocusRoundTrip}{12}{23}{Linear.FocusErasure.restore}
\lean{Graph/FocusRoundTrip}{25}{29}{Linear.FocusErasure.position_not_output}
\lean{Graph/FocusRoundTrip}{32}{34}{Linear.FocusErasure.reinsert}
The bijection's inverse laws follow by retained/fresh cases; the fresh carrier
has exactly one member. At ports, restoring host inclusion agrees with $j$,
and restoring the fresh insertion vertex gives $t$. Restoring the insertion's
source-redirection map agrees with $a$: the test $x=p$ is equivalent to
$j(x)=s$, by inclusion injectivity and $j(p)=s$.
\lean{Graph/FocusRoundTrip}{36}{38}{Linear.FocusErasure.restore_host}
\lean{Graph/FocusRoundTrip}{40}{41}{Linear.FocusErasure.restore_fresh}
\lean{Graph/FocusRoundTrip}{43}{55}{Linear.FocusErasure.restore_split}

\begin{theorem}[Focus-preserving round trip]
Reinserting after erasure is isomorphic to the original graph. The isomorphism
restores the focus edge, source, and target, preserves boundary order, and
preserves literal edge labels, including complete boxed bodies.
\end{theorem}
\begin{proof}
Use the restoration bijections. Retained internal labels are unchanged; the
fresh insertion vertex has term label, equal to the original label of $t$.
Retained edge labels are unchanged, and the new focus label equals the original
focus label by its certificate.
For a retained source list, restoration after insertion-redirection after
contraction acts as $a\circ c$, which is the identity on these sources because
they avoid $s$. For a retained target list, restoration acts as $j\circ c$,
which is the identity because these targets avoid $t$. Map these pointwise
identities over the lists. The new focus source restores to $j(p)=s$, and its
target restores to $t$, giving the original singleton incidence lists.
The fresh edge restores to $e_f$ by definition, and the preceding two port
identities explicitly identify its endpoints.
\end{proof}
\lean{Graph/FocusRoundTrip}{59}{104}{Linear.FocusErasure.reinsert_iso}
\lean{Graph/FocusRoundTrip}{107}{107}{Linear.FocusErasure.reinsert_focus_edge}
\lean{Graph/FocusRoundTrip}{110}{116}{Linear.FocusErasure.reinsert_focus_endpoints}

\begin{lemma}[Every closed focussed graph admits erasure]
The focus target of any graph $G:I\to\epsilon$ is internal. Therefore every
such focussed graph has a globally focus-free underlying graph, preserving
hole-freedom, and reinserting its mode recovers it up to the proved isomorphism.
\end{lemma}
\begin{proof}
The focus target occurs in a target list, so cannot be an input. There are no
outputs. The remaining port constructor is internal. Choose its internal name
with its equality certificate, and instantiate the preceding erasure,
invariant, and round-trip constructions. This use of choice supplies a witness
already proved to exist; it does not assume an erasure law.
\end{proof}
\lean{Graph/FocusRoundTrip}{127}{133}{Linear.Focussed.target_internal}
\lean{Graph/FocusRoundTrip}{135}{135}{Linear.Focussed.targetVertex}
\lean{Graph/FocusRoundTrip}{137}{137}{Linear.Focussed.targetVertex_spec}
\lean{Graph/FocusRoundTrip}{140}{141}{Linear.Focussed.underlying}
\lean{Graph/FocusRoundTrip}{143}{144}{Linear.Focussed.underlying_focusFree}
\lean{Graph/FocusRoundTrip}{146}{147}{Linear.Focussed.underlying_holeFree}
\lean{Graph/FocusRoundTrip}{151}{153}{Linear.Focussed.underlying_reinsert_iso}
The theorem proves the erasure-then-insertion direction for arbitrary closed
focussed graphs. The opposite insertion-then-erasure isomorphism is proved in
Section~\ref{sec:insertion-cancellation}, and naturality under source
isomorphisms in Section~\ref{sec:erasure-renaming}.

\Needspace{16\baselineskip}
\subsection{Focus-preserving isomorphisms and state recovery}
\begin{definition}[Focussed isomorphism]\label{term:focussed-isomorphism}
A focussed isomorphism is a literal-label graph isomorphism together with
equality of the two focus modes and an equation mapping the distinguished
focus edge to the distinguished focus edge. Identity, inverse, and composition
are obtained from the corresponding graph isomorphisms, using the edge
bijection's inverse laws and composition for the focus equations.
\end{definition}
\lean{Graph/StateSearch}{12}{15}{Linear.Focussed.Iso}
\lean{Graph/StateSearch}{20}{20}{Linear.Focussed.Iso.refl}
\lean{Graph/StateSearch}{22}{27}{Linear.Focussed.Iso.symm}
\lean{Graph/StateSearch}{29}{34}{Linear.Focussed.Iso.trans}

\begin{lemma}[Endpoints and observations are invariant]
A focussed isomorphism maps both focus endpoints correctly and preserves
initiality and finality.
\end{lemma}
\begin{proof}
Specialize the isomorphism's source-list equation to the focus edge and rewrite
both lists as singletons. Singleton injectivity identifies the endpoints.
Repeat for target lists. Port maps preserve input status; together with the
mode equality, this yields both boundary-predicate equivalences.
\end{proof}
\lean{Graph/StateSearch}{36}{39}{Linear.Focussed.Iso.source}
\lean{Graph/StateSearch}{41}{44}{Linear.Focussed.Iso.target}
\lean{Graph/StateSearch}{46}{49}{Linear.Focussed.Iso.initial_iff}
\lean{Graph/StateSearch}{51}{54}{Linear.Focussed.Iso.final_iff}

\begin{definition}[Extracted state coordinate]\label{term:state-coordinate}
For an arbitrary state $s$, erase its focus and take the typed position whose
mode is the old mode and whose vertex is the joined port $p$. Its typing is
the erased-position typing lemma. Re-realizing this coordinate in the globally
focus-free, hole-free underlying graph yields a state focus-preservingly
isomorphic to $s$, using the general round trip. Thus the input need not be a
canonical insertion or carry an assumed representation certificate.
\end{definition}
\lean{Graph/StateSearch}{63}{65}{Linear.State.coordinate}
\lean{Graph/StateSearch}{69}{74}{Linear.State.recover}
The extracted coordinate is initial or final exactly when $s$ is. Compose the
previous realization boundary equivalence with the focussed isomorphism's
boundary equivalence.
\lean{Graph/StateSearch}{76}{81}{Linear.State.coordinate_boundary}


\Needspace{20\baselineskip}
\subsection{Insertion followed by erasure}
\label{sec:insertion-cancellation}
Fix a focus-free graph $G:I\to[]$, a term vertex $v$, and a mode $m$.
Write $S=\operatorname{insert}(G,v,m)$ for the focussed graph constructed
in Section~\ref{sec:realization}. The preceding round trip started with an
arbitrary focussed graph. We now establish the opposite direction, including
the position at which the focus was inserted.

\begin{definition}[Deleting a fresh singleton]
For any carrier $A$, deleting the fresh member from $A+\mathbf1$ leaves a
carrier canonically bijective with $A$. The forward map sends a retained
$\mathsf{inl}(a)$ to $a$; the inverse sends $a$ to that retained element.
An element of the right summand cannot survive, since its singleton index
must be zero. Apply general erasure to the fresh focus and fresh target of
$S$ to define $E(S)$.
\end{definition}
The inverse laws hold by cases on the summand: the retained case is reflexive,
and the fresh case contradicts its non-deletion certificate.
\lean{Graph/InsertionRoundTrip}{10}{21}{Linear.FocusInsertion.deleteFresh}
\lean{Graph/InsertionRoundTrip}{27}{28}{Linear.FocusInsertion.erasedInsertion}

\begin{lemma}[Cancellation on ports]
Let $D$ be the port map induced by deletion of the fresh singleton, $j$ the
inclusion of retained ports, and $\kappa$ the collapse map of insertion.
Then $D(x)=\kappa(j(x))$. The joined vertex maps to $v$, and for every port
$w$ of $S$, $D(c(w))=\kappa(w)$, where $c$ is general focus contraction.
\end{lemma}
\begin{proof}
The first identity follows by cases on the retained port. An internal
right-summand case is impossible; input ports are fixed. For the joined
vertex, inclusion gives the focus source, namely the hosted copy of $v$;
collapse sends it to $v$. Finally, if $w$ is the fresh target, contraction
gives the joined vertex and both sides equal $v$. Otherwise inclusion after
contraction is $w$, and the first identity applies.
\end{proof}
\lean{Graph/InsertionRoundTrip}{31}{40}{Linear.FocusInsertion.deleteFresh_port}
\lean{Graph/InsertionRoundTrip}{43}{47}{Linear.FocusInsertion.erase_position}
\lean{Graph/InsertionRoundTrip}{50}{59}{Linear.FocusInsertion.erase_contract}

\begin{theorem}[Insertion--erasure cancellation]
\termdef{insertion-cancellation}
There is an isomorphism $E(S)\cong G$ preserving full labels and ordered
incidences. The choice-based operation $S.\mathsf{underlying}$ has such an
isomorphism too, and one can choose it to send the joined vertex to $v$.
\end{theorem}
\begin{proof}
Use deletion of the singleton on both internal vertices and edges. All
surviving labels are inherited. Sources were first split by insertion and
then contracted by erasure; the port cancellation identity reduces their
composite to collapse after split, which is the identity. Targets were
hosted by insertion, so the analogous composite is collapse after hosting,
also the identity. These pointwise identities give equality of ordered lists.

The target selected by general erasure is necessarily the fresh internal
vertex: both it and the fresh vertex equal the focus target, and the internal
port constructor is injective. Substituting this equality identifies the
general underlying graph with $E(S)$. For the position-sensitive assertion,
carry an arbitrary target together with its equality certificate, identify
it with the fresh target, and use the explicit isomorphism and joined-port
identity together. This keeps graph and coordinate recovery synchronized.
\end{proof}
\lean{Graph/InsertionRoundTrip}{62}{100}{Linear.FocusInsertion.erase_iso}
\lean{Graph/InsertionRoundTrip}{103}{104}{Linear.FocusInsertion.targetVertex_eq}
\lean{Graph/InsertionRoundTrip}{107}{111}{Linear.FocusInsertion.underlying_iso}
\lean{Graph/InsertionRoundTrip}{115}{129}{Linear.FocusInsertion.underlying_position_exists}

\Needspace{18\baselineskip}
\subsection{Search coordinates respect renaming}
\label{sec:search-renaming}
Let $a:G\cong H$ be a literal-label isomorphism and write $a_*$ for its port
map. Internal wire labels are preserved by the isomorphism's label field;
boundary labels are preserved because boundary positions are fixed.
\lean{Graph/SearchRenaming}{11}{16}{Linear.iso_vertexLabel}

\begin{definition}[Transport of operation views and coordinates]
Transport an operation view by applying the edge bijection to its edge and
$a_*$ to its source and both ordered lists. Keep its operation symbol.
Source and target certificates follow from the isomorphism's incidence laws;
arity follows because mapping a list preserves length. Transport a typed
coordinate $(m,v)$ to $(m,a_*v)$. Wire-label preservation supplies its typing
certificate.
\end{definition}
\lean{Graph/SearchRenaming}{21}{30}{Linear.SearchOperation.rename}
\lean{Graph/SearchRenaming}{86}{87}{Linear.TypedPosition.rename}

\begin{lemma}[Transport of the local rules]
Operation entry and argument return commute with transport. Every coordinate
search step transports forwards, and a step between transported coordinates
exists if and only if the original step exists.
\end{lemma}
\begin{proof}
Entry depends on whether the argument list is empty or has a first element;
mapping preserves these alternatives. Return depends on the optional element
at index $i+1$; lookup in a mapped list is the mapped lookup. These facts prove
the two local calculation identities.

Induct on the search rule. The variable rule uses the transported label and
singleton source certificate. The operation-entry case uses the entry identity.
For argument return, map the certified lookup at index $i$ and use the return
identity. This proves forward transport for all five rules. Apply the same
result to $a^{-1}$ for reflection: the composite port map is the identity,
and neither map changes the mode. The same inverse law proves that transporting
a typed coordinate by $a$ and back by $a^{-1}$ returns that coordinate; subtype
extensionality discards the irrelevant typing certificates.
\end{proof}
\lean{Graph/SearchRenaming}{32}{35}{Linear.SearchOperation.rename_enter}
\lean{Graph/SearchRenaming}{37}{40}{Linear.SearchOperation.rename_resume}
\lean{Graph/SearchRenaming}{47}{62}{Linear.SearchStep.rename}
\lean{Graph/SearchRenaming}{65}{78}{Linear.SearchStep.rename_iff}
\lean{Graph/SearchRenaming}{90}{95}{Linear.TypedPosition.rename_symm}

\Needspace{18\baselineskip}
\subsection{Insertion respects renaming}
Extend the vertex and edge bijections of $a$ by the identity on each fresh
singleton. The extended port map sends hosted vertices to hosted images and
fixes the fresh vertex. It also commutes with splitting: injectivity gives
$a_*w=a_*v$ exactly when $w=v$, so both splits select the fresh branch in
exactly the same case.
\lean{Graph/InsertionRenaming}{10}{12}{Linear.FocusInsertion.rename_host}
\lean{Graph/InsertionRenaming}{14}{15}{Linear.FocusInsertion.rename_fresh}
\lean{Graph/InsertionRenaming}{17}{25}{Linear.FocusInsertion.rename_split}

\begin{theorem}[Naturality of insertion]
Insertion at $(m,v)$ in $G$ is focus-preservingly isomorphic to insertion at
$(m,a_*v)$ in $H$. In particular, realizing a typed coordinate commutes with
its transport through $a$.
\end{theorem}
\begin{proof}
Use the two extended bijections. Retained labels are preserved by $a$;
fresh vertex and edge labels agree by construction. On retained edges,
commutation with splitting proves the source equation and commutation with
hosting proves the target equation, after applying the corresponding
incidence equation of $a$. The fresh edge has exactly the hosted insertion
position as source and the fresh vertex as target, so the two port identities
prove its equations. The mode is unchanged, and the fresh focus edge maps
to itself. Packaging these fields gives a focussed isomorphism. Applying
this construction to typed coordinates gives the realization result;
focus-free and hole-free certificates ensure both endpoints are states.
\end{proof}
\lean{Graph/InsertionRenaming}{28}{67}{Linear.FocusInsertion.rename_iso}
\lean{Graph/InsertionRenaming}{70}{77}{Linear.FocusInsertion.rename_focussed}
\lean{Graph/InsertionRenaming}{86}{89}{Linear.TypedPosition.realize_rename}

\Needspace{20\baselineskip}
\subsection{Exact characterization of search from realized states}
\label{sec:realized-search}
\begin{lemma}[Recovery of a realized coordinate]
If $p$ is a typed coordinate of a focus-free, hole-free graph $G$, there is
an isomorphism $a$ from the underlying graph of its realization to $G$
such that $a_*\operatorname{coordinate}(\operatorname{realize}(p))=p$.
\end{lemma}
\begin{proof}
Use the position-sensitive insertion--erasure isomorphism. It restores the
vertex, while the extracted mode is definitionally the original mode. Equality
of these two components gives equality of coordinates and then of the typed
subtypes.
\end{proof}
\lean{Graph/RealizedSearch}{14}{22}{Linear.TypedPosition.coordinate_recovery}

\begin{theorem}[Coordinate characterization]
\termdef{coordinate-characterization}
For every actual target state $t$,
\[
 \operatorname{realize}(p)\longrightarrow t
 \quad\Longleftrightarrow\quad
 \exists q,\quad p\longrightarrow_G q\ \land\
       \operatorname{realize}(q)\cong_f t.
\]
Here the left relation is erasure-based state search, the middle relation is
coordinate search on the original graph $G$, and $\cong_f$ preserves the
focus and mode. Both realizations use the same original graph $G$.
\end{theorem}
\begin{proof}
Choose the recovery isomorphism $a$ from the previous lemma. For the forward
direction, a state-search witness consists of a successor coordinate on the
erased graph and an isomorphism from its realization to $t$. Transport this
coordinate and its step through $a$. Coordinate recovery identifies the source
with $p$. Naturality of realization, inverted and composed with the supplied
endpoint isomorphism, identifies the new realization with $t$.

Conversely, transport $q$ and its step through $a^{-1}$. Inverse transport and
coordinate recovery identify the source with the extracted coordinate of
$\operatorname{realize}(p)$. Invert the realization isomorphism and compose it
with the given isomorphism to $t$. The transported coordinate, transported
step, and composite isomorphism are exactly the witnesses required by the
state-search definition.
\end{proof}
\lean{Graph/RealizedSearch}{26}{48}{Linear.TypedPosition.search_realize_iff}
Taking $t=\operatorname{realize}(q)$ and the identity endpoint isomorphism
shows that every coordinate step realizes an actual state-search step between
its two prescribed realizations. This strengthens the earlier endpoint-only
realization theorem. It does not identify search with contextual graph rewriting.
\lean{Graph/RealizedSearch}{52}{55}{Linear.TypedPosition.search_realize}

\Needspace{20\baselineskip}
\subsection{Erasure and state search respect renaming}
\label{sec:erasure-renaming}
Fix a focussed isomorphism $a:f\cong_f g$. For any bijection carrying $x$ to
$y$, restrict it to the complements of $x$ and $y$. Its forward and inverse
maps are the original maps with non-deletion certificates: injectivity prevents
a retained point from mapping to the deleted point, and the inverse law proves
the analogous claim backwards. Subtype extensionality lifts both inverse laws.
\lean{Graph/ErasureRenaming}{10}{15}{Linear.exceptOfEq}

The graph vertex bijection carries the selected internal target of $f$ to
that of $g$. Indeed, the focussed isomorphism identifies the targets, their
selection certificates make them internal, and the internal constructor is
injective. Restrict this vertex bijection at the selected targets and the edge
bijection at the focus edges. These give bijections of the erased carriers.
\lean{Graph/ErasureRenaming}{21}{24}{Linear.Focussed.Iso.targetVertex}
\lean{Graph/ErasureRenaming}{27}{28}{Linear.Focussed.Iso.erasedVertices}
\lean{Graph/ErasureRenaming}{31}{32}{Linear.Focussed.Iso.erasedEdges}

\begin{lemma}[Naturality of contraction]
Inclusion of retained ports commutes with the restricted vertex bijection.
The joined positions correspond, and contraction commutes with renaming.
\end{lemma}
\begin{proof}
Commutation with inclusion is a case analysis on ports. To compare joined
positions, apply the injective inclusion on the target side. The resulting
equation is precisely preservation of the focus source. For contraction,
separate the deleted target from other vertices. In the first case both sides
are the corresponding joined positions. In the second, injectivity ensures
that the image is not the target either. Apply inclusion on the target side:
both contraction--inclusion composites reduce to their original vertices,
and the inclusion commutation identity proves equality.
\end{proof}
\lean{Graph/ErasureRenaming}{34}{37}{Linear.Focussed.Iso.erased_lift}
\lean{Graph/ErasureRenaming}{40}{45}{Linear.Focussed.Iso.erased_position}
\lean{Graph/ErasureRenaming}{48}{63}{Linear.Focussed.Iso.erased_contract}

\begin{theorem}[Naturality of erasure and state search]
The restricted bijections form an isomorphism of the erased graphs, carrying
the extracted coordinate of the first state to that of the second. Search
is invariant under focus-preserving isomorphisms in both arguments.
\end{theorem}
\begin{proof}
Retained labels inherit the original isomorphism's label equations. For each
incidence list, use its original equation and commute contraction past the
port map pointwise. This proves the erased graph isomorphism. Its coordinate
equation follows from preservation of the mode and the joined position.

Given a search step from the first state, transport its successor coordinate
through the erased isomorphism. The coordinate equation identifies the new
source, and search transport supplies the step. Invert the realization
isomorphism and compose with the old endpoint isomorphism to retain the same
actual target. This proves source invariance. Combine it with the earlier
target invariance to obtain invariance in both arguments; use the inverse
isomorphisms for the reverse implication.
\end{proof}
\lean{Graph/ErasureRenaming}{66}{86}{Linear.Focussed.Iso.erased}
\lean{Graph/ErasureRenaming}{94}{99}{Linear.State.coordinate_iso}
\lean{Graph/ErasureRenaming}{102}{112}{Linear.State.search_source_iso}
\lean{Graph/ErasureRenaming}{115}{122}{Linear.State.search_iso_iff}

\Needspace{20\baselineskip}
\subsection{Search on isomorphism classes of states}
\label{sec:state-quotient}
\begin{definition}[State isomorphism class]
\termdef{state-isomorphism-class}
Two states are equivalent when there exists a focus-preserving literal-label
isomorphism between them. Reflexivity, symmetry, and transitivity follow
from the corresponding isomorphism constructions. A state class is the
quotient by this equivalence; write $[s]$ for the class of $s$.
\end{definition}
This quotient preserves full edge labels literally, including box bodies.
It does not identify independently renamed contents of nested boxes.
Section~\ref{sec:recursive-state-search} constructs the recursive state quotient
and proves projection and exact-length search lifting between the two quotients.
\lean{Graph/StateQuotient}{10}{13}{Linear.stateSetoid}
\lean{Graph/StateQuotient}{16}{16}{Linear.StateClass}
\lean{Graph/StateQuotient}{22}{22}{Linear.StateClass.ofState}

\begin{definition}[Quotient search and observations]
\termdef{quotient-search}
Define $[s]\longrightarrow[t]$ by state search $s\longrightarrow t$.
The preceding two-argument invariance proves independence of representatives,
so quotient elimination defines a proposition on classes. Initiality and
finality descend in the same way, using their focussed-isomorphism invariance.
\end{definition}
On supplied representatives, all three descended propositions agree with the
original propositions by definition. The representative equations are therefore
reflexive equivalences, not extra assumptions.
\lean{Graph/StateQuotient}{25}{27}{Linear.StateClass.Search}
\lean{Graph/StateQuotient}{29}{30}{Linear.StateClass.search_iff}
\lean{Graph/StateQuotient}{39}{41}{Linear.StateClass.Initial}
\lean{Graph/StateQuotient}{44}{46}{Linear.StateClass.Final}
\lean{Graph/StateQuotient}{48}{48}{Linear.StateClass.initial_iff}
\lean{Graph/StateQuotient}{50}{50}{Linear.StateClass.final_iff}

\begin{theorem}[Deterministic quotient search and its boundaries]
If $S\longrightarrow T$ and $S\longrightarrow U$, then $T=U$.
A final class has no search successor, and no search step reaches an initial
class.
\end{theorem}
\begin{proof}
Eliminate the three quotient arguments to representatives $s,t,u$.
The two representative steps have a common source, so state-level determinism
gives a focussed isomorphism $t\cong_f u$. Quotient soundness identifies their
classes. For each boundary theorem eliminate its two class arguments;
the goal becomes the already proved state-level boundary law.
\end{proof}
\lean{Graph/StateQuotient}{33}{36}{Linear.StateClass.search_deterministic}
\lean{Graph/StateQuotient}{53}{56}{Linear.StateClass.final_no_search}
\lean{Graph/StateQuotient}{59}{62}{Linear.StateClass.no_search_initial}
These are results for the search relation alone. Substitution, computation,
rootedness and refocusing remain necessary for
the full concrete UAM argument.

\Needspace{18\baselineskip}
\subsection{Concrete quotient-search regression}
The passive nullary value step now runs from the previously defined
\texttt{initialValue} to the specific \texttt{finalValue}, using the general
coordinate-to-state theorem. Erase and reconstruct \texttt{initialValue} to
define \texttt{recoveredValue}. Its carriers are rebuilt by deletion and
insertion, but state recovery supplies a focussed isomorphism to the original.
Source invariance therefore gives a step from this reconstructed state to
\texttt{finalValue}. Quotient soundness identifies its class with the original
initial class. Finally, quotient determinism compares any successor of that
initial class with the concrete value step: it must equal the class of
\texttt{finalValue}.
\lean{Graph/QuotientSearchExamples}{9}{12}{Linear.Examples.value_search_final}
\lean{Graph/QuotientSearchExamples}{15}{17}{Linear.Examples.recoveredValue}
\lean{Graph/QuotientSearchExamples}{19}{20}{Linear.Examples.recovered_value_search}
\lean{Graph/QuotientSearchExamples}{22}{23}{Linear.Examples.recovered_value_class}
\lean{Graph/QuotientSearchExamples}{26}{29}{Linear.Examples.value_class_successor}

\Needspace{20\baselineskip}
\subsection{Recursive renaming of focussed states}
\label{sec:recursive-state-search}
The preceding quotient fixes box-body labels literally. We now extend search
to the recursive graph isomorphisms of Section~23, allowing independent changes
of carrier names inside every nested box. All constructions stay at one fixed
finite depth bound. Compatibility across depth bounds and with the paper's
set-based representation remains a separate obligation.

\begin{lemma}[Primitive labels and focus freedom]
Recursive label equivalence preserves each specified primitive symbol exactly.
It also preserves and reflects absence of focuses throughout a label, and
recursive graph isomorphisms preserve and reflect global focus freedom.
\end{lemma}
\begin{proof}
At depth zero label equivalence is equality. At positive depth a primitive can
only relate to the same primitive: the primitive--box alternatives are false.
This proves the first assertion, including the variable and operation symbols
needed by search.

For focus freedom, induct on depth. Primitive cases again reduce to equality;
mixed cases are impossible. For related boxes, take the supplied recursive
isomorphism of their bodies. In the forward direction, choose the inverse
image of each target edge under its edge bijection, apply the induction
hypothesis to the related labels, and cancel the bijection with its inverse.
In the reverse direction, apply the induction hypothesis to the image of each
source edge. Applying this label result edgewise, with the same inverse-image
argument, gives the graph-level equivalence.
\end{proof}
\lean{Graph/DeepFocus}{10}{19}{Linear.labelEquiv_primitive}
\lean{Graph/DeepFocus}{22}{43}{Linear.labelEquiv_focusFree}
\lean{Graph/DeepFocus}{45}{52}{Linear.DeepIso.focusFree_iff}

\begin{definition}[Focussed recursive isomorphism]
\termdef{focussed-recursive-isomorphism}
A focussed recursive isomorphism consists of a recursive graph isomorphism,
equality of the modes, and an equation carrying the distinguished focus edge
to the distinguished focus edge. Literal-label focussed isomorphisms embed
by treating equality of labels as recursive label equivalence.
\end{definition}
Identity uses recursive identity; inversion uses the inverse graph bijections
and the reversed mode equation; composition composes the graph maps and both
focus equations. Thus these constructions retain the distinguished data,
not only the unpointed graph.
\lean{Graph/DeepFocus}{57}{60}{Linear.Focussed.DeepIso}
\lean{Graph/DeepFocus}{65}{65}{Linear.Focussed.DeepIso.refl}
\lean{Graph/DeepFocus}{67}{72}{Linear.Focussed.DeepIso.symm}
\lean{Graph/DeepFocus}{74}{79}{Linear.Focussed.DeepIso.trans}
\lean{Graph/DeepFocus}{82}{82}{Linear.Focussed.DeepIso.ofIso}

The focus source and target correspond: apply the ordered incidence equations
to the distinguished edge and use its singleton source and target lists.
The port maps preserve input status, so these endpoint equations and mode
equality preserve and reflect both initiality and finality.
\lean{Graph/DeepFocus}{84}{87}{Linear.Focussed.DeepIso.source}
\lean{Graph/DeepFocus}{89}{92}{Linear.Focussed.DeepIso.target}
\lean{Graph/DeepFocus}{94}{97}{Linear.Focussed.DeepIso.initial_iff}
\lean{Graph/DeepFocus}{99}{102}{Linear.Focussed.DeepIso.final_iff}

\Needspace{18\baselineskip}
\subsection{Search and insertion under recursive renaming}
Let $a:G\cong_d H$ be a recursive graph isomorphism and $a_*$ its outer port
map. Wire labels still agree: internal labels agree by the graph-isomorphism
field, and boundary labels agree because boundary positions are fixed.
\lean{Graph/DeepSearchRenaming}{11}{16}{Linear.deepIso_vertexLabel}

Transport an operation view by the outer edge bijection and port map, retaining
its operation symbol. The primitive-label lemma proves the operation-label
certificate even when other edges carry independently renamed boxes. Ordered
incidence certificates follow from the graph equations, and arity follows
from preservation of list length. Transport a typed coordinate by keeping
its mode and applying $a_*$ to its vertex; wire-label preservation proves
the new typing certificate.
\lean{Graph/DeepSearchRenaming}{21}{30}{Linear.SearchOperation.deepRename}
\lean{Graph/DeepSearchRenaming}{86}{87}{Linear.TypedPosition.deepRename}

\begin{lemma}[Recursive transport of coordinate search]
Operation entry and return commute with recursive transport. A search step
between transported coordinates exists exactly when the original step does.
Transporting a typed coordinate forwards and then backwards returns it.
\end{lemma}
\begin{proof}
Entry examines the empty/nonempty argument-list alternatives, which mapping
preserves. Return uses optional lookup at index $i+1$, which commutes with
mapping. For search, induct on the rule: the variable case uses exact
preservation of the variable label and singleton source; entry uses its
commutation equation; return also maps the certified lookup at index $i$.
This proves forward transport. Apply it to the inverse recursive isomorphism
and cancel its port bijection to reflect steps. The same cancellation proves
the typed-coordinate inverse law, by equality of the mode and vertex and
then subtype extensionality.
\end{proof}
\lean{Graph/DeepSearchRenaming}{32}{35}{Linear.SearchOperation.deepRename_enter}
\lean{Graph/DeepSearchRenaming}{37}{40}{Linear.SearchOperation.deepRename_resume}
\lean{Graph/DeepSearchRenaming}{47}{62}{Linear.SearchStep.deepRename}
\lean{Graph/DeepSearchRenaming}{65}{78}{Linear.SearchStep.deepRename_iff}
\lean{Graph/DeepSearchRenaming}{90}{95}{Linear.TypedPosition.deepRename_symm}

\begin{theorem}[Recursive naturality of insertion]
Inserting a focus at $(m,v)$ in $G$ is recursively focus-preservingly
isomorphic to inserting it at $(m,a_*v)$ in $H$. Realization of typed
coordinates respects this isomorphism.
\end{theorem}
\begin{proof}
Extend the outer vertex and edge bijections by fixing the inserted singleton
names. Hosted vertices map to hosted images, and the fresh vertex is fixed.
Splitting commutes with this map because injectivity says $a_*w=a_*v$ exactly
when $w=v$; hence both sides take the fresh branch in the same case.

Retained edge labels inherit their recursive equivalence from $a$. The fresh
focus labels agree literally, so reflexivity supplies their recursive
equivalence. Retained source lists use commutation with splitting, retained
target lists use commutation with hosting, and the fresh focus equations use
the hosted and fresh port identities directly. The mode and fresh focus edge
are fixed. These fields give the focussed recursive isomorphism, and the
state-realization specialization follows by supplying its focus-free and
hole-free certificates.
\end{proof}
\lean{Graph/DeepInsertionRenaming}{10}{12}{Linear.FocusInsertion.deepRename_host}
\lean{Graph/DeepInsertionRenaming}{14}{15}{Linear.FocusInsertion.deepRename_fresh}
\lean{Graph/DeepInsertionRenaming}{17}{25}{Linear.FocusInsertion.deepRename_split}
\lean{Graph/DeepInsertionRenaming}{28}{67}{Linear.FocusInsertion.deepRename_iso}
\lean{Graph/DeepInsertionRenaming}{70}{77}{Linear.FocusInsertion.deepRename_focussed}
\lean{Graph/DeepInsertionRenaming}{86}{89}{Linear.TypedPosition.realize_deepRename}

\Needspace{18\baselineskip}
\subsection{Erasure under recursive renaming}
Let $a:f\cong_{f,d}g$ be a focussed recursive isomorphism. Its outer vertex
bijection carries the selected internal target to the selected internal target:
use target preservation, rewrite both selection certificates, and cancel the
internal-port constructor. Restrict this bijection to the undeleted vertices
and restrict the edge bijection to the non-focus edges, using the existing
complement-bijection construction.
\lean{Graph/DeepErasureRenaming}{13}{16}{Linear.Focussed.DeepIso.targetVertex}
\lean{Graph/DeepErasureRenaming}{19}{20}{Linear.Focussed.DeepIso.erasedVertices}
\lean{Graph/DeepErasureRenaming}{23}{24}{Linear.Focussed.DeepIso.erasedEdges}

\begin{lemma}[Retained ports and contraction]
Inclusion of retained ports commutes with these restricted maps. The joined
vertices correspond, and contraction commutes with recursive renaming.
\end{lemma}
\begin{proof}
The inclusion equation holds by cases on ports. Apply the injective inclusion
on the target side to compare joined vertices: their images are the focus
sources, already identified by $a$. For contraction, if the port is the
deleted target, both sides are corresponding joined vertices. Otherwise
injectivity ensures that its image is also not deleted. Apply inclusion and
use the contraction--inclusion identities on both sides to reduce to the
inclusion commutation equation.
\end{proof}
\lean{Graph/DeepErasureRenaming}{26}{29}{Linear.Focussed.DeepIso.erased_lift}
\lean{Graph/DeepErasureRenaming}{32}{37}{Linear.Focussed.DeepIso.erased_position}
\lean{Graph/DeepErasureRenaming}{40}{55}{Linear.Focussed.DeepIso.erased_contract}

The restricted maps form a recursive isomorphism of the erased graphs.
Internal wire labels and retained recursive edge-label relations are inherited
from $a$. For sources and targets, rewrite the original incidence equation
and commute contraction with the port map pointwise along its ordered list.
Nothing requires literal equality of boxed bodies.
\lean{Graph/DeepErasureRenaming}{58}{78}{Linear.Focussed.DeepIso.erased}

\Needspace{20\baselineskip}
\subsection{State search modulo recursive endpoint renaming}
\begin{definition}[Recursively renamed state search]
\termdef{recursive-state-search}
Write $s\longrightarrow_d t$ when a typed coordinate $q$ on the erased
underlying graph of $s$ satisfies the coordinate search rule from the extracted
coordinate of $s$, and the realization of $q$ is recursively
focus-preservingly isomorphic to $t$.
\end{definition}
Taking the actual realization as target gives a step by the identity
isomorphism. Composing the supplied endpoint isomorphism with another proves
closure under recursive target renaming.
\lean{Graph/DeepStateSearch}{12}{16}{Linear.State.DeepSearch}
\lean{Graph/DeepStateSearch}{19}{24}{Linear.State.deepSearch_realize}
\lean{Graph/DeepStateSearch}{27}{30}{Linear.State.deepSearch_target_iso}

\begin{theorem}[Recursive search invariance, determinism, and boundaries]
Two recursive-search successors of the same state are focussed recursively
isomorphic. Recursive search respects recursive isomorphisms in both arguments.
Final states have no such search step, and no such step reaches an initial
state.
\end{theorem}
\begin{proof}
For determinism, compare the two witnessing coordinate steps on the identical
erased source graph. Coordinate determinism identifies their results; subtype
extensionality identifies their typed coordinates. Compose the inverse of
one endpoint isomorphism with the other.

For source invariance, erasure naturality identifies the extracted coordinates:
the mode equation gives the first component and the joined-position equation
the second. Transport the successor coordinate and step through the erased
recursive isomorphism. Naturality of realization identifies the transported
realization with the old one, so compose its inverse with the old endpoint
isomorphism. Combine this with target closure; apply inverse isomorphisms to
obtain the converse and hence invariance in both arguments.

For a final source, its extracted coordinate is backtrack at an input, which
has no coordinate search successor. For an initial target, transport its
initiality backwards through the endpoint isomorphism and then through the
realization boundary law. The successor coordinate would be search at an
input, excluded by the coordinate no-initial-target theorem.
\end{proof}
\lean{Graph/DeepStateSearch}{33}{39}{Linear.State.deepSearch_deterministic}
\lean{Graph/DeepStateSearch}{62}{67}{Linear.State.coordinate_deepIso}
\lean{Graph/DeepStateSearch}{70}{80}{Linear.State.deepSearch_source_iso}
\lean{Graph/DeepStateSearch}{83}{90}{Linear.State.deepSearch_iso_iff}
\lean{Graph/DeepStateSearch}{42}{48}{Linear.State.final_no_deepSearch}
\lean{Graph/DeepStateSearch}{51}{59}{Linear.State.no_deepSearch_initial}

\begin{theorem}[Exact relation to literal-label state search]
Every literal-label state-search step is a recursive step. More precisely,
\[
 s\longrightarrow_d t\quad\Longleftrightarrow\quad
 \exists u,\ s\longrightarrow u\ \land\ u\cong_{f,d}t.
\]
\end{theorem}
\begin{proof}
A literal endpoint isomorphism embeds into a recursive one, giving the first
claim. For the forward implication of the equivalence, choose $u$ to be the
actual realization of the witnessing successor coordinate. This is a literal
state-search successor by identity, and the supplied recursive endpoint
isomorphism identifies it with $t$. Conversely, embed the literal step and
compose its endpoint witness with the supplied recursive isomorphism.
\end{proof}
\lean{Graph/DeepStateSearch}{93}{96}{Linear.State.search_deepSearch}
\lean{Graph/DeepStateSearch}{99}{106}{Linear.State.deepSearch_iff}
Thus the new relation changes the permitted endpoint representatives; it does
not introduce another local search rule. It is not a claim that literal-label
search between two fixed representatives is itself invariant under arbitrary
changes of box-body labels.

\Needspace{18\baselineskip}
\subsection{Recursive state classes and their search relation}
\begin{definition}[Recursive state class]
\termdef{recursive-state-class}
Quotient states by existence of a focussed recursive isomorphism. The identity,
inverse, and composition constructions prove this is an equivalence relation.
Write $[s]_d$ for the class of $s$. Descend recursive search to pairs of classes
using its proved two-argument invariance. Descend initiality and finality using
their recursive-isomorphism invariance.
\end{definition}
\lean{Graph/DeepStateQuotient}{10}{13}{Linear.deepStateSetoid}
\lean{Graph/DeepStateQuotient}{16}{16}{Linear.DeepStateClass}
\lean{Graph/DeepStateQuotient}{22}{22}{Linear.DeepStateClass.ofState}
\lean{Graph/DeepStateQuotient}{25}{27}{Linear.DeepStateClass.Search}
\lean{Graph/DeepStateQuotient}{39}{41}{Linear.DeepStateClass.Initial}
\lean{Graph/DeepStateQuotient}{44}{46}{Linear.DeepStateClass.Final}
Each descended proposition reduces to its representative proposition by
definition, giving the three representative equivalences.
\lean{Graph/DeepStateQuotient}{29}{30}{Linear.DeepStateClass.search_iff}
\lean{Graph/DeepStateQuotient}{48}{48}{Linear.DeepStateClass.initial_iff}
\lean{Graph/DeepStateQuotient}{50}{50}{Linear.DeepStateClass.final_iff}

\begin{theorem}[Deterministic search on recursive classes]
Recursive quotient search has at most one successor class. Final classes have
no search successor, and no search step reaches an initial class.
\end{theorem}
\begin{proof}
Eliminate the quotient arguments to representatives. Two steps from a common
representative have recursively isomorphic targets by state-level determinism;
quotient soundness equates their classes. Each boundary assertion reduces to
the corresponding representative-level boundary theorem.
\end{proof}
\lean{Graph/DeepStateQuotient}{33}{36}{Linear.DeepStateClass.search_deterministic}
\lean{Graph/DeepStateQuotient}{53}{56}{Linear.DeepStateClass.final_no_search}
\lean{Graph/DeepStateQuotient}{59}{62}{Linear.DeepStateClass.no_search_initial}

\Needspace{20\baselineskip}
\subsection{Projection and lifting between the two state quotients}
\begin{definition}[Forgetting literal body representatives]
\termdef{state-quotient-projection}
Define $\pi([s])=[s]_d$, mapping literal-label state classes to recursive state
classes. This is well-defined because every literal focussed isomorphism
embeds into a recursive one. On a specified representative the equation is
reflexive. The map is surjective: choose the literal class of any representative
of a recursive class.
\end{definition}
\lean{Graph/StateQuotientBridge}{10}{12}{Linear.StateClass.toDeep}
\lean{Graph/StateQuotientBridge}{14}{15}{Linear.StateClass.toDeep_ofState}
\lean{Graph/StateQuotientBridge}{18}{20}{Linear.StateClass.toDeep_surjective}

\begin{theorem}[Search projection and lifting]
A literal-class step $S\longrightarrow U$ gives
$\pi(S)\longrightarrow_d\pi(U)$. Conversely, for every specified literal
class $S$, a step $\pi(S)\longrightarrow_d T$ lifts to some step
$S\longrightarrow U$ with $\pi(U)=T$. Consequently,
\[
 \pi(S)\longrightarrow_d T\quad\Longleftrightarrow\quad
 \exists U,\ S\longrightarrow U\ \land\ \pi(U)=T.
\]
\end{theorem}
\begin{proof}
For projection, eliminate both literal classes to representatives and embed
the literal state-search step. For lifting, eliminate the source literal class
and target recursive class. The exact endpoint characterization supplies a
literal successor $u$ of the fixed source representative together with a
recursive isomorphism to the chosen target representative. Its literal class
is the desired $U$, and quotient soundness proves the projection equation.
The displayed equivalence combines projection and lifting.
\end{proof}
\lean{Graph/StateQuotientBridge}{23}{26}{Linear.StateClass.search_toDeep}
\lean{Graph/StateQuotientBridge}{30}{36}{Linear.StateClass.search_lift}
\lean{Graph/StateQuotientBridge}{39}{46}{Linear.StateClass.search_toDeep_iff}

Initiality and finality agree before and after projection, by quotient
elimination and their representative definitions. Absence of search successors
also agrees: a literal successor projects to a recursive successor, and a
recursive successor lifts to a literal successor. Taking the contrapositive
in each direction proves the normal-form equivalence.
\lean{Graph/StateQuotientBridge}{48}{50}{Linear.StateClass.initial_toDeep}
\lean{Graph/StateQuotientBridge}{52}{54}{Linear.StateClass.final_toDeep}
\lean{Graph/StateQuotientBridge}{57}{64}{Linear.StateClass.no_search_toDeep}

\Needspace{20\baselineskip}
\subsection{Exact-length executions across the quotients}
\begin{definition}[Search-only machines]
\termdef{search-only-machine}
On each state quotient, form a \texttt{Machine} whose step relation is its
search relation and whose final predicate is its boundary finality predicate.
The search boundary theorem proves the required final-normal law. Its
search-determinism theorem proves determinism of this machine.
\end{definition}
These machines deliberately contain only search. They are not the full UAM:
substitution and computation are absent. Their finite executions use the
already defined \texttt{Machine.Steps} relation, so the length counts search
steps exactly, without introducing a second execution datatype.
\lean{Graph/SearchQuotientExecutions}{10}{13}{Linear.StateClass.searchMachine}
\lean{Graph/SearchQuotientExecutions}{16}{19}{Linear.DeepStateClass.searchMachine}
\lean{Graph/SearchQuotientExecutions}{21}{24}{Linear.StateClass.searchMachine_deterministic}
\lean{Graph/SearchQuotientExecutions}{26}{29}{Linear.DeepStateClass.searchMachine_deterministic}

\begin{theorem}[Exact-length projection and lifting]
For every $n\in\mathbb N$, literal executions project with length unchanged,
and every recursive execution of length $n$ from $\pi(S)$ lifts from the
specified literal class $S$ with length unchanged. Thus
\[
 \pi(S)\longrightarrow_d^{n} T\quad\Longleftrightarrow\quad
 \exists U,\ S\longrightarrow^{n}U\ \land\ \pi(U)=T.
\]
\end{theorem}
\begin{proof}
Projection is induction on the execution: the zero-step execution remains
zero, and a step followed by a shorter run maps to its projected step followed
by the projected shorter run.

For lifting, induct on $n$, allowing arbitrary source and target classes.
At zero, the execution identifies the target with $\pi(S)$, so take $U=S$.
At $n+1$, lift the first recursive step from the specified source to a literal
successor $V$. Rewrite the recursive tail's source using $\pi(V)$, apply the
induction hypothesis from $V$, and prepend the lifted first step. The endpoint
equality and the exact length are retained. The equivalence combines both
constructions.
\end{proof}
\lean{Graph/SearchQuotientExecutions}{35}{40}{Linear.StateClass.steps_toDeep}
\lean{Graph/SearchQuotientExecutions}{43}{56}{Linear.StateClass.steps_lift}
\lean{Graph/SearchQuotientExecutions}{59}{65}{Linear.StateClass.steps_toDeep_iff}

\begin{theorem}[Exact-length search termination]
For each prescribed $n$, a successful search run of length $n$ exists from
$S$ if and only if one exists from $\pi(S)$.
\end{theorem}
\begin{proof}
Project or lift the run using the previous theorem. In the lifting direction,
rewrite finality along the endpoint projection equality and use reflection
of finality. In the projection direction, use preservation of finality.
\end{proof}
\lean{Graph/SearchQuotientExecutions}{68}{76}{Linear.StateClass.final_run_toDeep_iff}

\Needspace{18\baselineskip}
\subsection{Recursive-renaming and execution regression proofs}
The concrete two-hole box is globally focus-free because every edge in its
body is a hole rather than a focus. The established recursive isomorphism to
the box with its body's edge names exchanged transports this focus freedom.
These two boxed graphs were already proved to admit no literal-label
isomorphism, so the test exercises genuinely recursive label comparison.
\lean{Graph/DeepSearchExamples}{10}{12}{Linear.Examples.boxedTwoHoles_focusFree}
\lean{Graph/DeepSearchExamples}{16}{17}{Linear.Examples.renamedBoxedTwoHoles_focusFree}

The earlier concrete value step embeds into recursive search. Prepending that
step to the zero-step execution gives an execution of exactly one step from
its initial recursive class to its final recursive class. To test lifting from
a reconstructed source, identify its literal class with the original initial
class, use this recursive execution, and apply the exact-length termination
equivalence backwards. The result is a successful literal search run of exactly
one step from the reconstructed source. Finally, the final value's boundary
finality and the recursive no-search-from-final theorem exclude every recursive
search successor.
\lean{Graph/DeepSearchExamples}{20}{23}{Linear.Examples.value_deep_run}
\lean{Graph/DeepSearchExamples}{27}{33}{Linear.Examples.recovered_value_lift}
\lean{Graph/DeepSearchExamples}{36}{39}{Linear.Examples.final_value_deep_normal}

\Needspace{22\baselineskip}
\subsection{Canonical roots of concrete states}
\label{sec:concrete-roots}
This milestone connects the existing search semantics to Appendix~C.
We construct the paper's rooted predicate, prove that search preserves it,
and establish concrete versions of Lemmas~23 and~26(1). These results use
our erasure-based search semantics on recursive state classes; identifying that
semantics with the paper's contextual local rewrites remains an adequacy
obligation. The proofs below do not assume the abstract specimen locality laws.

\begin{definition}[Canonical initial attachment]
\termdef{concrete-root}
For a concrete state $s$, erase its focus, attach a searching focus at the
sole input of the resulting graph, and call the resulting state $r(s)$.
Focus freedom and hole freedom of erasure supply the certificates required
for realization. Its focus is initial by the initial-realization theorem.
\end{definition}
\lean{Graph/StateRoots}{12}{14}{Linear.State.root}
\lean{Graph/StateRoots}{16}{17}{Linear.State.root_initial}

\begin{lemma}[Roots respect renaming and initiality]
Recursively isomorphic underlying graphs give recursively focus-preservingly
isomorphic roots. In particular, a focussed recursive isomorphism between
states gives one between their roots. If $s$ is already initial, $r(s)$ is
recursively focus-preservingly isomorphic to $s$.
\end{lemma}
\begin{proof}
The coordinate ``search at input zero'' is fixed by every boundary-preserving
port map. Apply recursive naturality of realization to the given underlying
isomorphism. For focussed state isomorphisms, first use erasure naturality.

For an initial state, the extracted mode is search and its extracted vertex
is an input. The state interface has exactly one input, so that vertex is
input zero. Thus the extracted typed coordinate equals the initial coordinate,
by component equality and subtype extensionality. Rewrite the general recovery
isomorphism using this coordinate equality, and embed the literal-label
isomorphism into a recursive one.
\end{proof}
\lean{Graph/StateRoots}{20}{25}{Linear.State.root_of_underlying}
\lean{Graph/StateRoots}{27}{28}{Linear.State.root_iso}
\lean{Graph/StateRoots}{31}{40}{Linear.State.coordinate_initial}
\lean{Graph/StateRoots}{43}{47}{Linear.State.root_initial_iso}

\begin{lemma}[Search preserves the underlying graph]
If $s\longrightarrow_d t$, their erased underlying graphs are recursively
isomorphic.
\end{lemma}
\begin{proof}
Let $q$ be the witnessing successor coordinate on the erased graph of $s$.
Insertion--erasure cancellation identifies the underlying graph of
$\operatorname{realize}(q)$ with the underlying graph of $s$. Invert that
isomorphism and compose it with erasure of the supplied recursive endpoint
isomorphism from the realization to $t$.
\end{proof}
\lean{Graph/StateRoots}{51}{56}{Linear.State.deepSearch_underlying}

Descend $r$ to recursive state classes using root isomorphism invariance.
Its value on a representative is its representative root class, by definition.
Every root class is initial. An initial class is fixed by $r$, using quotient
soundness and the initial-state recovery lemma. In particular $r(r(S))=r(S)$.
Search preserves $r(S)$ because it preserves the underlying graph; induction
on a search execution then proves preservation of the root along the whole run.
\lean{Graph/StateRoots}{63}{65}{Linear.DeepStateClass.root}
\lean{Graph/StateRoots}{67}{67}{Linear.DeepStateClass.root_ofState}
\lean{Graph/StateRoots}{69}{71}{Linear.DeepStateClass.root_initial}
\lean{Graph/StateRoots}{73}{75}{Linear.DeepStateClass.root_eq_of_initial}
\lean{Graph/StateRoots}{77}{78}{Linear.DeepStateClass.root_idempotent}
\lean{Graph/StateRoots}{81}{86}{Linear.DeepStateClass.search_root}
\lean{Graph/StateRoots}{88}{92}{Linear.DeepStateClass.steps_root}

\Needspace{20\baselineskip}
\subsection{Concrete rootedness and its search invariant}
\begin{definition}[Concrete rooted state]
\termdef{concrete-rootedness}
A recursive state class $S$ is rooted when there is a finite search-only
execution from $r(S)$ to $S$:
\[
 \mathsf{Rooted}(S)\quad\Longleftrightarrow\quad
 \exists n\in\mathbb N,\ r(S)\longrightarrow_d^n S.
\]
A named state is rooted when its recursive isomorphism class is rooted.
This definition depends on the operation classification used by search.
\end{definition}
\lean{Graph/StateRoots}{95}{96}{Linear.DeepStateClass.Rooted}
\lean{Graph/RootedStatePaths}{96}{97}{Linear.State.Rooted}

\begin{theorem}[Initial states are rooted; search preserves rootedness]
Every initial class is rooted, every canonical root is rooted, and search
steps and finite search executions preserve rootedness.
\end{theorem}
\begin{proof}
An initial class equals its root, so the zero-step run witnesses rootedness.
Every root is initial, giving the second claim. For a step $S\longrightarrow_d T$,
append that step to the witnessing run from $r(S)$ to $S$. Search preservation
of roots identifies $r(S)$ with $r(T)$, so this is the required run to $T$.
Induction on a finite execution gives the final claim.
\end{proof}
\lean{Graph/StateRoots}{98}{102}{Linear.DeepStateClass.initial_rooted}
\lean{Graph/StateRoots}{104}{105}{Linear.DeepStateClass.root_rooted}
\lean{Graph/StateRoots}{108}{113}{Linear.DeepStateClass.rooted_search}
\lean{Graph/StateRoots}{115}{119}{Linear.DeepStateClass.rooted_steps}

\begin{theorem}[Equivalent initial-run characterization]
A class $S$ is rooted exactly when some initial class has a finite search run
to $S$. The initial class in any such run must equal $r(S)$.
\end{theorem}
\begin{proof}
The forward direction chooses $r(S)$, which is initial. Conversely, suppose
$I$ is initial and reaches $S$. Initiality gives $I=r(I)$, and preservation
of roots along the run gives $r(I)=r(S)$. Rewrite the run's source by these
equalities to obtain the rootedness witness.
\end{proof}
\lean{Graph/ConcreteRefocusing}{10}{17}{Linear.DeepStateClass.rooted_iff_initial_run}

\Needspace{20\baselineskip}
\subsection{Operation reachability under coordinate search}
\begin{definition}[Operation-reachable vertex]
\termdef{operation-reachable}
A vertex $v$ is operation-reachable if there are an input index $i$ and an
edge list $E$ with an operation trace from $v$ backwards to input $i$.
Equivalently, reversing $E$ gives an operation path from that input to $v$.
An input is reachable using the empty trace. Each certified operation view
has an actual operation label, by its label certificate and the definition
of primitive operation labels at each depth.
\end{definition}
\lean{Graph/SearchReachability}{10}{11}{Linear.OperationReachable}
\lean{Graph/SearchReachability}{13}{14}{Linear.OperationReachable.input}
\lean{Graph/SearchReachability}{16}{18}{Linear.SearchOperation.isOperation}

\begin{lemma}[Argument and source reachability]
If the source of an operation is reachable, every term argument is reachable.
Conversely, if one of its arguments is reachable, its source is reachable.
\end{lemma}
\begin{proof}
Forwards, prepend the operation edge to the reverse trace from its source;
its certified argument membership supplies the target incidence. Backwards,
inspect a witnessing trace from the argument to an input. It cannot be empty,
since an edge cannot target an input. Its first edge equals the given
operation edge by uniqueness of incoming incidence. The operation's singleton
source then identifies the next vertex in the trace with its source. The
remaining trace proves source reachability.
\end{proof}
\lean{Graph/SearchReachability}{21}{25}{Linear.OperationReachable.argument}
\lean{Graph/SearchReachability}{28}{39}{Linear.OperationReachable.source}

\begin{theorem}[Search preserves operation reachability]
Every coordinate search rule preserves operation reachability of its vertex.
Recursive renaming preserves and reflects this property.
\end{theorem}
\begin{proof}
A variable trigger keeps the vertex. Operation entry either keeps its source,
for no term arguments, or moves to the first argument, using forward argument
reachability. Return from an argument first gives reachability of the operation
source by the converse lemma. It then either stays there on completion or
moves to the next argument by the forward lemma.

For renaming, transport the reverse operation trace and its edge list through
the recursive graph isomorphism. The boundary input is fixed. Apply the same
construction to the inverse isomorphism and cancel the port maps for reflection.
\end{proof}
\lean{Graph/SearchReachability}{42}{56}{Linear.SearchStep.operationReachable}
\lean{Graph/SearchReachability}{59}{62}{Linear.OperationReachable.rename}
\lean{Graph/SearchReachability}{64}{71}{Linear.OperationReachable.rename_iff}

\Needspace{20\baselineskip}
\subsection{Incoming paths and insertion of the focus}
\begin{definition}[Incoming incidence trace]
\termdef{incoming-trace}
An incoming trace is a reverse path with explicit vertex incidences but no
label restriction. It is either empty at a vertex or consists of an edge
whose target contains the current vertex, a chosen source vertex of that edge,
and a remaining incoming trace from that source.
\end{definition}
\lean{Graph/RootedPaths}{11}{14}{Linear.IncomingTrace}

\begin{lemma}[Incoming paths to reachable vertices]
An edge targeting an operation-reachable vertex is an operation, and each of
its sources is operation-reachable. Consequently, every incoming trace ending
at a reachable vertex is an operation trace with the same vertices and edge list.
\end{lemma}
\begin{proof}
Compare the incoming edge with the first edge of a witnessing reverse
operation trace to an input. The trace cannot be empty because inputs have no
incoming edge. Target-incidence uniqueness identifies the two edges, proving
its operation label. An operation has exactly one source, so every chosen
source is the next vertex in the witnessing trace and is reachable by its tail.
Induct on the arbitrary incoming trace and apply this argument at each step.
\end{proof}
\lean{Graph/RootedPaths}{18}{31}{Linear.OperationReachable.predecessor}
\lean{Graph/RootedPaths}{34}{41}{Linear.IncomingTrace.operations}
\lean{Graph/RootedPathObstruction}{11}{19}{Linear.OperationReachable.incoming_operation}

\begin{lemma}[A path never departs from its final vertex earlier]
Suppose an operation trace from $v$ to an input has length $k$. In any operation
trace to an input with length at most $k$, no edge can have $v$ as a source.
In particular the original trace never has $v$ as a source of one of its edges.
\end{lemma}
\begin{proof}
Induct on the shorter trace's edge list. If its first edge has $v$ as source,
its singleton-source property identifies the next vertex with $v$. The tail
and the original length-$k$ trace therefore both run from $v$ to inputs.
Uniqueness of operation traces to inputs equates their edge lists and hence
their lengths. But the tail is shorter than the trace of length at most $k$,
a contradiction. For a later edge, apply the induction hypothesis to the tail,
whose length is still at most $k$.
\end{proof}
\lean{Graph/RootedPaths}{45}{68}{Linear.operationTrace_sources_avoid}

\begin{theorem}[Lifting the operation path through focus insertion]
An operation path from an input to $v$ survives insertion of a focus at $v$:
it becomes a path from the same input to the hosted copy of $v$, using the
retained copies of its original edges.
\end{theorem}
\begin{proof}
The previous lemma says no source occurrence on this path equals $v$.
Insertion therefore maps all its source occurrences by hosting rather than
to the fresh focus target. Targets are always hosted. Induct on the reverse
trace, retaining each operation label and mapping each source and target
membership into its incidence list. The empty trace and the boundary input
are fixed. This proves a path in the actual inserted graph, not just in its
focus erasure.
\end{proof}
\lean{Graph/RootedPaths}{72}{95}{Linear.operationTrace_insert}

\Needspace{22\baselineskip}
\subsection{Rooted states have actual operation paths: Lemma 23}
\label{sec:rooted-operation-path}
Define the intermediate invariant $U(s)$ to mean that the extracted coordinate
of $s$ is operation-reachable in its erased graph. It is invariant under
focussed recursive isomorphisms: erasure naturality transports the graph,
coordinate naturality identifies the vertex, and reachability is invariant
under this transport.
\lean{Graph/RootedStatePaths}{12}{13}{Linear.State.UnderlyingReachable}
\lean{Graph/RootedStatePaths}{15}{21}{Linear.State.underlyingReachable_iso}

\begin{lemma}[Realization and search preserve the intermediate invariant]
Realizing a reachable typed coordinate satisfies $U$. Every canonical root
satisfies $U$, and recursive state search preserves $U$.
\end{lemma}
\begin{proof}
Coordinate recovery provides one isomorphism that restores both the original
graph and the chosen coordinate after realizing and erasing. Reflect operation
reachability along this isomorphism to prove the realization assertion.
The initial coordinate is an input, hence reachable, which proves the root
assertion. For a search step, coordinate search preserves reachability; apply
the realization assertion to its successor coordinate and transport the
result through its supplied recursive endpoint isomorphism.
\end{proof}
\lean{Graph/RootedStatePaths}{24}{33}{Linear.State.underlyingReachable_realize}
\lean{Graph/RootedStatePaths}{35}{36}{Linear.State.root_underlyingReachable}
\lean{Graph/RootedStatePaths}{39}{42}{Linear.State.deepSearch_underlyingReachable}

The invariant descends to recursive classes using its isomorphism invariance.
Quotient elimination proves the root and one-step preservation laws there.
Induction on the witnessing root run proves that every rooted class satisfies
$U$. This is the search-sequence induction used for the path argument in
Appendix~C, expressed through an invariant of erased coordinates.
\lean{Graph/RootedStatePaths}{66}{68}{Linear.DeepStateClass.UnderlyingReachable}
\lean{Graph/RootedStatePaths}{70}{72}{Linear.DeepStateClass.root_underlyingReachable}
\lean{Graph/RootedStatePaths}{74}{77}{Linear.DeepStateClass.search_underlyingReachable}
\lean{Graph/RootedStatePaths}{80}{89}{Linear.DeepStateClass.rooted_underlyingReachable}

\begin{theorem}[Concrete rooted-focus operation path]
If a concrete state is rooted, its actual focus source is reachable from
its unique input by an operation path. If the focus source is not an input,
this path is nonempty.
\end{theorem}
\begin{proof}
First consider any state satisfying $U$. Lift its witnessing erased operation
trace through insertion at its extracted coordinate, by the insertion theorem.
The general recovery isomorphism carries this inserted graph to the actual
state and carries the inserted focus source to its actual focus source.
Transport the lifted operation trace through that isomorphism. Boundary
preservation fixes the input. This gives operation reachability in the actual
state graph.

A rooted state satisfies $U$ by the invariant proof, so the construction applies.
The interface has one input, identifying its index with zero. Reverse the trace
list to obtain the existing forward operation-path representation. If this
list were empty, the zero-step reverse trace would identify the focus source
with the input; hence the path is nonempty whenever the source is not an input.
\end{proof}
\lean{Graph/RootedStatePaths}{45}{59}{Linear.State.focus_operationReachable}
\lean{Graph/RootedStatePaths}{100}{102}{Linear.State.rooted_focus_operationReachable}
\lean{Graph/RootedStatePaths}{106}{121}{Linear.State.rooted_focus_path}
This establishes the graph-path conclusion of Lemma~23 for the present concrete
search model. Unlike the paper's presentation, the proof does not first require
a separate uniqueness theorem for whole search runs: it uses the operation-path
uniqueness already proved and a direct search invariant.

\Needspace{20\baselineskip}
\subsection{The non-operation-path obstruction: Lemma 26(1)}
\label{sec:rooted-path-obstruction}
\begin{theorem}[Incoming-path obstruction]
Every incoming incidence trace to a rooted focus source is an operation trace.
Thus an incoming trace that is not an operation trace proves that the state
is not rooted.
\end{theorem}
\begin{proof}
The preceding rooted-focus theorem gives reachability of the actual source.
Apply the incoming-trace theorem at that vertex. The obstruction is its
contrapositive.
\end{proof}
\lean{Graph/RootedStatePaths}{124}{128}{Linear.State.rooted_incoming_operations}
\lean{Graph/RootedStatePaths}{132}{135}{Linear.State.not_rooted_of_nonoperation_path}

\begin{lemma}[The existing edge-path representation]
For any nonempty edge path whose last edge targets a reachable vertex,
every edge of the path is an operation, and every source of each edge is
reachable. Repeated edges are allowed in the supplied path datatype.
\end{lemma}
\begin{proof}
Induct on the path. For a single edge, use the incoming-operation and source
reachability lemmas at its target. For a path beginning with adjacent edges
$e,f$, apply the induction hypothesis to the tail beginning at $f$. The
adjacency witness is a source of $f$, hence reachable, and a target of $e$.
The incoming-edge lemmas then prove the required claims for $e$, while the
induction hypothesis proves them for all remaining edges.
\end{proof}
\lean{Graph/RootedPathObstruction}{23}{45}{Linear.path_operations_of_reachable}

\begin{theorem}[Concrete Lemma 26(1)]
If a nonempty graph path ends at the focus source of a rooted state, all of
its edges are operations. Therefore a path ending there and containing any
non-operation edge is a certificate that the state is not rooted.
\end{theorem}
\begin{proof}
Apply the edge-path lemma to the actual focus-source reachability supplied
by rootedness. For the obstruction, select the non-operation edge from the
path and contradict the resulting operation-label assertion.
\end{proof}
\lean{Graph/RootedPathObstruction}{52}{57}{Linear.State.rooted_path_operations}
\lean{Graph/RootedPathObstruction}{61}{65}{Linear.State.not_rooted_of_path}
This is the first part of the paper's obstruction lemma in our concrete path
and search representations. The one-way replacement argument in Lemma~26(2)
is proved separately in Section~\ref{sec:one-way-obstruction}. The general
output-closed specimen argument in Lemma~26(3) is proved for shallow
specimens in Section~\ref{sec:shallow-specimens}; its deep extension remains open.

\Needspace{20\baselineskip}
\subsection{Concrete refocusing and the remaining rewrite obligation}
\begin{definition}[Concrete refocusing]
\termdef{concrete-refocusing}
A transition relation on recursive state classes is refocusing when each
transition from a concretely rooted state leads to a concretely rooted state.
The rooted predicate is the one constructed from actual search above; it is
not the abstract specimen interface's supplied rooted predicate.
\end{definition}
\lean{Graph/ConcreteRefocusing}{21}{23}{Linear.DeepStateClass.Refocusing}

\begin{theorem}[Search refocusing and reduction to rewriting]
Search is refocusing. For any supplied rewrite relation $R$, the relation
$\mathsf{Search}\cup R$ is refocusing if and only if $R$ preserves rootedness.
For any machine on these concrete classes whose step relation is refocusing,
every execution from an initial class ends in a rooted class.
\end{theorem}
\begin{proof}
Search refocusing is the proved one-step rootedness invariant. If the union
is refocusing, apply its law to transitions from its $R$ branch. Conversely,
case-split a union transition: use search preservation in the search branch
and the supplied preservation law in the rewrite branch. Finally induct on
a machine execution, starting with the theorem that initial classes are rooted
and applying its refocusing hypothesis at each step.
\end{proof}
\lean{Graph/ConcreteRefocusing}{26}{28}{Linear.DeepStateClass.search_refocusing}
\lean{Graph/ConcreteRefocusing}{32}{40}{Linear.DeepStateClass.refocusing_with_search}
\lean{Graph/ConcreteRefocusing}{44}{53}{Linear.DeepStateClass.refocusing_execution}
The last two statements are conditional on the indicated rewrite or machine
preservation hypothesis. They do not construct substitution or compute
transitions, and do not assert that arbitrary rewrite relations preserve roots.

\Needspace{18\baselineskip}
\subsection{Rooted and non-rooted concrete examples}
The initial value state is rooted for any operation classification, by initiality.
Its successful passive search execution gives a rooted final value state by
preservation along search runs.
\lean{Graph/RootedExamples}{9}{10}{Linear.Examples.initial_value_rooted}
\lean{Graph/RootedExamples}{13}{14}{Linear.Examples.final_value_rooted}

The previously constructed interior-focus state is valid and hole-free but
is not rooted, for any operation classification. Its underlying closed chain
has a variable before the focus position and a value after it. In the actual
inserted graph, the retained variable edge targets the focus source. The
single-edge path consisting of that variable therefore ends at the source,
and its label is not an operation. The concrete Lemma~26(1) obstruction rules
out rootedness. This separates graph validity from the operational invariant
and exercises the obstruction on an actual named state.
\lean{Graph/RootedExamples}{18}{26}{Linear.Examples.interior_state_not_rooted}

\Needspace{18\baselineskip}
\subsection{One-way hypernets and the term-boundary restriction}
\label{sec:one-way}
This milestone proves the single-hole case of Appendix~C, Lemma~26(2), for
actual shallow substitution and the concrete rooted-state predicate. The
proof follows the paper's reason for introducing one-way hypernets: an
operation path from the global input cannot pass through such a filling and
emerge at a term output. We make the path-extraction argument explicit and
allow the replaced edge to have a self-loop.

\begin{lemma}[Sources of operation paths have term type]
Every source vertex of an actual operation edge has type $\star$. If a
nonempty reverse operation trace ends at an input, that input has type
$\star$.
\end{lemma}
\begin{proof}
The signature of every primitive operation has exactly one source, of type
$\star$. Apply the source-typing equation to any member of its source list.
For the second assertion, induct on the trace's edge list. If its tail is
empty, its endpoint is the source of the first operation and the first
assertion applies. Otherwise apply the induction hypothesis to the tail.
Boxes, variables, focuses, and holes are excluded by the operation predicate.
\end{proof}
\lean{Graph/OneWay}{10}{22}{Linear.operation_source_term}
\lean{Graph/OneWay}{26}{36}{Linear.operationTrace_input_term}

\begin{definition}[One-way hypernet]\label{term:one-way}
A hypernet $H:I\to O$ is \emph{one-way} if, for every term input $i$ and term
output $j$, there is no operation path from $i$ to $j$. Paths to thunk outputs
are deliberately not excluded. This is the restriction in the paper's
Section~6.1, expressed using our existing operation-path datatype.
\end{definition}
\lean{Graph/OneWay}{40}{43}{Linear.OneWay}

\begin{theorem}[Reachability characterization]
A hypernet is one-way if and only if none of its term outputs is
operation-reachable from an input.
\end{theorem}
\begin{proof}
A reverse trace from an output to an input cannot be empty: the boundary
constructors are disjoint. Its input therefore has term type by the preceding
lemma. Reverse its recorded edge list to obtain the forbidden forward
operation path. Conversely any forbidden forward path supplies precisely a
reverse trace witnessing operation reachability of its term output.
\end{proof}
\lean{Graph/OneWay}{46}{54}{Linear.oneWay_iff}

One-wayness is invariant under recursive graph isomorphism. Such an
isomorphism fixes the ordered boundary and preserves and reflects operation
reachability, so the characterization applies in both directions. A graph
with no outputs is one-way vacuously, since its output index type is empty.
\lean{Graph/OneWay}{66}{73}{Linear.DeepIso.oneWay_iff}
\lean{Graph/OneWay}{76}{78}{Linear.oneWay_no_outputs}

\Needspace{18\baselineskip}
\subsection{Extracting an operation path from a filling}
Fix a host $C$, a selected edge $h$, and a filling $H$ with the interface of
$h$. Write $P=C[H]$, $\iota_C$ for the retained-host vertex map, and
$\iota_H$ for the filling map. These are the existing maps of shallow
replacement; $\iota_H$ need not be injective when $h$ has a self-loop.

The elementary decomposition used in the induction says that a nonempty
reverse operation trace with edge list $e::E$, from $v$ to $z$, consists
exactly of an intermediate vertex $w$, an operation label on $e$, incidences
$v\in\operatorname{target}(e)$ and $w\in\operatorname{source}(e)$, and a
trace with edge list $E$ from $w$ to $z$. Both directions follow directly
from the trace constructor.
\lean{Graph/OneWay}{57}{63}{Linear.operationTrace_cons_iff}

\begin{lemma}[An inserted non-input stays a non-input]
If a vertex $v$ of $H$ is not an input, its image $\iota_H(v)$ is not an
input of $P$.
\end{lemma}
\begin{proof}
Target coverage in $H$ supplies an edge targeting $v$. Its inserted copy
targets $\iota_H(v)$ in $P$. Inputs of a well-formed graph have no incoming
edge, which proves the claim, including when $v$ is a filling output.
\end{proof}
\lean{Graph/OneWay}{88}{92}{Linear.OneWayReplacement.insert_not_input}

\begin{theorem}[Extraction of a filling-local input path]
If $\iota_H(v)$ is operation-reachable in $P$, then $v$ is
operation-reachable inside $H$.
\end{theorem}
\begin{proof}
Choose a reverse operation trace from $\iota_H(v)$ to a global input and
induct on its length, generalizing $v$. If $v$ is already an input of $H$,
the empty local trace is the required witness. Otherwise, in the zero-length
case, $\iota_H(v)$ would be a global input, contradicting the previous lemma.

In the nonempty case, target coverage supplies an edge $f$ of $H$ targeting
$v$. Its inserted copy and the first edge of the global trace target the
same vertex. Unique incoming incidence identifies these two edges. The
first edge is consequently a filling edge and retains its operation label.
Its chosen source is the image of some source $w$ of $f$. Apply the induction
hypothesis to the shorter global trace from that image; this yields a local
operation trace from $w$ to an input of $H$. Prepending $f$ yields the desired
local trace from $v$.

The induction stops on first reaching a filling input. It neither reflects
an arbitrary whole edge path through $\iota_H$ nor assumes injectivity of
that map. This is why the proof remains valid for self-looping holes, unlike
unrestricted path reflection discussed earlier.
\end{proof}
\lean{Graph/OneWay}{98}{124}{Linear.OneWayReplacement.extract}

\Needspace{18\baselineskip}
\subsection{Reflecting host reachability through a one-way filling}
\begin{theorem}[Host term reachability reflects]
Suppose $H$ is one-way and $v$ is a term vertex of $C$. If $\iota_C(v)$ is
operation-reachable in $C[H]$, then $v$ is operation-reachable in $C$.
\end{theorem}
\begin{proof}
Induct on the edge list of a reverse operation trace from $\iota_C(v)$ to
a global input. In the empty case injectivity of the retained-host map
identifies $v$ with that same input.

For a nonempty trace, inspect its first edge. If it is a surviving host
edge, reflect its target incidence by injectivity of $\iota_C$ and express
its source as $\iota_C(w)$. The corresponding host edge is an operation,
so $w$ has term type. Apply the induction hypothesis to the shorter trace
and prepend the host edge.

If the first edge is inserted from $H$, its target is some
$\iota_H(x)=\iota_C(v)$. A filling target cannot be an input. The boundary
classification of replacement therefore forces $x$ to be an output of $H$.
Both replacement maps preserve vertex labels, so this output has term type.
Extraction applied to the original trace now makes it operation-reachable
inside $H$, contradicting one-wayness. Thus only the retained-host case
can occur.
\end{proof}
\lean{Graph/OneWay}{128}{158}{Linear.OneWayReplacement.reflect}

\begin{theorem}[Incoming host paths consist of operations]
Under the same hypotheses, every edge of any host path ending at $v$ is an
operation. The given host path is allowed to contain the replaced edge.
\end{theorem}
\begin{proof}
Reflect operation reachability to $C$, then apply the earlier theorem that
all edges of any path ending at an operation-reachable vertex are operations.
That theorem uses the existing nonempty edge-path representation and allows
repeated edges. No assertion that the supplied host path itself survives
filling is needed.
\end{proof}
\lean{Graph/OneWay}{162}{168}{Linear.OneWayReplacement.path_operations}

\begin{theorem}[One-way filling closure]
If both $C$ and $H$ are one-way, then $C[H]$ is one-way.
\end{theorem}
\begin{proof}
Suppose a term output of $C[H]$ were operation-reachable. It is the retained
copy of the corresponding term output of $C$. Reflection would make that
output operation-reachable in $C$, contradicting the host's one-wayness.
\end{proof}
\lean{Graph/OneWay}{171}{176}{Linear.OneWayReplacement.oneWay_plug}

\Needspace{18\baselineskip}
\subsection{The concrete one-way obstruction: Lemma 26(2)}
\label{sec:one-way-obstruction}
For a focussed host $C:\star\to\epsilon$, replace an edge other than its
focus by a globally focus-free filling $H$. Once the resulting graph is
proved globally hole-free, package it with its retained focus as an actual
state. The constructor requires these focus-freedom and hole-freedom proofs;
it does not infer them merely from the interfaces. Its focus source is
$\iota_C(s)$, where $s$ is the host's focus source.
\lean{Graph/OneWayObstruction}{12}{16}{Linear.Focussed.filledState}

\begin{theorem}[Concrete Lemma 26(2)]
Suppose the focussed host has a path ending at its focus source and that
path contains a non-operation edge. If the filling $H$ is one-way, the
resulting actual state is not rooted.
\end{theorem}
\begin{proof}
Assume the filled state is rooted. The concrete Lemma~23 gives operation
reachability of its focus source $\iota_C(s)$ in the actual filled graph.
The source $s$ of the host focus has term type. Apply the preceding host-path
theorem: every edge of the supplied host path must be an operation. This
contradicts its specified non-operation edge.
\end{proof}
\lean{Graph/OneWayObstruction}{21}{32}{Linear.Focussed.oneWay_not_rooted}

This proves the single-hole one-way obstruction for the current concrete
search semantics. It follows the uniqueness-and-extraction argument of
Proof~C.8, packaged as a reflection lemma rather than a split on whether the
supplied path contains the hole. In fact the Lean result permits replacement
at any non-focus edge with the right interface, so it includes the paper's
hole-labelled context case. The example below uses a genuine hole label.

\begin{theorem}[Paired single-hole obstruction]
If $H$ and $K$ are two admissible fillings of the same focussed host and at
least one is one-way, then at least one of the two filled states is not rooted,
provided the host has the same non-operation path to its focus source.
\end{theorem}
\begin{proof}
Split the hypothesis that $H$ or $K$ is one-way. Apply the preceding theorem
to that filling and introduce the corresponding side of the conclusion.
\end{proof}
\lean{Graph/OneWayObstruction}{37}{50}{Linear.Focussed.oneWay_pair_not_rooted}
This is a single-hole instance of the output-closure argument.
Section~\ref{sec:shallow-specimens} defines concrete pre-templates,
output-closure, and shallow multi-hole specimens and proves their
Lemma~26(3) obstruction. Deep specimen substitution remains open.

\Needspace{18\baselineskip}
\subsection{One-way examples and a filled-state obstruction}
The variable corolla is one-way. Its term output is targeted by a variable
edge. If that output were operation-reachable, the incoming-edge theorem
would force this variable edge to be an operation, a contradiction. By
contrast, the existing operation corolla with two term arguments is not
one-way: its sole edge is a one-step operation path from the input to the
first term output.
\lean{Graph/OneWayExamples}{9}{17}{Linear.Examples.variable_oneWay}
\lean{Graph/OneWayExamples}{21}{24}{Linear.Examples.operation_not_oneWay}

Define an operation signature with no term arguments and one thunk argument,
and take its one-edge corolla. This graph has a term input and only a thunk
output. It is one-way because no output has term type, yet its sole operation
edge explicitly witnesses reachability of its thunk output. Together these
examples verify both the nontriviality of the predicate and the importance
of retaining the paper's term-output restriction.
\lean{Graph/OneWayExamples}{27}{27}{Linear.Examples.thunkOnlyOperations}
\lean{Graph/OneWayExamples}{29}{31}{Linear.Examples.thunkOnlyGraph}
\lean{Graph/OneWayExamples}{34}{39}{Linear.Examples.thunk_only_oneWay}
\lean{Graph/OneWayExamples}{42}{43}{Linear.Examples.thunk_output_reachable}

For a nonvacuous state-level application, start with the earlier closed chain
whose interior focus follows a variable and precedes a value. Relabel that
variable by a hole with one term input and one term output, preserving all
incidences. The other labels remain the focus and the value operation. Their
signatures are unchanged, so the relabelled graph is well formed. Retain the
same distinguished focus, mode, and endpoints. The only other edges are a
hole and an operation, both focus-free; the focus therefore remains globally
unique and satisfies its original endpoint conditions.
\lean{Graph/OneWayExamples}{47}{58}{Linear.Examples.oneWayContext}
\lean{Graph/OneWayExamples}{61}{77}{Linear.Examples.oneWayContext_focus}

Fill the genuine hole by the variable corolla. The exact hole inventory of
shallow substitution shows the resulting graph is hole-free: the filling is
hole-free and no hole survives in the host. Indeed the only host edge that
could carry a hole is exactly the removed edge. Global focus freedom of the
variable filling was already checked.
\lean{Graph/OneWayExamples}{80}{97}{Linear.Examples.oneWayContext_filled_closed}

Finally, take the host path consisting only of its hole edge. Its target is
the focus source, and its label is not an operation. The variable filling is
one-way, so the concrete Lemma~26(2) theorem proves that this actual filled
state is not rooted, for every operation classification. The offending host
edge is removed by substitution: this exercises the new extraction argument,
not merely preservation of a path of retained host edges.
\lean{Graph/OneWayExamples}{101}{109}{Linear.Examples.oneWay_filled_not_rooted}

\Needspace{32\baselineskip}
\section{Proof checking and remaining obligations}
\Needspace{24\baselineskip}
\subsection{The two axiom audits}
The file \leanfile{Hypernets/Audit.lean}{Audit.lean} contains six
\texttt{\#print axioms} commands. They inspect counting-simulation soundness,
the forward and converse simulation constructions, the two contextual sufficiency
theorems, and the zero-step converse regression theorem. These diagnostic commands
introduce no new definitions or proofs.
\lean{Audit}{5}{10}{Audit: abstract theorem dependencies}
The graph audit inspects shallow plugging, host and inserted path equivalence,
plugging congruence, exact trace uniqueness, typed operation-path uniqueness,
finite deep enumeration, unique holes under boxing, deep plugging and its other-label
locality, and the self-loop and deep-replacement regression proofs. Twelve further commands
audit occurrence-coordinate observation preservation, exact hole inventories,
sharp freshness and consumption, isomorphism invariance of contexts, deep site
provenance, deep context closure, and the new reuse and collision regressions.
\lean{Graph/Audit}{5}{17}{Graph.Audit: concrete theorem dependencies}
\lean{Graph/Audit}{19}{30}{Graph.Audit: context calculus dependencies}
Nine further commands audit the exact deep inventory, both versions of sharp
context closure, hole-freeness, selected-label consumption, certified deep
plugging, and the deep reuse and collision regressions.
\lean{Graph/Audit}{32}{40}{Graph.Audit: exact deep inventory dependencies}
Eleven further commands audit recursive label transitivity, operation-path
invariance, occurrence embeddings, hole and context invariance, deep congruence,
quotient plugging and context validity, and the strict-extension regressions.
\lean{Graph/Audit}{42}{52}{Graph.Audit: recursive isomorphism dependencies}
Ten further commands audit shallow host congruence, site transport and its
specification, simultaneous congruence, quotient invariance, retained holes,
freshness, result validity, and the moved-address and wrong-address regressions.
\lean{Graph/Audit}{54}{63}{Graph.Audit: site transport dependencies}
Sixteen further commands inspect shallow associativity, both corolla units,
deep associativity, quotient composition, composed retained inventories, sharp
composite context closure and its certified constructor, three-stage coherence,
deep restoration, quotient identity, and the self-loop, repair, and mixed-depth
regressions.
\lean{Graph/Audit}{65}{80}{Graph.Audit: composition and unit dependencies}

Six further commands inspect global focus-free substitution, focussed shallow
and deep substitution, finite-filling boundary agreement, the initiality
regression, and rejection of a boxed focus.
\lean{Graph/Audit}{82}{87}{Graph.Audit: concrete focus dependencies}

Eight further commands inspect operation-view existence, search determinism,
search boundary laws, shallow filling preservation, and the thunk, value, and
variable regressions.
\lean{Graph/Audit}{89}{96}{Graph.Audit: search rule dependencies}

Ten further commands inspect explicit focus insertion and certification, both
incidence-recovery laws, hole freedom, term-typing preservation, initial
attachment, search endpoint realization, and the interior and value regressions.
\lean{Graph/Audit}{98}{107}{Graph.Audit: focus realization dependencies}

Twelve further commands inspect general erasure and its invariants, the
reinsertion isomorphism and focus endpoints, arbitrary-state recovery, state
search determinism and boundary laws, and the execution/interior regressions.
\lean{Graph/Audit}{109}{120}{Graph.Audit: erasure and state search dependencies}

Fifteen further commands inspect insertion--erasure cancellation, coordinate
transport and realization, erasure naturality, source invariance, quotient
search determinism and boundaries, and the concrete quotient regressions.
\lean{Graph/Audit}{122}{136}{Graph.Audit: round trips and quotient search dependencies}

Fourteen further commands inspect recursive focus freedom, coordinate search,
insertion and erasure, recursive state search and quotient determinism, the
quotient projection and lifting laws, exact-length executions, and regressions.
\lean{Graph/Audit}{138}{151}{Graph.Audit: recursive state search and execution dependencies}

Sixteen further commands inspect canonical roots, search preservation of
rootedness, operation reachability, insertion of rooted paths, the concrete
path obstruction, refocusing, and the positive and negative state examples.
\lean{Graph/Audit}{153}{168}{Graph.Audit: rooted-state and path-obstruction dependencies}

In the pinned snapshot, counting-simulation soundness, finite deep enumeration,
unique holes under boxing, occurrence-coordinate observation preservation,
isomorphism invariance of unique holes, and uniqueness of the two-hole example
report only \texttt{propext} and \texttt{Quot.sound}.
Recursive label transitivity, recursive operation-path invariance, the
nonexistence of a literal-label isomorphism between the example boxes, and their
equality as quotient classes likewise report only these two axioms. Existence
of transported sites also reports only these two axioms; the chosen transport
and substitution results additionally use classical choice.
All audited composition and unit results report these three standard
axioms. The focus results do too, except the boxed-focus rejection, which uses
only \texttt{propext} and \texttt{Quot.sound}. The search audits use only \texttt{propext} and \texttt{Quot.sound}, except
shallow filling preservation, which also uses classical choice. The new realization audits use all three standard axioms, except term-typing
preservation, which uses only propositional extensionality and quotient soundness.
All new general-erasure and state-search audit entries use those three
standard axioms. The round-trip and quotient-search audits do too, except
coordinate search renaming, which uses only \texttt{propext} and
\texttt{Quot.sound}. The remaining audited declarations additionally report \texttt{Classical.choice}.
The recursive-search audits use all three axioms, except label and graph
focus-freedom invariance, coordinate search transport, and the renamed-box
focus-freedom example, which use only \texttt{propext} and \texttt{Quot.sound}.
The rooted-state audit entries use all three axioms, except coordinate
operation-reachability preservation and the incoming-trace operation theorem,
which use only \texttt{propext} and \texttt{Quot.sound}.
These are standard Lean logical axioms: propositional extensionality, quotient
soundness, and classical choice, respectively. Neither audit reports
\texttt{sorryAx} or a project-specific axiom.

Absence of project-specific axioms does not discharge a theorem's hypotheses.
Specimen laws and initial lifting remain parameters of the abstract sufficiency
results. They do not appear as named global axioms in the report. In contrast,
the concrete graph constructions explicitly prove the invariant fields needed
to construct their results. Neither fact supplies the missing bridge to a
concrete UAM.

\Needspace{24\baselineskip}
\subsection{Reproducing the Lean and document checks}
The snapshot pins Lean \texttt{4.35.0} in
\leanfile{lean-toolchain}{lean-toolchain} and uses only the bundled \texttt{Std}.
With that toolchain available, run
\begin{quote}\ttfamily
lake build\\
lake env lean Hypernets/Audit.lean\\
lake env lean Hypernets/Graph/Audit.lean
\end{quote}
The script \leanfile{scripts/check.sh}{scripts/check.sh} runs the build and both
audits. Run it using \texttt{bash} with the script path as its argument. It selects
the workspace-local
macOS ARM compiler if present, and otherwise uses \texttt{lake} from the environment.
The audits are checked separately from the root import.

The commentary's source checker recursively scans all Lean modules under
\texttt{Hypernets/}, including nested graph modules, attributed declarations,
noncomputable definitions, abbreviations, and anonymous examples. It verifies
that each declaration has a source-fragment link, all linked ranges are valid,
and each working source file equals the version at the pinned commit. Diagnostic
audit ranges are linked separately. Run it with
\begin{quote}\ttfamily\small
python3 docs/mathematical-commentary/check\_sources.py
\end{quote}
This syntactic coverage check complements the mathematical exposition; it does
not itself check the correctness of English explanations. Source links are
immutable GitHub commit links, so future edits cannot move the referenced proof.

From the commentary project directory, \texttt{make} checks the links, compiles
with \texttt{latexmk}, and exports the PDF to the repository's
\texttt{output/pdf/} directory. The committed PDF can therefore be read without
a TeX installation, and its editable sources and reproduction instructions are
included alongside it.

\Needspace{24\baselineskip}
\subsection{What remains for the paper's concrete theorem}
This is a to-do list only: completed obligations are removed. The remaining
work needed to claim Theorem~6 with its original hypotheses is:
\begin{enumerate}
\item Prove representation adequacy between the typed depth-indexed hierarchy
with explicit binder metadata and the paper's set-based graph-labelled
hypernets, including compatibility across different depth bounds.
\item Construct a quotient of host/site pairs and prove interchange of
substitutions at independent sites.
\item Identify erasure-based state search with the paper's contextual local
graph rewrites in both directions.
\item Define the concrete substitution and compute transitions, with their
typing and context restrictions, and assemble the full UAM transition relation.
\item Prove rootedness preservation for substitution and the required
full-UAM structural consequences of determinism and refocusing. Establish the
one-way and output-closed specimen cases of Lemma~26(2--3), and the stationary
rewrite criterion where needed.
\item Define concrete pre-templates, specimens, quasi-specimens, output-closure,
input-safety, and robustness with the paper's restrictions.
\item Construct the abstract specimen interface from these objects. Connect
concrete specimen substitution to the finite-filling boundary theorem for
Lemma~39(1), and prove the rootedness and context-only traversal laws needed
for Lemma~39(2--3), including entering/exiting classifications. Prove the safety
and robustness bridges and the concrete observation/context membership
witnesses needed for initial lifting.
\item Instantiate the checked contextual sufficiency results using the
concrete interface, observations, and lifting. Discharge every remaining
concrete obligation to obtain the target robustness theorem.
\end{enumerate}

\Needspace{14\baselineskip}
\begin{thebibliography}{9}
\bibitem{paper}
Dan R. Ghica, Koko Muroya, and Todd Waugh Ambridge.
\emph{A robust graph-based approach to observational equivalence}.
Logical Methods in Computer Science 21(2:8), 2025.
Reference version: \href{https://arxiv.org/abs/1907.01257v5}{arXiv:1907.01257v5}.
\bibitem{code}
\emph{Hypernets in Lean 4}.
\href{https://github.com/danghica/hypernets-lean/tree/fe1b8904b1e394450111eabffaf4a7eded297740}{Repository at revision \texttt{fe1b890}}.
\end{thebibliography}
\clearpage
\appendix
\section{Technical terms}
\label{app:terms}

This appendix is alphabetical. Click a term or its page number to jump to its
definition in this PDF. Related terms may share a definition. Descriptions of
abstract interface predicates state what the interface supplies; they do not
claim that a concrete UAM interpretation has already been constructed.

\glossaryentry{path}{Adjacency}{Edges $e,f$ are adjacent when a vertex occurs in both the target list of $e$ and the source list of $f$.}

\glossaryentry{observations}{Allowed context}{A context satisfying the supplied predicate $\mathsf{Allowed}$ in an observation interface.}

\glossaryentry{box}{Alphabet}{A type of edge labels equipped with a signature assigning their input and output wire types.}

\glossaryentry{focus}{Backtracking mode}{The abstract focus mode $\mathsf{backtrack}$, one of the three constructors of the focus datatype.}

\glossaryentry{bijection}{Bijection}{A pair of forward and inverse functions satisfying both inverse equations.}

\glossaryentry{ports}{Boundary port}{An input or output constructor carrying an index in the corresponding ordered boundary list; input and output constructors are disjoint.}

\glossaryentry{box}{Box}{A label storing bound/free counts $n,m$ and a body of type $*\to *^{n+m}$; its outer edge has type $T^n(*)\to *^m$.}

\glossaryentry{box-replacement}{Box-label replacement}{Changing the body of one selected box edge while retaining its binding metadata and all outer incidence data.}

\glossaryentry{fibre}{Child occurrence}{An occurrence inside a box body. Primitive labels have no child occurrences.}

\glossaryentry{compatible}{Compatible provenance}{The condition that no hole label in the replacement occurs at a retained host origin of the provenance map.}

\glossaryentry{net}{Concrete typed net}{A finite labelled hypergraph with ordered boundaries and certified distinctness, incidence uniqueness, coverage, and signature typing.}

\glossaryentry{context}{Context}{A hypernet in which each complete typed hole label occurs at most once over all nesting depths.}

\glossaryentry{specimen}{Context-only traversal}{The specimen-interface law supplying paired context steps for a rooted, settled specimen whenever either endpoint can step.}

\glossaryentry{contextual}{Contextual equivalence}{Mutual contextual refinement: both directions hold for every pair of observations in the same allowed context.}

\glossaryentry{contextual}{Contextual refinement}{The requirement that observations of two fragments in each common allowed context produce states related by the chosen state refinement.}

\glossaryentry{corolla}{Corolla}{A graph with one operation edge, no internal vertices, and precisely that operation's ordered input and output boundaries.}

\glossaryentry{cost-orders}{Cost converse}{The relation $Q^{-1}$ defined by $a\,Q^{-1}\,b$ exactly when $b\,Q\,a$.}

\glossaryentry{cost-preorder}{Cost preorder}{A reflexive and transitive relation on natural-number execution counts; additional algebraic properties are imposed separately.}

\glossaryentry{execution}{Counted execution}{The relation $s\steps{n}t$, recording exactly $n$ machine transitions, including the zero-step execution from a state to itself.}

\glossaryentry{simulation}{Counting simulation}{A relation preserving finality and matching steps using related execution counts and an up-to relation, with an explicit stuck-execution alternative.}

\glossaryentry{all-edges}{Deep occurrence}{An address that records an enclosing box edge and recursively an occurrence in its body, distinguishing equal bodies at different locations.}

\glossaryentry{deep-plugging}{Deep plugging}{Typed replacement at an addressed site, recursively descending through boxes and rebuilding the selected enclosing labels.}

\glossaryentry{hypernet}{Depth bound}{The index bounding permitted box nesting in a hypernet; a graph need not use all of the available depth.}

\glossaryentry{promotion}{Depth promotion}{Raising a hypernet's available depth bound by relabelling primitives and recursively promoting box bodies, without changing outer incidence.}

\glossaryentry{determinism}{Determinism}{The property that any two one-step successors of the same state are equal.}

\glossaryentry{hole-presence}{Disjoint hole inventories}{The condition that no complete typed hole label occurs in both graphs, including occurrences inside boxes.}

\glossaryentry{net}{Edge}{An element of a net's finite edge carrier, carrying an operation label and ordered source and target vertex lists.}

\glossaryentry{empty-graph}{Empty graph}{The net with no internal vertices, no edges, and empty input and output boundaries.}

\glossaryentry{specimen}{Entering specimen}{A specimen satisfying the interface's supplied entering predicate. Its concrete interpretation by hole incidence is not yet instantiated.}

\glossaryentry{cost-orders}{Exact cost order}{The preorder $E$ with $a\,E\,b$ exactly when $a=b$, requiring identical successful execution counts in refinement.}

\glossaryentry{fibre-provenance}{Fibre provenance}{An injective, hole-observation-preserving map from a changed label's occurrence fibre to retained old fibre addresses or replacement occurrences.}

\glossaryentry{machine}{Final state}{A state satisfying a machine's successful-finality predicate; the machine laws require that it has no outgoing transition.}

\glossaryentry{enumeration}{Finite enumeration}{A finite list containing every element of a carrier, together with a completeness proof; the list may contain repetitions.}

\glossaryentry{focus}{Focus mode}{One of $\mathsf{search}$, $\mathsf{backtrack}$, or $\mathsf{rewrite}$; this datatype alone does not specify a graph position.}

\glossaryentry{observations}{Fragment}{An element of the family $F_i$ in an observation interface, matched with contexts of the same interface index $i$.}

\glossaryentry{ports}{Graph signature}{The assignment of an ordered input-label list and output-label list to each operation label.}

\glossaryentry{hole-presence}{Hole inventory}{The collection of complete hole labels $h$ for which some shallow or deep occurrence observes $h$.}

\glossaryentry{wire}{Hole label}{A natural-number name together with input and output wire-type lists. Equality compares the complete typed label, not just its name.}

\glossaryentry{hole-lookup}{Hole lookup}{The recursive function returning the typed hole label at an occurrence, or no label when that occurrence is not a hole.}

\glossaryentry{hole-presence}{Hole presence}{The predicate $\mathsf{HasHole}(G,h)$, witnessed by an occurrence whose hole lookup returns $h$.}

\glossaryentry{hole-presence}{Hole-free graph}{A graph having no hole label at any occurrence or nesting depth.}

\glossaryentry{embeddings}{Host embedding}{The map retaining host vertices during shallow replacement, tagging internal names and fixing each external boundary position.}

\glossaryentry{hypernet}{Hypernet}{A typed net over a finite-depth alphabet whose labels may include boxes with bodies over the preceding alphabet.}

\glossaryentry{cost-orders}{Improvement cost order}{The preorder $I$ with $a\,I\,b$ exactly when $b\leq a$; the target successful execution takes no more steps than the source.}

\glossaryentry{net}{Incidence coverage}{The net laws saying that exactly non-output vertices occur in some source list and exactly non-input vertices occur in some target list.}

\glossaryentry{net}{Incidence uniqueness}{The net laws allowing a vertex to occur in the sources of at most one edge and in the targets of at most one edge.}

\glossaryentry{initial-lifting}{Initial lifting}{The contextual-lifting hypothesis placing initial observations of related fragments into the specified lifted state relation.}

\glossaryentry{safety}{Input-safety}{The abstract condition on rooted, entering search specimens: a stuck alternative, a cost-related execution to a settled specimen, or a quasi-specimen alternative.}

\glossaryentry{embeddings}{Insertion map}{The map tagging replacement internal vertices freshly and identifying its boundary positions with the selected host edge's indexed ports.}

\glossaryentry{net}{Interface}{For a concrete net, the ordered input and output wire-label lists; matching lists enforce type-correct replacement.}

\glossaryentry{ports}{Internal vertex}{A vertex from the internal carrier, tagged separately from all input and output boundary positions.}

\glossaryentry{isomorphism}{Isomorphism}{Bijections of internal vertices and edges preserving labels and ordered incidence lists while fixing the external boundary positions.}

\glossaryentry{lifting}{Lifted state relation}{The relation witnessed by a remote specimen with rooted endpoints equal to the two states being compared.}

\glossaryentry{machine}{Machine}{A transition relation on states, a successful-finality predicate, and a proof that final states cannot step.}

\glossaryentry{non-self-loop}{Non-self-loop hole}{A selected edge whose source and target lists are disjoint; this condition permits reflection of inserted paths.}

\glossaryentry{determinism}{Normal state}{A state with no outgoing transition, whether or not it is successfully final.}

\glossaryentry{observations}{Observation interface}{Typed fragment and context families, an allowed-context predicate, and a possibly partial, nonfunctional relation producing observed states.}

\glossaryentry{all-edges}{Occurrence}{An individual shallow edge or an addressed edge inside boxes; occurrences distinguish positions even when their labels agree.}

\glossaryentry{occurrence-coordinates}{Occurrence coordinates}{A dependent pair consisting of an outer edge and an address in the occurrence fibre of that edge's label.}

\glossaryentry{fibre}{Occurrence fibre}{The disjoint sum of the edge's own shallow occurrence and all child occurrences determined by its label.}

\glossaryentry{fibre-lookup}{Occurrence observation}{For a fibre address, the hole lookup obtained either from its shallow label or recursively from its child address.}

\glossaryentry{operation-path}{Operation path}{A forward, endpoint-sensitive path defined by a reverse operation trace with its edge list reversed; empty paths are permitted.}

\glossaryentry{wire}{Operation specification}{The assignment of a term-argument count and an ordered list of thunk binder counts to each operation.}

\glossaryentry{unary}{Operation step}{A step through an edge selected by the operation predicate, from a vertex in its source list to one in its target list.}

\glossaryentry{operation-trace}{Operation trace}{A reverse traversal recording each actual operation edge, starting at a target vertex and continuing from that edge's source.}

\glossaryentry{opposite}{Opposite specimen system}{The specimen system obtained by swapping its left and right endpoint maps while retaining modes and entering and rootedness predicates.}

\glossaryentry{cost-orders}{Overhead cost order}{The preorder $H$ with $a\,H\,b$ exactly when $a\leq b$; the target successful execution takes at least as many steps as the source.}

\glossaryentry{path}{Path}{A nonempty edge list whose consecutive edges are adjacent. Repeated edges are allowed; this predicate does not specify endpoint vertices.}

\glossaryentry{site}{Plugging site}{An address selecting an edge, possibly inside boxes, indexed by the required local depth bound and replacement interface.}

\glossaryentry{ports}{Port}{A vertex represented as an internal name, an indexed input, or an indexed output, with the three constructors kept disjoint.}

\glossaryentry{primitive}{Primitive label}{An operation, variable, focus-mode label, or typed hole; a primitive label has no box body.}

\glossaryentry{primitive-signature}{Primitive signature}{The typed port assignment for primitives: one term input for non-holes, operation outputs determined by arities, and stored interfaces for holes.}

\glossaryentry{provenance}{Provenance certificate}{An injective map assigning every result occurrence a retained host or inserted replacement origin with the same hole observation.}

\glossaryentry{quasi}{Quasi-specimen}{The abstract relation requiring a non-rewrite specimen witness, with prescribed endpoint relations, whenever the compared states are rooted.}

\glossaryentry{reasonable}{Reasonable triple}{Three cost preorders satisfying the addition, decrease, inclusion, and sandwich conditions used by counting-simulation soundness.}

\glossaryentry{reindexing}{Reindexing}{Renaming a net's internal and edge carriers along bijections, transporting labels and incidence while preserving boundary order.}

\glossaryentry{relabelling}{Relabelling}{Changing operation labels along a signature-preserving map, leaving the graph's carriers, vertex labels, and incidence unchanged.}

\glossaryentry{lifting}{Remote specimen}{A specimen whose mode is not rewrite or which is not entering; only an entering rewrite focus is excluded.}

\glossaryentry{replacement-carriers}{Replacement carriers}{The disjoint union of host and replacement internal names, and the union of surviving host edges with inserted replacement edges.}

\glossaryentry{retained}{Retained hole label}{A label observed below a surviving host edge, excluding the selected removed edge and all of its descendants.}

\glossaryentry{reverse-traversal}{Reverse operation traversal}{The machine traversing selected operation edges backwards, taking boundary inputs as successful final states.}

\glossaryentry{focus}{Rewrite mode}{The abstract focus mode $\mathsf{rewrite}$, used in the robustness premise for a remote specimen.}

\glossaryentry{robustness}{Robustness}{The abstract condition requiring a rooted, remote rewrite specimen's source step to lead to stuckness or a cost-related quasi-specimen match.}

\glossaryentry{specimen}{Rooted state}{A state satisfying the specimen interface's supplied predicate $\mathsf{Root}$, which its laws require to be preserved by transitions.}

\glossaryentry{focus}{Search mode}{The abstract focus mode $\mathsf{search}$, used with entering specimens in the input-safety premise.}

\glossaryentry{lifting}{Settled specimen}{A specimen in backtracking mode, or in search mode without entering; every settled specimen is remote.}

\glossaryentry{all-edges}{Shallow occurrence}{The occurrence of an edge at the current graph layer, without descending into a box body.}

\glossaryentry{shallow-plugging}{Shallow plugging}{Replacing one outer edge by a graph with its exact signature interface, gluing corresponding boundary ports and retaining the outer interface.}

\glossaryentry{freshness}{Sharp plugging freshness}{The requirement that no retained host hole label occurs in the replacement; labels confined to the removed region may be reused.}

\glossaryentry{indexed-port}{Source port}{A source-list vertex selected by a finite index in the edge label's input signature; the arity equation justifies the lookup.}

\glossaryentry{specimen}{Specimen}{An element of the abstract specimen type, equipped with left and right state endpoints, a mode, and an entering predicate.}

\glossaryentry{specimen}{Specimen system}{The interface packaging specimens, endpoints, modes and predicates, with rootedness preservation, final agreement, and context-traversal laws.}

\glossaryentry{refinement}{State equivalence}{Mutual state refinement for the same cost preorder; it compares successful termination and execution counts in both directions.}

\glossaryentry{refinement}{State refinement}{The requirement that each successful source execution has a successful target execution with counts related by the chosen cost preorder.}

\glossaryentry{determinism}{Stuck state}{A normal state that is not successfully final; it cannot step and does not count as successful termination.}

\glossaryentry{indexed-port}{Target port}{A target-list vertex selected by a finite index in the edge label's output signature.}

\glossaryentry{wire}{Term wire}{The wire type $*$, distinguished from a thunk wire and used for ordinary operation inputs and term arguments.}

\glossaryentry{wire}{Thunk wire}{The wire type $T^n(*)$, carrying the number $n$ of variables bound by a corresponding box.}

\glossaryentry{unary}{Unary operations}{The condition that every selected operation edge has a source list of length one; no one-output condition is imposed.}

\glossaryentry{context}{Unique holes}{The property that any two occurrences observing the same complete typed hole label are equal, at every nesting depth.}

\glossaryentry{cost-orders}{Universal cost order}{The preorder $U$ relating every pair of natural-number counts, so refinement ignores quantitative cost.}

\glossaryentry{up-to}{Up-to closure}{The relation obtained by placing a candidate relation between a left auxiliary state refinement and a right auxiliary state refinement.}

\glossaryentry{vertex-label}{Vertex label}{The internal vertex's assigned label, or the wire label at the indexed position in the relevant boundary list.}

\glossaryentry{wire}{Wire type}{Either the term type $*$ or a thunk type $T^n(*)$ with its binding count.}

\end{document}
